% Acquisition d’une maîtrise des médias et de l’informat — Informatique 3ème — Cours signé OnBuch+
% Programme officiel MINESEC. Compiler avec : tectonic N09-acquisition-maitrise-medias-informat.tex
\newcommand{\DOCMATIERE}{Informatique}
\newcommand{\DOCNIVEAU}{3ème}
\newcommand{\DOCMODULE}{Informatique – classe de 3ème}
\newcommand{\DOCLECON}{N09}
\newcommand{\DOCTITRE}{Acquisition d’une maîtrise des médias et de l’informat}
% OnBuch+ — préambule « cours signés » ENRICHI (compilé avec tectonic / XeLaTeX)
% Système visuel « Pop » d'OnBuch+ : fond crème (jamais blanc), bordures encre
% épaisses, ombres « dures » décalées, titres Archivo Black en capitales.
% Version MATHÉMATIQUES : graphes (pgfplots/tikz), boîtes pédagogiques
% (définition, propriété, méthode, attention, le savais-tu, exercice résolu,
% exercice, TP, bilan), page de garde et signature OnBuch+ ancrée en bas.
%
% Macros attendues dans main.tex AVANT \input{preamble} :
%   \DOCMATIERE  \DOCNIVEAU  \DOCMODULE  \DOCLECON (numéro)  \DOCTITRE
\documentclass[11pt]{article}

\usepackage{fontspec}
\usepackage[french]{babel}
\usepackage[a4paper,margin=2cm,top=3.4cm,bottom=2.8cm,headheight=1.4cm,footskip=1.3cm]{geometry}
\usepackage{amsmath,amssymb,amsfonts}
\usepackage{mathtools} % \xmapsto, \coloneqq... souvent utilisés par les modèles
\usepackage{cancel} % simplifications barrées (bilans de forces)
% \abs{z} (module, valeur absolue) et \norm{u} (norme), utilisés par les modèles
\providecommand{\abs}{}\let\abs\relax\DeclarePairedDelimiter{\abs}{\lvert}{\rvert}
\providecommand{\norm}{}\let\norm\relax\DeclarePairedDelimiter{\norm}{\lVert}{\rVert}
\usepackage[table]{xcolor} % [table] : fournit aussi \rowcolors (alternance de lignes)
\usepackage{pagecolor}
\usepackage{tikz}
\usetikzlibrary{arrows.meta,positioning,calc,shapes.geometric,shapes.misc,shapes.arrows,shapes.symbols,decorations.pathmorphing,decorations.pathreplacing,decorations.markings,patterns,shadows,mindmap,trees,fit,backgrounds,matrix,intersections,angles,quotes}
\usepackage{pgfplots}
\usepackage[version=4]{mhchem} % équations nucléaires \ce{^{235}_{92}U}
\usepackage[european,siunitx]{circuitikz} % schémas électriques
\pgfplotsset{compat=1.18}
\usepgfplotslibrary{fillbetween,groupplots} % aires entre courbes (calcul intégral)
\usepackage{enumitem}
\usepackage{fancyhdr}
% [most] charge toutes les bibliothèques tcolorbox (listings, minted, raster,
% documentation...) et gonfle inutilement la mémoire TeX sur un document long
% (~300 boîtes) : on ne charge que ce qui est réellement utilisé.
\usepackage[skins,breakable]{tcolorbox}
\usepackage{graphicx}
\usepackage{colortbl}
\usepackage{booktabs}
\usepackage{tabularx}
\usepackage{diagbox} % case d'en-tête diagonale des tableaux à double entrée
% Colonnes extensibles centrée / gauche / droite (C, L, R), souvent utilisées par les modèles
\newcolumntype{C}{>{\centering\arraybackslash}X}
\newcolumntype{L}{>{\raggedright\arraybackslash}X}
\newcolumntype{R}{>{\raggedleft\arraybackslash}X}
\usepackage{array}
\usepackage{multirow}
\usepackage{multicol}
\usepackage{siunitx}
% parse-numbers=false : accepte aussi \SI{2,5 \times 10^{-3}}{...}
\sisetup{locale=FR,output-decimal-marker={,},group-minimum-digits=4,parse-numbers=false}
\DeclareSIUnit\ohm{\ensuremath{\symup{\Omega}}} % U+2126 absent de la police
\DeclareSIUnit\division{div} % réglages d'oscilloscope (ms/div, V/div)
\DeclareSIUnit\tr{tr} % tours (vitesse de rotation, ex. \SI{1500}{\tr\per\minute})
\DeclareSIUnit\bit{bit}
\DeclareSIUnit\byte{o} % octet
\DeclareSIUnit\inch{po} % pouce
\DeclareSIUnit\ppp{ppp} % points par pouce (résolution d'image)
\usepackage{qrcode}
% Blocs de code source (PHP, C, HTML, JavaScript, SQL) — spécifique à la
% matière Informatique : listings + habillage aux couleurs OnBuch+ (voir le
% style \lstset ci-dessous, à côté des autres boîtes pédagogiques).
\usepackage{listings}
% \degree : commande fréquemment utilisée par les modèles pour un angle ou une
% température en dehors de \SI{}{} (non fournie par les paquets chargés).
\providecommand{\degree}{^{\circ}}
% \qed (fin de démonstration), utilisé par les modèles sans amsthm
\providecommand{\qed}{\hfill$\square$}
% \barre : double barre des valeurs interdites dans un tableau de variations (tkz-tab)
\providecommand{\barre}{\ensuremath{\|}}
% environnement proof (amsthm non chargé) utilisé par les modèles
\ifdefined\proof\else\newenvironment{proof}[1][Démonstration]{\par\noindent\textit{#1.}\ }{\qed\par}\fi
\usepackage{lastpage}
\usepackage{hyperref}
\hypersetup{hidelinks}

% ============================================================
% Palette « Pop » OnBuch+ — mêmes hex que la classe P (plus_ui.dart)
% ============================================================
\definecolor{poporange}{HTML}{F59321}
\definecolor{popdark}{HTML}{E07A0C}
\definecolor{popcream}{HTML}{FFF6E8}
\definecolor{popink}{HTML}{1C1714}
\definecolor{poppurple}{HTML}{7A5AE0}
\definecolor{popgreen}{HTML}{2E9E6A}
\definecolor{poppink}{HTML}{E0457B}
\definecolor{popblue}{HTML}{2F7DE1}
\definecolor{popgold}{HTML}{C98A12}
\definecolor{popmuted}{HTML}{8A7F76}
\definecolor{popline}{HTML}{EADFD2}
\definecolor{popgray}{HTML}{8A7F76}
\colorlet{popgrey}{popgray}
% Alias tolérants (le modèle nomme parfois les couleurs différemment)
\colorlet{popwhite}{white}
\colorlet{popblack}{popink}
\colorlet{popred}{poppink}
\colorlet{popyellow}{popgold}
% Teintes claires pour fonds de tableaux / zones
\colorlet{popblueL}{popblue!10!white}
\colorlet{poporangeL}{poporange!14!white}
\colorlet{popgreenL}{popgreen!12!white}
\colorlet{poppurpleL}{poppurple!12!white}
\colorlet{poppinkL}{poppink!12!white}
\colorlet{popgoldL}{popgold!14!white}

\pagecolor{popcream}

% ============================================================
% Typographie
% ============================================================
\defaultfontfeatures{Ligatures=TeX}
\setmainfont{PlusJakartaSans-Medium.ttf}[Path=../../../fonts/,BoldFont=PlusJakartaSans-Bold.ttf]
% Maths en sans-serif (Fira Math) avec chiffres et lettres droites de Plus
% Jakarta, pour que les formules se fondent dans le texte.
\usepackage[math-style=ISO,bold-style=ISO]{unicode-math}
\setmathfont{FiraMath-Regular.otf}[Path=../../../fonts/]
\setmathfont{PlusJakartaSans-Medium.ttf}[Path=../../../fonts/,range={up/num,up/{Latin,latin},\mathup}]
\usepackage{icomma} % virgule décimale sans espace en mode math
% Caractères mathématiques Unicode absents des polices de texte (Plus Jakarta,
% Archivo Black) : rendus par leur équivalent mathématique, en texte comme en
% formule — sinon ils s'affichent en carré vide (ex. « Arithmétique dans ℤ »).
\usepackage{newunicodechar}
\newunicodechar{ℝ}{\ensuremath{\mathbb{R}}}
\newunicodechar{ℤ}{\ensuremath{\mathbb{Z}}}
\newunicodechar{ℂ}{\ensuremath{\mathbb{C}}}
\newunicodechar{ℕ}{\ensuremath{\mathbb{N}}}
\newunicodechar{ℚ}{\ensuremath{\mathbb{Q}}}
\newunicodechar{𝔻}{\ensuremath{\mathbb{D}}}
\newunicodechar{α}{\ensuremath{\alpha}}
\newunicodechar{β}{\ensuremath{\beta}}
\newunicodechar{γ}{\ensuremath{\gamma}}
\newunicodechar{δ}{\ensuremath{\delta}}
\newunicodechar{ε}{\ensuremath{\varepsilon}}
\newunicodechar{θ}{\ensuremath{\theta}}
\newunicodechar{λ}{\ensuremath{\lambda}}
\newunicodechar{μ}{\ensuremath{\mu}}
\newunicodechar{σ}{\ensuremath{\sigma}}
\newunicodechar{τ}{\ensuremath{\tau}}
\newunicodechar{φ}{\ensuremath{\varphi}}
\newunicodechar{ψ}{\ensuremath{\psi}}
\newunicodechar{ω}{\ensuremath{\omega}}
\newunicodechar{Σ}{\ensuremath{\Sigma}}
% majuscules produites par \MakeUppercase dans les titres (α-aminés -> Α)
\newunicodechar{Α}{\ensuremath{\alpha}}
\newunicodechar{Β}{\ensuremath{\beta}}
\newunicodechar{Γ}{\ensuremath{\Gamma}}
\newunicodechar{Δ}{\ensuremath{\Delta}}
\newunicodechar{Ε}{\ensuremath{\varepsilon}}
\newunicodechar{Θ}{\ensuremath{\Theta}}
\newunicodechar{Λ}{\ensuremath{\Lambda}}
\newunicodechar{Μ}{\ensuremath{\mu}}
\newunicodechar{Τ}{\ensuremath{\tau}}
\newunicodechar{Φ}{\ensuremath{\Phi}}
\newunicodechar{Ψ}{\ensuremath{\Psi}}
\newunicodechar{Ω}{\ensuremath{\Omega}}
\newunicodechar{↦}{\ensuremath{\mapsto}}
\newunicodechar{∈}{\ensuremath{\in}}
\newunicodechar{∉}{\ensuremath{\notin}}
\newunicodechar{∧}{\ensuremath{\wedge}}
\newunicodechar{⇔}{\ensuremath{\Leftrightarrow}}
\newunicodechar{⟺}{\ensuremath{\Longleftrightarrow}}
\newunicodechar{⇒}{\ensuremath{\Rightarrow}}
\newunicodechar{⟹}{\ensuremath{\Longrightarrow}}
\newunicodechar{∘}{\ensuremath{\circ}}
\newunicodechar{∪}{\ensuremath{\cup}}
\newunicodechar{∩}{\ensuremath{\cap}}
\newunicodechar{≡}{\ensuremath{\equiv}}
\newunicodechar{⊂}{\ensuremath{\subset}}
\newunicodechar{⊥}{\ensuremath{\perp}}
\newunicodechar{∃}{\ensuremath{\exists}}
\newunicodechar{∀}{\ensuremath{\forall}}
\newunicodechar{∛}{\ensuremath{\sqrt[3]{\,}}}
\newunicodechar{∜}{\ensuremath{\sqrt[4]{\,}}}
\newunicodechar{⃗}{\ensuremath{{}^{\to}}}
\newunicodechar{ⁿ}{\ifmmode^{n}\else\textsuperscript{\ensuremath{n}}\fi}
\newunicodechar{ˣ}{\ifmmode^{x}\else\textsuperscript{\ensuremath{x}}\fi}
\newunicodechar{ᵇ}{\ifmmode^{b}\else\textsuperscript{\ensuremath{b}}\fi}
\newunicodechar{ʳ}{\ifmmode^{r}\else\textsuperscript{\ensuremath{r}}\fi}
\newunicodechar{ᵘ}{\ifmmode^{u}\else\textsuperscript{\ensuremath{u}}\fi}
\newunicodechar{ʸ}{\ifmmode^{y}\else\textsuperscript{\ensuremath{y}}\fi}
\newunicodechar{ᵃ}{\ifmmode^{a}\else\textsuperscript{\ensuremath{a}}\fi}
\newunicodechar{ᵉ}{\ifmmode^{e}\else\textsuperscript{\ensuremath{e}}\fi}
\newunicodechar{ᵏ}{\ifmmode^{k}\else\textsuperscript{\ensuremath{k}}\fi}
\newunicodechar{ᵖ}{\ifmmode^{p}\else\textsuperscript{\ensuremath{p}}\fi}
\newunicodechar{ᴺ}{\ifmmode^{N}\else\textsuperscript{\ensuremath{N}}\fi}
\newunicodechar{ᶜ}{\ifmmode^{c}\else\textsuperscript{\ensuremath{c}}\fi}
\newunicodechar{ʰ}{\ifmmode^{h}\else\textsuperscript{\ensuremath{h}}\fi}
\newunicodechar{ᵠ}{\ifmmode^{q}\else\textsuperscript{\ensuremath{q}}\fi}
\newunicodechar{ₙ}{\ifmmode_{n}\else\textsubscript{\ensuremath{n}}\fi}
\newunicodechar{ₐ}{\ifmmode_{a}\else\textsubscript{\ensuremath{a}}\fi}
\newunicodechar{ᵢ}{\ifmmode_{i}\else\textsubscript{\ensuremath{i}}\fi}
\newunicodechar{ᵣ}{\ifmmode_{r}\else\textsubscript{\ensuremath{r}}\fi}
\newunicodechar{ₖ}{\ifmmode_{k}\else\textsubscript{\ensuremath{k}}\fi}
\newunicodechar{ᵦ}{\ifmmode_{\beta}\else\textsubscript{\ensuremath{\beta}}\fi}
\newunicodechar{ₓ}{\ifmmode_{x}\else\textsubscript{\ensuremath{x}}\fi}
\newfontfamily\popdisplay{ArchivoBlack-Regular.ttf}[Path=../../../fonts/]
\newfontfamily\poplogo{SpaceGrotesk-Bold.ttf}[Path=../../../fonts/]
\newfontfamily\popsemi{PlusJakartaSans-SemiBold.ttf}[Path=../../../fonts/]
\newfontfamily\popxbold{PlusJakartaSans-ExtraBold.ttf}[Path=../../../fonts/]
\newfontfamily\popxboldls{PlusJakartaSans-ExtraBold.ttf}[Path=../../../fonts/,LetterSpace=3.0]

\setlist[enumerate]{leftmargin=1.7em,itemsep=5pt,topsep=4pt}
\setlist[itemize]{leftmargin=1.7em,itemsep=4pt,topsep=4pt,label={\color{poporange}\textbullet}}
\setlength{\parindent}{0pt}
\setlength{\parskip}{5pt}
\renewcommand{\arraystretch}{1.25}

% pgfplots : style Pop par défaut
\pgfplotsset{
  every axis/.append style={axis line style={popink,line width=1pt},tick style={popink},
    grid style={popline,line width=0.5pt},label style={font=\small},tick label style={font=\footnotesize}},
  popcurve/.style={poporange,line width=1.8pt,no markers,smooth},
}
% Même style disponible dans un \draw TikZ simple (hors pgfplots)
\tikzset{popcurve/.style={poporange,line width=1.8pt}}
\tikzset{popgrid/.style={popline,very thin}}
\colorlet{popcurve}{poporange} % le nom du style est parfois utilisé comme couleur

% ============================================================
% Pill / squiggle
% ============================================================
\newcommand{\poppill}[3]{%
  \tikz[baseline=-0.55ex]\node[draw=popink,line width=1.2pt,rounded corners=2.6pt,%
    fill=#2,inner xsep=6.5pt,inner ysep=3pt]{%
    \popxboldls\fontsize{7}{7}\selectfont\color{#3}\MakeUppercase{#1}};%
}
\newcommand{\popsquiggle}[1]{%
  \tikz[baseline=-0.4ex]{
    \draw[#1,line width=2.6pt,line cap=round]
      plot [smooth] coordinates {(0,0.09) (0.34,0.01) (0.68,0.21) (1.02,0.01) (1.36,0.10)};
  }%
}
% Mot-clé surligné (marqueur orange sous le texte)
\newcommand{\cle}[1]{{\popxbold\color{popdark}#1}}

% ============================================================
% En-tête / pied
% ============================================================
\pagestyle{fancy}
\fancyhf{}
\renewcommand{\headrulewidth}{0pt}
\renewcommand{\footrulewidth}{0pt}
\lhead{\raisebox{-0.15ex}{\poplogo\fontsize{13}{13}\selectfont\color{popink}OnBuch\color{poporange}+}}
\rhead{\poppill{\DOCMATIERE\ \textbullet\ \DOCNIVEAU\ \textbullet\ Leçon \DOCLECON}{white}{popink}}
\lfoot{\popxbold\fontsize{7.6}{7.6}\selectfont\color{popmuted}\MakeUppercase{Cours signé OnBuch+}\ \textbullet\ {\popsemi\DOCTITRE}}
\rfoot{\tikz[baseline=-0.6ex]\node[rounded corners=8.5pt,fill=popink,minimum height=17pt,minimum width=17pt,inner xsep=5pt,inner sep=0]{\popxbold\fontsize{8}{8}\selectfont\color{popcream}\thepage\,/\,\pageref*{LastPage}};}

% ============================================================
% Titres
% ============================================================
\newcommand{\chaptitre}[2]{%
  \begin{tcolorbox}[enhanced,colback=poporange,colframe=popink,boxrule=2.6pt,arc=11pt,
      drop shadow=popink,left=15pt,right=15pt,top=12pt,bottom=11pt,before skip=3pt,after skip=18pt]
    \poppill{#2}{popink}{popcream}\par\vspace{9pt}
    {\raggedright\hyphenpenalty=10000\exhyphenpenalty=10000\popdisplay\fontsize{22}{24}\selectfont\color{popcream}\MakeUppercase{#1}\par}
  \end{tcolorbox}
}
% Section numérotée : pastille orange avec le numéro + titre Archivo
\newcounter{popsec}
\newcommand{\coursec}[1]{%
  \stepcounter{popsec}\setcounter{popsub}{0}%
  \par\vspace{16pt}\noindent
  \begin{minipage}{\linewidth}
  \tikz[baseline=-0.7ex]\node[draw=popink,line width=1.6pt,fill=poporange,rounded corners=4pt,minimum size=22pt,inner sep=0]{\popdisplay\fontsize{12}{12}\selectfont\color{popcream}\thepopsec};%
  \hspace{8pt}{\popdisplay\fontsize{15}{17}\selectfont\color{popink}\MakeUppercase{#1}}\par
  \vspace{4pt}\noindent{\color{poporange}\rule{36pt}{3.6pt}}
  \end{minipage}\par\nopagebreak\vspace{8pt}
}
\newcounter{popsub}
\newcommand{\courssub}[1]{%
  \stepcounter{popsub}%
  \par\vspace{9pt}\noindent{\popxbold\fontsize{11.5}{13}\selectfont\color{popink}{\color{poporange}\thepopsec.\thepopsub}\hspace{6pt}#1}\par\nopagebreak\vspace{4pt}
}

% ============================================================
% Boîtes pédagogiques — même recette : fond blanc, bordure épaisse colorée,
% ombre dure encre, badge plein. Titre optionnel [#1].
% ============================================================
\tcbset{popbase/.style={breakable,enhanced,colback=white,boxrule=2.3pt,arc=9pt,
  shadow={3pt}{-3pt}{0pt}{popink},left=14pt,right=14pt,top=11pt,bottom=11pt,before skip=11pt,after skip=15pt,
  fontupper=\popsemi\fontsize{10.6}{14.5}\selectfont\color{popink},
  fontlower=\popsemi\fontsize{10.6}{14.5}\selectfont\color{popink}}}
% Coche dessinée (le glyphe ✓ n'existe pas dans les polices du document)
\newcommand{\popcheck}[1]{\tikz[baseline=-0.5ex]\draw[#1,line width=1.6pt,line cap=round,line join=round](-0.13,0.0)--(-0.04,-0.09)--(0.13,0.1);}
\providecommand{\coche}{\popcheck{popgreen}}
\providecommand{\resultat}[1]{\par\smallskip\noindent\fcolorbox{popgreen}{popgreenL}{\parbox{\dimexpr\linewidth-2\fboxsep-2\fboxrule\relax}{\textbf{Résultat.} #1}}\par\smallskip}
\newcommand{\popboxhead}[3]{\poppill{#1}{#2}{white}\ifx\relax#3\relax\else\hspace{6pt}{\popxbold\fontsize{10.6}{12}\selectfont\color{popink}#3}\fi\par\vspace{7pt}}

\newtcolorbox{aretenir}[1][]{popbase,colframe=popgreen,before upper={\popboxhead{À retenir}{popgreen}{#1}}}
\newtcolorbox{exemplebox}[1][]{popbase,colframe=poporange,before upper={\popboxhead{Exemple}{poporange}{#1}}}
\newtcolorbox{definition}[1][]{popbase,colframe=popblue,before upper={\popboxhead{Définition}{popblue}{#1}}}
\newtcolorbox{propriete}[1][]{popbase,colframe=poppurple,before upper={\popboxhead{Propriété}{poppurple}{#1}}}
\newtcolorbox{methode}[1][]{popbase,colframe=popgold,before upper={\popboxhead{Méthode}{popgold}{#1}}}
\newtcolorbox{attention}[1][]{popbase,colframe=poppink,before upper={\popboxhead{Attention}{poppink}{#1}}}
\newtcolorbox{experience}[1][]{popbase,colframe=popblue,colback=popblueL,before upper={\popboxhead{Expérience}{popblue}{#1}}}
% « Le savais-tu ? » : carte violette pleine (contexte, culture, Cameroun)
\newtcolorbox{savaistu}[1][]{popbase,colback=poppurple,colframe=popink,
  fontupper=\popsemi\fontsize{10.4}{14.2}\selectfont\color{white},
  before upper={\poppill{Le savais-tu\,?}{white}{poppurple}\ifx\relax#1\relax\else\hspace{6pt}{\popxbold\color{white}#1}\fi\par\vspace{7pt}}}
% Exercice résolu : énoncé en haut, \tcblower puis la solution sur fond crème
\newcounter{popexo}\newcounter{popexr}
\newtcolorbox{exoresolu}[1][]{popbase,colframe=poporange,colbacklower=popcream,
  segmentation style={popink,line width=1.2pt,dashed},
  before upper={\stepcounter{popexr}\popboxhead{Exercice résolu \thepopexr}{poporange}{#1}},
  before lower={\poppill{Solution}{popink}{popcream}\par\vspace{6pt}}}
% Exercice (non résolu en place) — la correction est dans la partie Corrigés
\newtcolorbox{exercice}[1][]{popbase,colframe=popink,
  before upper={\stepcounter{popexo}\popboxhead{Exercice \thepopexo}{popink}{#1}}}
% Corrigé
\newtcolorbox{corrige}[1][]{popbase,colframe=popgreen,colback=popgreenL,
  before upper={\popboxhead{Corrigé}{popgreen}{#1}}}
% Objectifs / prérequis / bilan
\newtcolorbox{objectifs}{popbase,colframe=popink,colback=poporangeL,
  before upper={\popboxhead{Objectifs de la leçon}{popink}{}}}
\newtcolorbox{prerequis}{popbase,colframe=popink,colback=popblueL,
  before upper={\popboxhead{Prérequis}{popblue}{}}}
\newtcolorbox{bilan}{popbase,colframe=popink,colback=white,boxrule=2.8pt,
  before upper={\popboxhead{Fiche bilan}{poporange}{L'essentiel à connaître pour l'examen}}}
% Figure encadrée avec légende
\newtcolorbox{popfigure}[1][]{enhanced,breakable=false,colback=white,colframe=popline,boxrule=1.4pt,arc=8pt,
  left=10pt,right=10pt,top=10pt,bottom=8pt,before skip=10pt,after skip=12pt,halign=center,
  lower separated=false,halign lower=center,
  fontlower=\popsemi\fontsize{9}{11.5}\selectfont\color{popmuted},
  #1}
\newcommand{\legende}[1]{\par\vspace{6pt}{\popsemi\fontsize{9}{11.5}\selectfont\color{popmuted}#1}}

% ============================================================
% Bloc de code source — \begin{lstlisting}[language=Php]...\end{lstlisting}
% (ou language=C, HTML, JavaScript, SQL ; option omise = pseudo-code brut).
% Même recette visuelle que les figures (\popfigure) : cadre encre épais,
% fond clair (popcream), légende possible juste après via \legende{...}
% comme pour une figure (listings ne sait pas arrondir ses coins : on garde
% un cadre rectangulaire, cohérent avec les autres tableaux du document).
% CONTRAT : les modèles ne doivent utiliser QUE \begin{lstlisting}[...] tel
% quel (pas de \begin{tcblisting}, pas de \verb pour un bloc multi-lignes).
% ============================================================
\lstdefinestyle{poplst}{
  backgroundcolor=\color{popcream},
  basicstyle=\popsemi\fontsize{8.6}{11}\selectfont\color{popink},
  keywordstyle=\bfseries\color{popblue},
  commentstyle=\itshape\color{popmuted},
  stringstyle=\color{popgreen},
  numberstyle=\tiny\color{popmuted},
  identifierstyle=\color{popink},
  frame=l,
  rulecolor=\color{popink},
  framerule=1.6pt,
  framesep=7pt,
  framexleftmargin=7pt,
  framexrightmargin=7pt,
  xleftmargin=2pt,
  xrightmargin=2pt,
  breaklines=true,
  breakatwhitespace=false,
  showstringspaces=false,
  showspaces=false,
  showtabs=false,
  tabsize=2,
  columns=flexible,
  keepspaces=true,
  numbers=none,
  aboveskip=10pt,
  belowskip=10pt,
  captionpos=b,
}
\lstset{style=poplst, language=PHP}
% La distribution TeX embarquée ne fournit pas le fichier de langue
% JavaScript de listings (lstlang*.dat) : on la redéfinit nous-mêmes
% pour que \begin{lstlisting}[language=JavaScript] fonctionne.
\lstdefinelanguage{JavaScript}{
  keywords={break,case,catch,class,const,continue,debugger,default,delete,do,
    else,export,extends,finally,for,function,if,import,in,instanceof,let,new,
    return,static,super,switch,this,throw,try,typeof,var,void,while,with,yield,
    async,await,of,get,set},
  keywordstyle=\bfseries\color{popblue},
  ndkeywords={true,false,null,undefined,NaN,Infinity},
  ndkeywordstyle=\bfseries\color{poporange},
  identifierstyle=\color{popink},
  sensitive=true,
  comment=[l]{//},
  morecomment=[s]{/*}{*/},
  string=[b]",
  morestring=[b]',
  morestring=[b]`,
}
% Même limitation pour CSS et les fichiers de configuration (apacheconf,
% nginx, ini...) : ils ne sont pas fournis. CSS et un style générique de
% type "config" (commentaires # ou ;, pas de mots-clés) couvrent les besoins.
\lstdefinelanguage{CSS}{
  keywords={color,background,background-color,margin,padding,border,
    border-radius,font,font-size,font-weight,font-family,width,height,
    display,position,top,left,right,bottom,float,clear,text-align,
    line-height,list-style,cursor,overflow,box-shadow,transition,
    flex,grid,align-items,justify-content,z-index},
  keywordstyle=\bfseries\color{popblue},
  identifierstyle=\color{popink},
  sensitive=false,
  comment=[s]{/*}{*/},
  string=[b]",
  morestring=[b]',
}
\lstdefinelanguage{apacheconf}{
  sensitive=false,
  morecomment=[l]{\#},
}
\lstdefinelanguage{Apache}{
  sensitive=false,
  morecomment=[l]{\#},
}
\lstdefinelanguage{text}{sensitive=false}
\lstdefinelanguage{powershell}{
  keywords={New-Object,New-SmbShare,Get-Acl,Set-Acl,Get-ChildItem,Set-Content,
    Get-Content,Write-Host,ForEach-Object,Where-Object,Select-Object,
    Test-Path,Remove-Item,Copy-Item,Move-Item,Start-Process,Stop-Process,
    If,Else,ElseIf,ForEach,While,Function,Return,Try,Catch},
  keywordstyle=\bfseries\color{popblue},
  identifierstyle=\color{popink},
  sensitive=false,
  comment=[l]{\#},
  string=[b]",
  morestring=[b]',
}
\lstdefinelanguage{ini}{
  sensitive=false,
  morecomment=[l]{;},
  morecomment=[l]{\#},
}

% ============================================================
% Signature OnBuch+
% ============================================================
\newcommand{\poplogoinline}[1]{{\poplogo\fontsize{#1}{#1}\selectfont\color{popink}OnBuch\color{poporange}+}}
\newcommand{\signatureonbuch}{%
  \begin{tcolorbox}[enhanced,colback=popink,colframe=popink,boxrule=2.2pt,arc=10pt,
      drop shadow=poporange,left=14pt,right=14pt,top=10pt,bottom=10pt,before skip=0pt,after skip=0pt]
    \tikz[baseline=-0.7ex]\node[circle,fill=popgreen,draw=popcream,line width=1.3pt,minimum size=22pt,inner sep=0]{\popcheck{white}};%
    \hspace{9pt}\begin{minipage}[c]{0.62\linewidth}
      {\popxbold\fontsize{12}{13}\selectfont\color{popcream}Signé\ \textbullet\ {\poplogo OnBuch\color{poporange}+}}\par\vspace{2pt}
      {\raggedright\popsemi\fontsize{8}{10.5}\selectfont\color{popline}Ludovic A.\\\MakeUppercase{Conforme au programme officiel MINESEC}\par}
    \end{minipage}\hfill
    {\poplogo\fontsize{22}{22}\selectfont\color{popcream}OnBuch\color{poporange}+}
  \end{tcolorbox}%
}

% ============================================================
% Page de garde
% #1 titre · #2 matière · #3 niveau · #4 module · #5 n° leçon · #6 durée ·
% #7 description / objectifs (texte ou itemize)
% ============================================================
\newcommand{\pagegarde}[7]{%
  \thispagestyle{empty}
  \newgeometry{a4paper,margin=1.7cm,top=1.6cm,bottom=1.4cm}%
  \begin{tikzpicture}[remember picture,overlay]
    \fill[poporange!18] ([xshift=-3.2cm,yshift=-3.6cm]current page.north east) circle (4.6cm);
    \fill[poppurple!14] ([xshift=2.4cm,yshift=7.6cm]current page.south west) circle (3.8cm);
  \end{tikzpicture}%
  {\poplogo\fontsize{30}{30}\selectfont\color{popink}OnBuch\color{poporange}+}\hfill
  \raisebox{0.6ex}{\poppill{Cours signé \textbullet\ Premium}{popink}{popcream}}\par
  \vspace{1pt}\popsquiggle{poppurple}\par
  \vspace{8pt}
  \begin{tcolorbox}[enhanced,colback=poporange,colframe=popink,boxrule=2.8pt,arc=13pt,
      drop shadow=popink,left=18pt,right=18pt,top=12pt,bottom=13pt,after skip=10pt]
    \poppill{#2 \textbullet\ #3}{popink}{popcream}\hspace{5pt}\poppill{#4}{white}{popink}\par\vspace{8pt}
    {\popdisplay\fontsize{14}{14}\selectfont\color{popink}LEÇON #5}\par\vspace{4pt}
    {\raggedright\hyphenpenalty=10000\exhyphenpenalty=10000\popdisplay\fontsize{25}{27}\selectfont\color{popcream}\MakeUppercase{#1}\par}
    \vspace{8pt}
    {\popxbold\fontsize{9.5}{9.5}\selectfont\color{popink}\MakeUppercase{Durée conseillée : #6}}
  \end{tcolorbox}
  \begin{tcolorbox}[enhanced,colback=white,colframe=popink,boxrule=2.2pt,arc=10pt,
      drop shadow=popink,left=16pt,right=16pt,top=10pt,bottom=10pt,after skip=8pt]
    \poppill{Dans cette leçon}{poppurple}{white}\par\vspace{5pt}
    \popsemi\fontsize{9.3}{12.6}\selectfont\color{popink}#7
  \end{tcolorbox}
  \noindent\begin{minipage}[t]{0.487\linewidth}
    \begin{tcolorbox}[enhanced,colback=popgreen,colframe=popink,boxrule=2.2pt,arc=10pt,
        drop shadow=popink,left=11pt,right=11pt,top=10pt,bottom=10pt,height=3.1cm,valign=center]
      \tcbox[colback=white,colframe=popink,boxrule=1.3pt,arc=5pt,boxsep=0pt,nobeforeafter,
        left=3pt,right=3pt,top=3pt,bottom=3pt,valign=center]{\qrcode[height=1.9cm]{https://play.google.com/store/apps/details?id=com.luvvix.onbuch&hl=fr}}%
      \hspace{8pt}\begin{minipage}[c]{0.5\linewidth}
        {\popdisplay\fontsize{11}{12.5}\selectfont\color{white}\MakeUppercase{Télécharge OnBuch}}\par\vspace{4pt}
        {\raggedright\popsemi\fontsize{8.4}{11}\selectfont\color{white}Gratuit \textbullet\ Google Play\\Tuteur IA, annales, concours, résultats\par}
      \end{minipage}
    \end{tcolorbox}
  \end{minipage}\hfill
  \begin{minipage}[t]{0.487\linewidth}
    \begin{tcolorbox}[enhanced,colback=poppurple,colframe=popink,boxrule=2.2pt,arc=10pt,
        drop shadow=popink,left=11pt,right=11pt,top=10pt,bottom=10pt,height=3.1cm,valign=center]
      \tikz[baseline=-0.7ex]\node[circle,fill=white,draw=popink,line width=1.3pt,minimum size=1.6cm,inner sep=0]{\popdisplay\fontsize{24}{24}\selectfont\color{poppurple}+};%
      \hspace{8pt}\begin{minipage}[c]{0.6\linewidth}
        {\popdisplay\fontsize{11}{12.5}\selectfont\color{white}\MakeUppercase{Passe à OnBuch+}}\par\vspace{4pt}
        {\raggedright\popsemi\fontsize{8.4}{11}\selectfont\color{white}Tous les cours signés, Tuteur IA illimité, fiches PDF — onglet OnBuch+ de l'app\par}
      \end{minipage}
    \end{tcolorbox}
  \end{minipage}
  \vspace{8pt}
  \popcontenu
  \vfill
  \signatureonbuch
  \restoregeometry
  \newpage
}

% Rangée « Ce cours contient » (page de garde)
\newcommand{\poptile}[3]{% #1 couleur #2 gros texte #3 légende
  \begin{tcolorbox}[enhanced,colback=#1,colframe=popink,boxrule=2pt,arc=9pt,shadow={2.5pt}{-2.5pt}{0pt}{popink},
      left=6pt,right=6pt,top=9pt,bottom=9pt,halign=center,valign=center,height=1.75cm,width=\linewidth,nobeforeafter]
    {\popdisplay\fontsize{12.5}{13}\selectfont\color{white}\MakeUppercase{#2}}\par\vspace{3pt}
    {\centering\popsemi\fontsize{7.8}{9.5}\selectfont\color{white}#3\par}
  \end{tcolorbox}}
\newcommand{\popcontenu}{%
  \par\noindent{\popxboldls\fontsize{7.5}{7.5}\selectfont\color{popmuted}CE COURS SIGNÉ CONTIENT}\par\vspace{7pt}
  \noindent\begin{minipage}[t]{0.235\linewidth}\poptile{popblue}{Cours}{complet, illustré, pas à pas}\end{minipage}\hfill
  \begin{minipage}[t]{0.235\linewidth}\poptile{poporange}{Méthodes}{et exercices résolus}\end{minipage}\hfill
  \begin{minipage}[t]{0.235\linewidth}\poptile{poppink}{Exercices}{type examen + corrigés}\end{minipage}\hfill
  \begin{minipage}[t]{0.235\linewidth}\poptile{popgreen}{Fiche bilan}{l'essentiel à retenir}\end{minipage}\par}

% Sommaire
\newtcolorbox{sommaire}{enhanced,colback=white,colframe=popink,boxrule=2.2pt,arc=10pt,
  drop shadow=popink,left=18pt,right=18pt,top=13pt,bottom=13pt,before skip=6pt,after skip=20pt,
  before upper={\poppill{Sommaire}{poppurple}{white}\par\vspace{9pt}\popsemi\fontsize{10.6}{14.5}\selectfont\color{popink}}}

% Signature de fin de cours — poussée tout en bas de la dernière page
\newcommand{\findecours}{%
  \par\vfill\nopagebreak
  \begin{center}\popsquiggle{poporange}\end{center}\vspace{4pt}
  \signatureonbuch
}
\AtBeginDocument{\def\llbracket{\mathopen{[\mkern-3.5mu[}}\def\rrbracket{\mathclose{]\mkern-3.5mu]}}} % intervalles d'entiers (glyphe ⟦ absent de la police)
% Glyphes absents de Fira Math : redéfinis à partir de glyphes disponibles
\AtBeginDocument{%
  \def\setminus{\mathbin{\backslash}}\let\smallsetminus\setminus
  \def\vdots{\mathord{\vbox{\baselineskip3.2pt\lineskiplimit0pt\kern4pt\hbox{.}\hbox{.}\hbox{.}}}}%
  \def\ddots{\mathinner{\mkern1mu\raise6.5pt\hbox{.}\mkern2mu\raise3.6pt\hbox{.}\mkern2mu\raise0.7pt\hbox{.}\mkern1mu}}%
  \def\triangle{\mathord{\scalebox{0.9}{$\symup{\Delta}$}}}%
  \def\nabla{\mathord{\raisebox{\depth}{\scalebox{1}[-1]{$\symup{\Delta}$}}}}%
  \def\checkmark{\popcheck{popdark}}%
  \def\star{\mathbin{\ast}}\def\bigstar{\mathord{\ast}}%
  \def\diamond{\mathbin{\scalebox{0.6}{$\Diamond$}}}%
  \def\bigcup{\mathop{\mathchoice{\raisebox{-0.3em}{\scalebox{1.7}{$\cup$}}}{\scalebox{1.25}{$\cup$}}{\cup}{\cup}}\displaylimits}%
  \def\bigcap{\mathop{\mathchoice{\raisebox{-0.3em}{\scalebox{1.7}{$\cap$}}}{\scalebox{1.25}{$\cap$}}{\cap}{\cap}}\displaylimits}%
  \def\lesseqqgtr{\mathrel{\lessgtr}}\def\bigoplus{\mathop{\oplus}\displaylimits}%
}
\providecommand{\ssi}{si et seulement si} % abréviation fréquente
\AtBeginDocument{\providecommand{\wideparen}{\overparen}} % arc de cercle

%% FIGFIX-BEGIN v5 — correctifs globaux des figures (docs/onbuch-generation/tools/figfix)
\usepackage{adjustbox}\usepackage{etoolbox}\tcbuselibrary{raster}
% toute figure TikZ/pgfplots plus large que la ligne est réduite à la largeur disponible
\BeforeBeginEnvironment{tikzpicture}{\begin{adjustbox}{max width=\linewidth}}
\AfterEndEnvironment{tikzpicture}{\end{adjustbox}}
% légendes pgfplots lisibles par-dessus les courbes
\pgfplotsset{every axis/.append style={legend style={fill=white,fill opacity=0.92,text opacity=1,draw=black!15,inner sep=3pt}}}
% étiquettes d'arêtes (syntaxe edge["..."]) avec fond blanc
\tikzset{every edge quotes/.append style={fill=white,inner sep=1.2pt}}
%% FIGFIX-END
\begin{document}

\pagegarde{Acquisition d’une maîtrise des médias et de l’informat}{Informatique}{3ème}{Informatique – classe de 3ème}{N09}{1 séance (fiche de progression harmonisée)}{Cette leçon permet aux élèves d'identifier les critères d'une bonne information et de les appliquer pour évaluer, rédiger et justifier une analyse critique de contenus médiatiques quotidiens.
\vspace{4pt}
\begin{multicols}{2}\small
\begin{itemize}
\item À la fin de cette leçon, je sais définir l'information et en distinguer les principales caractéristiques.
\item À la fin de cette leçon, je sais énumérer et expliquer les six critères de qualité d'une information (exactitude, pertinence, actualité, fiabilité, objectivité, accessibilité).
\item À la fin de cette leçon, je sais appliquer une grille d'évaluation à un message reçu sur les réseaux sociaux.
\item À la fin de cette leçon, je sais repérer les indices de biais, de manipulation ou de désinformation dans un contenu médiatique.
\item À la fin de cette leçon, je sais structurer une justification écrite de mon évaluation en citant les critères utilisés.
\item À la fin de cette leçon, je sais argumenter oralement mon jugement sur la crédibilité d'une information face à mes pairs.
\end{itemize}\end{multicols}}

\begin{sommaire}
\begin{multicols}{2}
\begin{enumerate}
\item 1. Qu'est-ce qu'une information ?
\item 2. Les six critères d'une bonne information
\item 3. Démarche pratique d'évaluation d'une information
\item 4. Rédiger et justifier une évaluation critique
\item 5. Application à des situations camerounaises concrètes
\item Activité d'intégration
\item Méthodes et exercices résolus
\item Exercices
\item Corrigés des exercices
\item Fiche bilan
\end{enumerate}
\end{multicols}
\end{sommaire}

\begin{objectifs}
\begin{itemize}
\item À la fin de cette leçon, je sais définir l'information et en distinguer les principales caractéristiques.
\item À la fin de cette leçon, je sais énumérer et expliquer les six critères de qualité d'une information (exactitude, pertinence, actualité, fiabilité, objectivité, accessibilité).
\item À la fin de cette leçon, je sais appliquer une grille d'évaluation à un message reçu sur les réseaux sociaux.
\item À la fin de cette leçon, je sais repérer les indices de biais, de manipulation ou de désinformation dans un contenu médiatique.
\item À la fin de cette leçon, je sais structurer une justification écrite de mon évaluation en citant les critères utilisés.
\item À la fin de cette leçon, je sais argumenter oralement mon jugement sur la crédibilité d'une information face à mes pairs.
\item À la fin de cette leçon, je sais transférer ces compétences à des situations réelles (actualité locale, publicités, annonces officielles).
\end{itemize}
\end{objectifs}

\begin{prerequis}
\begin{itemize}
\item Notion de source d'information (enseignée en 4ème).
\item Différence entre fait et opinion (début d'année 3ème).
\item Utilisation basique d'un moteur de recherche pour vérifier une information.
\item Structure d'un paragraphe argumentatif (français, 4ème).
\end{itemize}
\end{prerequis}

\newpage

\coursec{Qu'est-ce qu'une information ?}

\textbf{Situation d'accroche.} Il est 17 h 30 dans le quartier Melen à Yaoundé. Sur le groupe WhatsApp du quartier, un message circule à toute vitesse : \emph{« Alerte ! Le prix du carburant va doubler demain matin. Partagez massivement ! »}. En quelques minutes, le message est transféré des dizaines de fois. À 18 h, une longue file de motos-taxis et de voitures se forme devant la station-service du carrefour. Certains habitants ont quitté leur travail plus tôt, d'autres ont emprunté de l'argent pour faire le plein. Le lendemain matin... rien. Le prix n'a pas bougé. Le message venait d'un simple particulier qui avait mal compris une discussion à la radio. Temps perdu, argent dépensé, panique inutile.

Cette scène, tu l'as peut-être déjà vécue ou observée. Elle pose une question fondamentale : \textbf{comment distinguer une information fiable d'une fausse nouvelle avant d'agir ?} Pour répondre à cette question, il faut d'abord comprendre ce qu'est réellement une information, d'où elle vient et comment elle circule. C'est l'objectif de cette première partie de la leçon.

\courssub{Définition de l'information}

Dans la vie courante, nous utilisons le mot « information » à tout moment : information à la radio, information sur Internet, information de dernière minute... Mais en informatique et dans les médias, ce mot a un sens précis qu'il faut bien maîtriser.

\begin{definition}[Information]
Une \cle{information} est une \textbf{donnée traitée, contextualisée et porteuse de sens} pour un destinataire. Autrement dit :
\[
\text{Information} = \text{Donnée} + \text{Contexte} + \text{Sens}
\]
\begin{itemize}
\item la \cle{donnée} est l'élément brut (un chiffre, un mot, une image) ;
\item le \cle{contexte} précise où, quand, qui et dans quelles circonstances ;
\item le \cle{sens} est ce que le destinataire peut comprendre et utiliser pour décider ou agir.
\end{itemize}
\end{definition}

Prenons un exemple simple. Le nombre \texttt{850} écrit seul sur un papier est une simple donnée : il ne veut rien dire. Si l'on ajoute le contexte — « prix du litre de super à la station Total de Mokolo, le 12 octobre 2025 » — alors ce nombre devient une \textbf{information} : « le litre de super coûte 850 FCFA à Mokolo ». Le chauffeur de taxi qui lit cette information peut décider de faire le plein ou d'attendre : l'information a un sens pour lui et déclenche une action.

\begin{aretenir}[L'essentiel]
Une donnée brute ne devient une information que si elle est \textbf{placée dans un contexte} et si elle \textbf{a un sens} pour celui qui la reçoit. Sans contexte ni sens, il n'y a pas d'information, seulement du bruit.
\end{aretenir}

\courssub{Donnée brute, information et connaissance : trois niveaux à ne pas confondre}

Pour bien comprendre, il faut distinguer trois niveaux qui s'enchaînent comme les étages d'une maison.

\begin{definition}[Les trois niveaux]
\begin{itemize}
\item La \cle{donnée brute} : un fait élémentaire, non interprété (un chiffre, un nom, une date). Exemple : \texttt{37,8}.
\item L'\cle{information} : la donnée mise en contexte, qui devient compréhensible et utile. Exemple : « la température de l'élève Amina, mesurée ce matin à l'infirmerie du collège, est de 37,8 °C ».
\item La \cle{connaissance} : ce que l'on construit en reliant plusieurs informations à son expérience et à ses savoirs, et qui permet de juger et d'agir. Exemple : « une température supérieure à 38 °C indique une fièvre ; Amina est légèrement fiévreuse, il faut surveiller son état ».
\end{itemize}
\end{definition}

\begin{exemplebox}[Du mot à la décision]
Imaginons le mot \texttt{NGONG} affiché sur un écran.
\begin{enumerate}
\item \textbf{Donnée brute} : la suite de lettres N-G-O-N-G.
\item \textbf{Information} : « Ngong est un quartier de la ville de Ngaoundéré où se trouve un marché très fréquenté le samedi ».
\item \textbf{Connaissance} : un commerçant de Yaoundé qui connaît ce marché sait qu'il peut y acheter des oignons moins chers le samedi matin et les revendre à Yaoundé avec un bénéfice. Il transforme l'information en décision commerciale.
\end{enumerate}
\end{exemplebox}

La frontière est donc claire : la donnée est la matière première, l'information est le produit fabriqué, et la connaissance est l'usage intelligent que l'on en fait.

\courssub{Les sources de l'information : primaires et secondaires}

D'où vient une information ? C'est une question capitale, car la fiabilité d'une information dépend beaucoup de sa source. On distingue deux grandes familles de sources.

\begin{definition}[Source primaire et source secondaire]
\begin{itemize}
\item Une \cle{source primaire} est une source \textbf{directe et de première main} : un témoin direct d'un événement, un document officiel (communiqué, décret, arrêté), un procès-verbal, une photographie prise sur place, une donnée mesurée soi-même.
\item Une \cle{source secondaire} est une source qui \textbf{rapporte, commente ou transforme} une information venue d'ailleurs : un article de journal, un reportage radio, un billet de blog, une publication sur les réseaux sociaux, un message WhatsApp transféré.
\end{itemize}
\end{definition}

\begin{center}
\begin{tabularx}{\linewidth}{|l|X|X|}
\hline
\rowcolor{popblueL}
\textbf{Critère} & \textbf{Source primaire} & \textbf{Source secondaire} \\
\hline
Lien avec l'événement & Direct, de première main & Indirect, rapporté par un intermédiaire \\
\hline
Exemples & Communiqué officiel du MINCOMMERCE, témoin oculaire, acte de naissance, photo prise sur place & Article de journal en ligne, post Facebook d'un particulier, message WhatsApp transféré, billet de blog \\
\hline
Fiabilité générale & Généralement plus élevée (mais à vérifier : authenticité du document) & Variable : dépend du sérieux de l'intermédiaire \\
\hline
Risque principal & Document falsifié ou témoin de mauvaise foi & Déformation, raccourci, erreur de recopie, rumeur \\
\hline
\end{tabularx}
\end{center}

\begin{exemplebox}[Trois informations, trois sources]
Reprenons la rumeur de notre situation d'accroche et comparons trois messages :
\begin{enumerate}
\item \textbf{Communiqué du MINCOMMERCE} publié sur le site officiel du ministère : « Aucune hausse des prix des produits pétroliers n'est prévue. » $\rightarrow$ \textbf{source primaire} (document officiel, émanant de l'autorité compétente).
\item \textbf{Article d'un journal en ligne} qui cite ce communiqué et ajoute l'analyse d'un journaliste $\rightarrow$ \textbf{source secondaire} sérieuse, car elle vérifie et cite sa source.
\item \textbf{Post Facebook d'un particulier} : « On m'a dit que l'essence va doubler demain ! » $\rightarrow$ \textbf{source secondaire} très fragile : anonyme, sans preuve, sans vérification.
\end{enumerate}
Face à ces trois messages, lequel doit guider ta décision ? Le premier, évidemment.
\end{exemplebox}

\begin{attention}[Popularité n'est pas fiabilité]
Une erreur très fréquente consiste à \textbf{confondre la popularité d'un message et sa fiabilité}. Qu'un message soit partagé 10 000 fois sur WhatsApp ou Facebook ne prouve absolument pas qu'il est vrai : cela prouve seulement qu'il a été beaucoup transféré. Une rumeur spectaculaire se propage souvent plus vite qu'une information vérifiée, car elle fait peur ou fait rêver. La bonne question n'est pas « combien de personnes l'ont partagé ? » mais « \textbf{qui l'a dit, et sur quelle preuve ?} ».
\end{attention}

\courssub{Le cycle de vie de l'information}

Une information n'apparaît pas par magie sur ton téléphone : elle suit un parcours, que l'on appelle le \cle{cycle de vie de l'information}. Le comprendre permet de savoir à quelle étape on se situe quand on reçoit un message, et donc quelles vérifications faire.

\begin{definition}[Cycle de vie de l'information]
Le \cle{cycle de vie de l'information} comporte quatre grandes étapes :
\begin{enumerate}
\item la \cle{création} : l'information est produite (rédaction d'un communiqué, prise d'une photo, enregistrement d'un témoignage) ;
\item la \cle{diffusion} : elle est transmise au public (site web, radio, réseaux sociaux, bouche-à-oreille) ;
\item la \cle{réception} : elle est reçue, lue et interprétée par les destinataires, qui peuvent la relayer (c'est l'étape où toi, tu interviens) ;
\item l'\cle{archivage} : elle est conservée pour une consultation future (archives d'un journal, base de données d'une administration, historique d'un site web).
\end{enumerate}
\end{definition}

\begin{popfigure}
\begin{center}
\begin{tikzpicture}[
    etape/.style={rectangle, rounded corners=4pt, draw=popblue, fill=popblueL, thick,
                  minimum width=3.4cm, minimum height=1.1cm, align=center, font=\small\bfseries},
    descr/.style={font=\scriptsize, text=popdark, align=center},
    fleche/.style={-{Stealth[length=3mm]}, very thick, poporange}
]
\node[etape] (creation) at (0,3.2) {1. Création};
\node[descr, below=0.05cm of creation, align=center]{production de l'information\\(communiqué, photo, mesure)};
\node[etape] (diffusion) at (5.6,3.2) {2. Diffusion};
\node[descr, below=0.05cm of diffusion, align=center]{transmission au public\\(web, radio, réseaux sociaux)};
\node[etape] (reception) at (5.6,0) {3. Réception};
\node[descr, below=0.05cm of reception, align=center]{lecture, interprétation,\\relais par les destinataires};
\node[etape] (archivage) at (0,0) {4. Archivage};
\node[descr, above=0.05cm of archivage, align=center]{conservation pour\\consultation future};
\draw[fleche] (creation) -- (diffusion);
\draw[fleche] (diffusion) -- (reception);
\draw[fleche] (reception) -- (archivage);
\draw[fleche] (archivage) -- (creation)
    node[midway, left, descr, align=center]{réutilisation : une archive\\peut nourrir une\\nouvelle information};
\end{tikzpicture}
\end{center}
\legende{Le cycle de vie de l'information : création, diffusion, réception, archivage. Une information archivée peut resservir plus tard pour créer une nouvelle information (par exemple un article qui cite un ancien communiqué).}
\end{popfigure}

Pourquoi ce cycle est-il utile dans la vie quotidienne ? Parce que \textbf{chaque étape peut introduire une déformation}. À la création, l'auteur peut se tromper ou mentir. À la diffusion, un intermédiaire peut raccourcir le message. À la réception, le lecteur peut mal comprendre. À l'archivage, une vieille information peut ressortir hors de son contexte — par exemple une photo d'inondation de 2019 présentée en 2025 comme un événement actuel. Savoir à quelle étape on se trouve aide à poser les bonnes questions.

\begin{savaistu}[Le mot « infox »]
Le terme \cle{infox} est une contraction des mots \emph{information} et \emph{intoxication}. Il désigne une \textbf{fausse information diffusée volontairement} pour tromper le public : faire paniquer, salir une personne, influencer une élection ou vendre un produit. On parle aussi de « fake news » en anglais. Au Cameroun comme partout dans le monde, les infox circulent très vite sur les réseaux sociaux, surtout en période de crise ou d'élection. La meilleure arme contre l'infox, c'est l'esprit critique... que tu es en train de développer !
\end{savaistu}

\begin{exemplebox}[Un constat qui fait réfléchir]
Dans de nombreux pays, et le Cameroun ne fait pas exception, une grande partie des citoyens s'informe désormais \textbf{d'abord par les messageries et les réseaux sociaux} (WhatsApp, Facebook, TikTok), avant même les journaux, la radio ou les sites officiels. Or, dans une application de messagerie, n'importe qui peut écrire n'importe quoi, sans aucun contrôle. Ce constat montre pourquoi il est devenu indispensable, pour chaque élève, d'apprendre à évaluer une information avant de la croire et de la partager.
\end{exemplebox}

\courssub{Exemples concrets : analyser trois messages du quotidien}

Mettons en pratique tout ce que nous venons d'apprendre sur trois messages typiques de la vie numérique camerounaise.

\textbf{Message A — Communiqué du MINCOMMERCE.} « Le Ministère du Commerce informe l'opinion publique que les prix des produits de première nécessité restent inchangés. Toute information contraire est une rumeur. » Données : des prix, un ministère. Contexte : site officiel, daté et signé. Sens : rassurer et informer. Il s'agit d'une \textbf{information de source primaire}, à forte fiabilité.

\textbf{Message B — Post Facebook d'un particulier.} « Urgent !!! L'essence double demain, mon voisin qui travaille à la SCDP me l'a confirmé, partagez !!! » Ici, la donnée existe (une affirmation sur les prix), mais le contexte est flou (qui est ce voisin ? quelle preuve ?) et le sens pousse à la panique et au partage. C'est une \textbf{source secondaire non vérifiable} : typiquement le terreau d'une infox.

\textbf{Message C — Article de journal en ligne.} « D'après un communiqué du MINCOMMERCE publié ce jour, aucune hausse du carburant n'est prévue. Le ministre invite les populations au calme. » Source secondaire, mais \textbf{sérieuse} : le journaliste cite sa source primaire, date son article et signe de son nom. On peut vérifier en consultant directement le communiqué.

\begin{exoresolu}[Identifier donnée, information, source et étape du cycle]
\textbf{Énoncé.} Le principal d'un collège de Douala affiche au tableau d'affichage : « Les évaluations du 1\textsuperscript{er} trimestre commenceront le lundi 17 novembre 2025 à 7 h 30. » Un élève photographie l'affiche et la publie dans le groupe WhatsApp de la classe avec le commentaire : « Les évals commencent le 17, préparez-vous ! »
\begin{enumerate}
\item Quelle est la donnée brute contenue dans ce message ?
\item Pourquoi peut-on parler d'information ?
\item Quelle est la source primaire ? Quelle est la source secondaire ?
\item À quelles étapes du cycle de vie de l'information correspondent l'affichage au tableau et la publication WhatsApp ?
\end{enumerate}
\tcblower
\textbf{Solution détaillée.}
\begin{enumerate}
\item La donnée brute est la date \texttt{17 novembre 2025} (ainsi que l'heure \texttt{7 h 30}) : prise seule, cette date ne dit rien.
\item On parle d'information parce que la date est \textbf{contextualisée} (évaluations du 1\textsuperscript{er} trimestre, collège, heure de début) et \textbf{porteuse de sens} pour les élèves : ils savent qu'ils doivent réviser et arriver à l'heure.
\item La \textbf{source primaire} est l'affichage officiel du principal (document émanant de l'autorité compétente). La \textbf{source secondaire} est la publication WhatsApp de l'élève, qui rapporte et commente l'information d'origine.
\item L'affichage au tableau correspond aux étapes de \textbf{création} (rédaction de la décision) et de \textbf{diffusion} (mise à disposition du public). La publication WhatsApp relève de la \textbf{réception} suivie d'un \textbf{relais} (nouvelle diffusion), et la photo conservée dans le téléphone constitue une forme d'\textbf{archivage}.
\end{enumerate}
\textbf{Conclusion.} Le message WhatsApp est fidèle ici, mais la bonne habitude reste de vérifier la source primaire : aller lire soi-même l'affiche avant de modifier son emploi du temps.
\end{exoresolu}

\begin{attention}[Les trois pièges classiques du débutant]
\begin{itemize}
\item \textbf{Croire qu'un message bien écrit est forcément vrai.} Une infox peut être rédigée dans un français parfait, avec un logo officiel copié. La forme ne garantit pas le fond.
\item \textbf{Oublier la date.} Une information vraie en 2020 peut être fausse en 2025. Toujours vérifier la date de création, pas seulement la date de réception.
\item \textbf{Confondre les niveaux.} Dire « j'ai une information » alors qu'on n'a qu'une donnée brute (« 850 ») ou qu'une opinion (« je pense que... ») est une erreur de vocabulaire... et de méthode.
\end{itemize}
\end{attention}

\begin{aretenir}[Bilan de la section]
\begin{itemize}
\item Une \textbf{information} est une donnée \textbf{traitée, contextualisée et porteuse de sens} pour un destinataire : Information = Donnée + Contexte + Sens.
\item On distingue la \textbf{donnée brute} (fait élémentaire), l'\textbf{information} (donnée mise en contexte) et la \textbf{connaissance} (savoir construit à partir d'informations).
\item Les \textbf{sources primaires} (témoins directs, documents officiels) sont de première main ; les \textbf{sources secondaires} (médias, blogs, réseaux sociaux) rapportent l'information par un intermédiaire.
\item L'information suit un \textbf{cycle de vie} en quatre étapes : création, diffusion, réception, archivage ; chaque étape peut la déformer.
\item La \textbf{popularité d'un message ne prouve pas sa fiabilité} : c'est la qualité de la source qui compte.
\end{itemize}
Maintenant que tu sais ce qu'est une information et d'où elle vient, la prochaine étape consiste à apprendre les \textbf{six critères} qui permettent de juger si une information est bonne : c'est l'objet de la section suivante.
\end{aretenir}

\coursec{Les six critères d'une bonne information}

Dans la section précédente, nous avons vu qu'une \cle{information} est un renseignement qui nous apprend quelque chose et qui peut circuler sous de nombreuses formes : texte, image, son, vidéo, message WhatsApp, article de journal, affiche, etc. Mais toutes les informations ne se valent pas ! Imagine que tu prépares un exposé sur la réforme de l'examen du BEPC : un camarade te dit « j'ai entendu dire que l'épreuve de maths sera supprimée ». Dois-tu le croire ? Pour répondre à cette question, les spécialistes de l'information ont défini \cle{six critères} qui permettent de juger la qualité d'une information. C'est cette « grille de contrôle qualité » que nous allons étudier maintenant, critère par critère.

\begin{propriete}[Règle fondamentale des six critères]
Une information de qualité satisfait \textbf{simultanément} les six critères : \cle{exactitude}, \cle{pertinence}, \cle{actualité}, \cle{fiabilité}, \cle{objectivité} et \cle{accessibilité}. L'absence d'\textbf{un seul} de ces critères affaiblit la confiance que l'on peut accorder à l'information, comme une chaîne dont un seul maillon cassé ne retient plus rien.
\end{propriete}

\courssub{L'exactitude : l'information est-elle conforme aux faits ?}

\textbf{Intuition.} Avant de croire une information, la première question à se poser est : « Est-ce \emph{vrai} ? » Une information peut être bien écrite, récente, diffusée par un grand site... et pourtant fausse. L'exactitude est le critère de la \emph{conformité à la réalité}.

\begin{definition}[Exactitude]
L'\cle{exactitude} est la conformité d'une information aux \textbf{faits vérifiables}. Une information est exacte si elle peut être confirmée par des observations, des mesures, des documents officiels ou plusieurs sources indépendantes qui disent la même chose.
\end{definition}

\textbf{Comment vérifier l'exactitude en pratique ?}
\begin{enumerate}
\item Repérer les affirmations précises du contenu (chiffres, dates, noms, lieux).
\item Comparer avec au moins \textbf{deux autres sources} indépendantes (un autre journal, un site officiel, un document administratif).
\item Vérifier les chiffres : « 90 \% des élèves ont réussi » — d'où vient ce chiffre ? Qui l'a calculé ?
\item Se méfier des formulations vagues : « on dit que... », « tout le monde sait que... » ne sont pas des preuves.
\end{enumerate}

\begin{exemplebox}[Exact ou inexact ?]
Message reçu sur WhatsApp : « Le ministère des Enseignements secondaires a annoncé que le BEPC 2026 aura lieu en mars. » Pour vérifier son exactitude, on consulte le site officiel du MINESEC ou le communiqué lu à la radio nationale (CRTV). Si aucun document officiel ne confirme cette date, l'information est \textbf{inexacte} ou du moins \textbf{non vérifiée}.
\end{exemplebox}

\courssub{La pertinence : l'information répond-elle à mon besoin ?}

\textbf{Intuition.} Une information peut être parfaitement exacte mais totalement \emph{inutile} pour toi. Si tu cherches les dates du BEPC et qu'on te donne l'historique complet du système éducatif camerounais depuis 1960, tu perds ton temps : l'information n'est pas \emph{pertinente}.

\begin{definition}[Pertinence]
La \cle{pertinence} est l'adéquation d'une information au \textbf{besoin} ou à la \textbf{question} de l'utilisateur. Une information est pertinente si elle répond directement à ce que l'on cherche, sans détails inutiles ni hors-sujet.
\end{definition}

\textbf{Pour juger la pertinence}, pose-toi trois questions :
\begin{itemize}
\item Cette information répond-elle à \emph{ma} question précise ?
\item Est-elle adaptée à \emph{mon} niveau et à \emph{mon} objectif (exposé, décision, simple curiosité) ?
\item Me fait-elle gagner du temps ou m'en fait-elle perdre ?
\end{itemize}

\begin{attention}[Exact ne veut pas dire pertinent]
Un article exact sur la météo de Douala est inutile si tu prépares un voyage à Ngaoundéré. Avant même de lire un document, lis son \textbf{titre} et son \textbf{résumé} pour décider s'il mérite ton attention. C'est un réflexe d'élève efficace !
\end{attention}

\courssub{L'actualité : l'information est-elle fraîche ?}

\textbf{Intuition.} Certains renseignements « périment » comme le lait : un horaire de train de 2015, un classement de football de la saison passée, une loi abrogée... D'autres restent vrais très longtemps (la formule de l'aire d'un triangle, par exemple). L'actualité mesure ce degré de « fraîcheur ».

\begin{definition}[Actualité]
L'\cle{actualité} est le degré de fraîcheur d'une information par rapport à l'événement qu'elle décrit. Elle s'apprécie en comparant la \textbf{date de publication} de l'information à la \textbf{date où l'on s'en sert} et à la \textbf{vitesse d'évolution} du sujet.
\end{definition}

\begin{exemplebox}[Des durées de vie très différentes]
\begin{itemize}
\item Résultat d'un match joué hier : actualité excellente aujourd'hui, moyenne dans un mois.
\item Prix du litre d'essence à la pompe : change souvent, il faut une information très récente.
\item Définition d'un triangle rectangle : information stable, l'actualité importe peu.
\end{itemize}
\end{exemplebox}

\textbf{Réflexe pratique :} cherche toujours la \textbf{date de publication} ou de \textbf{dernière mise à jour} d'une page web ou d'un article. Une page sans date est suspecte.

\courssub{La fiabilité : peut-on faire confiance à la source ?}

\textbf{Intuition.} Si un inconnu dans la rue te dit « ton école est fermée demain », tu vérifies auprès du principal. Si c'est le principal lui-même qui l'annonce, tu le crois. La valeur d'une information dépend énormément de \emph{celui qui la donne} : c'est la fiabilité de la source.

\begin{definition}[Fiabilité]
La \cle{fiabilité} est la crédibilité de la \textbf{source} d'une information. Elle dépend de trois éléments principaux :
\begin{itemize}
\item l'\textbf{expertise} : l'auteur connaît-il bien le sujet (médecin pour la santé, ingénieur pour les ponts, enseignant pour les programmes scolaires) ?
\item l'\textbf{indépendance} : l'auteur a-t-il un intérêt caché (vendre un produit, défendre un parti) ?
\item l'\textbf{historique} : cette source s'est-elle trompée ou a-t-elle déjà menti par le passé ?
\end{itemize}
\end{definition}

\begin{center}
\begin{tabularx}{\linewidth}{|l|X|}
\hline
\rowcolor{popblueL} \textbf{Type de source} & \textbf{Fiabilité habituelle} \\
\hline
Site officiel (MINESEC, ministères, ONU) & Élevée : informations vérifiées et responsables \\
\hline
Journal professionnel reconnu (CRTV, quotidiens sérieux) & Bonne : journalistes identifiés, sources citées \\
\hline
Encyclopédie, manuel scolaire & Bonne : contenu relu par des spécialistes \\
\hline
Blog ou page personnelle & Moyenne : dépend de l'auteur, à vérifier \\
\hline
Message anonyme transféré sur WhatsApp & Faible : auteur inconnu, aucune garantie \\
\hline
\end{tabularx}
\end{center}

\courssub{L'objectivité : l'information est-elle présentée sans parti pris caché ?}

\textbf{Intuition.} Deux journaux peuvent raconter le même match : l'un écrit « notre glorieuse équipe a écrasé l'adversaire », l'autre écrit « l'équipe A a battu l'équipe B par 2 buts à 1 ». Les faits sont les mêmes, mais la \emph{présentation} diffère. L'objectivité concerne la manière de présenter les faits.

\begin{definition}[Objectivité]
L'\cle{objectivité} est la présentation \textbf{équilibrée} d'une information, \textbf{sans parti pris caché} : les faits sont séparés des opinions, les différents points de vue sont mentionnés, et le vocabulaire reste neutre (pas de mots exagérés ou insultants).
\end{definition}

\textbf{Indices d'un manque d'objectivité :}
\begin{itemize}
\item adjectifs élogieux ou insultants (« magnifique », « scandaleux », « honteux ») ;
\item un seul point de vue présenté, les opposants ignorés ou ridiculisés ;
\item mélange permanent entre faits et opinions sans distinction claire ;
\item publicité déguisée en information.
\end{itemize}

\begin{attention}[Objectivité n'est pas neutralité absolue]
Ne confonds pas \emph{objectivité} et \emph{neutralité}. Un \textbf{éditorial} ou une \textbf{tribune} assume ouvertement un point de vue : c'est acceptable \textbf{à condition de le signaler clairement} (mention « opinion », « éditorial », signature de l'auteur). Le problème n'est pas d'avoir une opinion, mais de la \emph{cacher} derrière une fausse apparence de fait. Un article de presse classique, lui, doit rester objectif.
\end{attention}

\courssub{L'accessibilité : l'information est-elle claire et disponible ?}

\textbf{Intuition.} La meilleure information du monde ne sert à rien si tu ne peux pas la lire : rédigée dans une langue que tu ne comprends pas, publiée dans un format que ton téléphone n'ouvre pas, ou écrite avec un vocabulaire de spécialistes incompréhensible. C'est le critère d'accessibilité.

\begin{definition}[Accessibilité]
L'\cle{accessibilité} est la facilité avec laquelle l'utilisateur peut \textbf{obtenir} et \textbf{comprendre} une information. Elle concerne :
\begin{itemize}
\item la \textbf{clarté du langage} : phrases simples, vocabulaire adapté au public ;
\item la \textbf{langue} utilisée : français, anglais, langues nationales... ;
\item le \textbf{format du support} : fichier lisible, image nette, son audible ;
\item la \textbf{disponibilité} : l'information est-elle facile à trouver et à consulter ?
\end{itemize}
\end{definition}

\begin{exemplebox}[Une information exacte mais inaccessible]
Un communiqué officiel sur les examens publié uniquement en anglais, dans un fichier illisible sur les téléphones anciens, est exact et fiable... mais \textbf{inaccessible} pour un élève francophone de 3ème sans ordinateur. Le critère d'accessibilité n'est pas satisfait.
\end{exemplebox}

\courssub{Vue d'ensemble : la check-list des six critères}

Pour retenir et utiliser facilement les six critères, on les rassemble dans une \textbf{check-list} (liste de contrôle) que l'on peut afficher en classe ou garder dans son cahier. La figure ci-dessous la résume visuellement.

\begin{popfigure}
\begin{center}
\begin{tikzpicture}[scale=1, every node/.style={transform shape}]
% Titre
\node[font=\bfseries\large, text=popdark] at (4.5,6.4) {Check-list : les 6 critères d'une bonne information};
% Ligne 1
\draw[fill=popblueL, draw=popblue, rounded corners=4pt, thick] (0,4.6) rectangle (4.3,6);
\node[font=\bfseries, text=popblue] at (2.15,5.62) {1. EXACTITUDE};
\node[font=\footnotesize, text=popdark, align=center] at (2.15,5.05) {Est-ce conforme\\aux faits vérifiables ?};
\draw[fill=popgreenL, draw=popgreen, rounded corners=4pt, thick] (4.7,4.6) rectangle (9,6);
\node[font=\bfseries, text=popgreen!60!black] at (6.85,5.62) {2. PERTINENCE};
\node[font=\footnotesize, text=popdark, align=center] at (6.85,5.05) {Répond-elle à\\mon besoin ?};
% Ligne 2
\draw[fill=poporangeL, draw=poporange, rounded corners=4pt, thick] (0,2.9) rectangle (4.3,4.3);
\node[font=\bfseries, text=poporange!70!black] at (2.15,3.92) {3. ACTUALITÉ};
\node[font=\footnotesize, text=popdark, align=center] at (2.15,3.35) {Est-elle récente\\et à jour ?};
\draw[fill=poppurpleL, draw=poppurple, rounded corners=4pt, thick] (4.7,2.9) rectangle (9,4.3);
\node[font=\bfseries, text=poppurple] at (6.85,3.92) {4. FIABILITÉ};
\node[font=\footnotesize, text=popdark, align=center] at (6.85,3.35) {La source est-elle\\crédible ?};
% Ligne 3
\draw[fill=poppinkL, draw=poppink, rounded corners=4pt, thick] (0,1.2) rectangle (4.3,2.6);
\node[font=\bfseries, text=poppink!70!black] at (2.15,2.22) {5. OBJECTIVITÉ};
\node[font=\footnotesize, text=popdark, align=center] at (2.15,1.65) {Est-elle présentée\\sans parti pris caché ?};
\draw[fill=popgoldL, draw=popgold, rounded corners=4pt, thick] (4.7,1.2) rectangle (9,2.6);
\node[font=\bfseries, text=popgold!60!black] at (6.85,2.22) {6. ACCESSIBILITÉ};
\node[font=\footnotesize, text=popdark, align=center] at (6.85,1.65) {Puis-je la lire et\\la comprendre ?};
% Cases à cocher
\foreach \y in {5.3, 3.6, 1.9}{
  \draw[draw=popdark, thick] (-0.55,\y-0.12) rectangle (-0.25,\y+0.18);
  \draw[draw=popdark, thick] (9.15,\y-0.12) rectangle (9.45,\y+0.18);
}
\node[font=\footnotesize\itshape, text=popmuted, align=center] at (4.5,0.5) {Coche chaque critère : si une case reste vide, méfie-toi de l'information !};
\end{tikzpicture}
\end{center}
\legende{Check-list visuelle des six critères d'une bonne information, à afficher en classe.}
\end{popfigure}

\begin{methode}[La grille « C.R.A.F.O.A » pour évaluer toute information]
Pour ne rien oublier, retiens l'acronyme \cle{C.R.A.F.O.A}, formé des initiales des six critères (dans l'ordre : Conformité/exactitude, Réponse au besoin, Actualité, Fiabilité, Objectivité, Accessibilité). Face à \emph{tout} contenu (article, message, vidéo), applique la démarche suivante :
\begin{enumerate}
\item \textbf{C — Exactitude} : les faits sont-ils vérifiables ? Confirmés par d'autres sources ?
\item \textbf{R — Pertinence} : le contenu répond-il à ma question ?
\item \textbf{A — Actualité} : quelle est la date de publication ? Le sujet évolue-t-il vite ?
\item \textbf{F — Fiabilité} : qui est l'auteur ? Est-il expert, indépendant, connu ?
\item \textbf{O — Objectivité} : faits et opinions sont-ils séparés ? Y a-t-il un parti pris caché ?
\item \textbf{A — Accessibilité} : le langage, la langue et le format me conviennent-ils ?
\end{enumerate}
Attribue ensuite une note sur 10 à chaque critère : tu obtiens un profil complet de l'information.
\end{methode}

\courssub{Noter une information : le graphique radar}

Une fois les six critères notés sur 10, on peut représenter le résultat sur un \textbf{graphique radar} : chaque axe correspond à un critère, et plus le point est éloigné du centre, meilleure est la note. La forme obtenue montre d'un coup d'œil les points forts et les points faibles de l'information.

\begin{exemplebox}[Exemple chiffré : un article sur la réforme du bac]
Un article publié \emph{hier} annonce une réforme du baccalauréat. Évaluation avec la grille C.R.A.F.O.A :
\begin{itemize}
\item \textbf{Exactitude} : 5/10 — les faits sont plausibles mais non confirmés ailleurs ;
\item \textbf{Pertinence} : 8/10 — le sujet concerne directement les élèves ;
\item \textbf{Actualité} : 9/10 — article publié hier, sujet d'actualité brûlante ;
\item \textbf{Fiabilité} : 3/10 — \textbf{aucune citation de source officielle}, auteur inconnu ;
\item \textbf{Objectivité} : 6/10 — ton alarmiste mais faits et opinions à peu près séparés ;
\item \textbf{Accessibilité} : 8/10 — français clair, article lisible sur téléphone.
\end{itemize}
\textbf{Conclusion :} malgré une excellente actualité, la \textbf{fiabilité très faible} (3/10) ruine la confiance : on ne doit pas diffuser cet article sans vérification auprès d'une source officielle.
\end{exemplebox}

\begin{popfigure}
\begin{center}
\begin{tikzpicture}[scale=0.85]
% Grilles concentriques (niveaux 2,4,6,8,10)
\foreach \r in {0.8, 1.6, 2.4, 3.2, 4.0}{
  \draw[popline, thin] (90:\r) \foreach \a in {150,210,270,330,30}{ -- (\a:\r) } -- cycle;
}
% Axes
\foreach \a in {90,150,210,270,330,30}{
  \draw[popline] (0,0) -- (\a:4.0);
}
% Étiquettes des axes
\node[font=\footnotesize\bfseries, text=popblue] at (90:4.7) {Exactitude};
\node[font=\footnotesize\bfseries, text=popgreen!60!black] at (150:4.9) {Pertinence};
\node[font=\footnotesize\bfseries, text=poporange!70!black] at (210:4.9) {Actualité};
\node[font=\footnotesize\bfseries, text=poppurple] at (270:4.7) {Fiabilité};
\node[font=\footnotesize\bfseries, text=poppink!70!black] at (330:4.9) {Objectivité};
\node[font=\footnotesize\bfseries, text=popgold!60!black] at (30:4.9) {Accessibilité};
% Échelle
\foreach \r/\n in {0.8/2, 1.6/4, 2.4/6, 3.2/8, 4.0/10}{
  \node[font=\tiny, text=popmuted] at (98:\r) {\n};
}
% Polygone des notes : Exact 5, Pert 8, Actu 9, Fiab 3, Obj 6, Acc 8
% 1 unité = 0.4 cm ; angles : Exact 90, Pert 150, Actu 210, Fiab 270, Obj 330, Acc 30
\draw[draw=popink, very thick, fill=popblue, fill opacity=0.25]
  (90:2.0) -- (150:3.2) -- (210:3.6) -- (270:1.2) -- (330:2.4) -- (30:3.2) -- cycle;
% Points
\foreach \a/\r in {90/2.0, 150/3.2, 210/3.6, 270/1.2, 330/2.4, 30/3.2}{
  \fill[popink] (\a:\r) circle (2.5pt);
}
% Notes affichées
\node[font=\scriptsize\bfseries, text=popink] at (90:2.45) {5};
\node[font=\scriptsize\bfseries, text=popink] at (150:3.65) {8};
\node[font=\scriptsize\bfseries, text=popink] at (210:4.05) {9};
\node[font=\scriptsize\bfseries, text=popink] at (270:0.75) {3};
\node[font=\scriptsize\bfseries, text=popink] at (330:2.85) {6};
\node[font=\scriptsize\bfseries, text=popink] at (30:3.65) {8};
\end{tikzpicture}
\end{center}
\legende{Graphique radar de l'article sur la réforme du bac : le creux vers « Fiabilité » (3/10) est immédiatement visible.}
\end{popfigure}

\begin{aretenir}[L'essentiel des six critères]
\begin{itemize}
\item \textbf{Exactitude} : conforme aux faits vérifiables.
\item \textbf{Pertinence} : répond à mon besoin précis.
\item \textbf{Actualité} : suffisamment récente pour le sujet.
\item \textbf{Fiabilité} : source experte, indépendante et sérieuse.
\item \textbf{Objectivité} : présentation équilibrée, sans parti pris caché.
\item \textbf{Accessibilité} : langage clair, langue et format adaptés.
\end{itemize}
Les six critères fonctionnent \textbf{ensemble} : un seul critère manquant suffit à rendre une information suspecte. L'acronyme \cle{C.R.A.F.O.A} et le graphique radar sont tes deux outils d'évaluation.
\end{aretenir}

\begin{exoresolu}[Évaluer une annonce de concours]
\textbf{Énoncé.} Amina lit sur une page Facebook anonyme : « Concours d'entrée dans les écoles de la fonction publique : inscription gratuite cette année, envoyez juste 2 000 FCFA de frais de dossier par mobile money au numéro ci-dessous. » Applique la grille C.R.A.F.O.A et conclus.
\tcblower
\textbf{Solution détaillée.}
\begin{itemize}
\item \textbf{C — Exactitude (1/10)} : « inscription gratuite » mais « 2 000 FCFA de frais » : contradiction interne, signe classique d'arnaque.
\item \textbf{R — Pertinence (6/10)} : le sujet intéresse Amina, mais l'annonce ne précise ni l'école ni le concours concerné.
\item \textbf{A — Actualité (5/10)} : aucune date de publication ni date limite indiquée.
\item \textbf{F — Fiabilité (0/10)} : page anonyme, aucun lien avec un site officiel ; les vrais concours sont publiés par communiqué officiel (presse, radio, sites des ministères).
\item \textbf{O — Objectivité (2/10)} : ton pressant (« envoyez juste... ») destiné à pousser à payer vite.
\item \textbf{A — Accessibilité (7/10)} : message clair et lisible — c'est justement ce qui le rend dangereux.
\end{itemize}
\textbf{Conclusion justifiée :} la fiabilité est nulle et l'exactitude contradictoire. C'est très probablement une \textbf{escroquerie}. Amina ne doit rien payer et doit vérifier tout concours sur les canaux officiels (communiqués ministériels, CRTV, affichage des établissements).
\end{exoresolu}

\begin{exercice}[À toi de jouer]
Ton oncle te transfère une vidéo affirmant qu'un fruit local guérit toutes les maladies. Note chacun des six critères C.R.A.F.O.A sur 10, justifie chaque note en une phrase, puis rédige une conclusion de trois lignes indiquant si l'on peut croire et partager cette vidéo.
\end{exercice}

\begin{savaistu}[Le fact-checking au Cameroun]
De grands médias et des organisations pratiquent le \emph{fact-checking} : la vérification systématique des informations qui circulent, surtout en période d'élections ou de crise sanitaire. Au Cameroun comme ailleurs en Afrique, des plateformes spécialisées démontent chaque semaine des rumeurs diffusées sur les réseaux sociaux. En appliquant la grille C.R.A.F.O.A, tu deviens toi-même un petit « vérificateur d'informations » : c'est une compétence précieuse pour le BEPC... et pour toute ta vie de citoyen !
\end{savaistu}

\coursec{Démarche pratique d'évaluation d'une information}

Dans la section précédente, nous avons découvert les \cle{six critères} qui permettent de juger la qualité d'une information : la \cle{fiabilité} de la source, l'\cle{exactitude}, l'\cle{actualité}, l'\cle{objectivité}, la \cle{pertinence} et la \cle{clarté}. Mais connaître les critères ne suffit pas : encore faut-il savoir les \emph{appliquer}, dans le bon ordre, face à une information concrète — un message WhatsApp, un post Facebook, un article de blog. C'est précisément l'objet de cette section : te donner une \cle{démarche} simple, en cinq étapes, que tu pourras réutiliser toute ta vie, comme un réflexe, avant de croire ou de partager quoi que ce soit.

\courssub{Pourquoi une démarche en étapes ?}

Imagine la situation suivante : tu reçois sur WhatsApp un message affirmant que « le BEPC 2026 est reporté au mois de septembre ». Ton cœur s'accélère. Deux réactions sont possibles : le \emph{réflexe émotionnel} (partager immédiatement à tous tes groupes) ou le \emph{réflexe méthodique} (vérifier d'abord). La démarche en cinq étapes transforme la vérification en une routine simple, un peu comme l'algorithme d'une recette de cuisine : on suit les étapes dans l'ordre, on note les résultats, et la conclusion s'impose presque d'elle-même.

\begin{definition}[Démarche d'évaluation d'une information]
La \cle{démarche d'évaluation d'une information} est une suite ordonnée d'étapes permettant d'examiner une information selon les six critères d'une bonne information, de lui attribuer une \cle{note globale} chiffrée, puis de prendre une \cle{décision} raisonnée : partager, archiver, signaler ou ignorer.
\end{definition}

\begin{popfigure}
\begin{center}
\begin{tikzpicture}[
  font=\small,
  etape/.style={rectangle, rounded corners=3pt, draw=popblue, fill=popblueL, thick, text width=4.6cm, align=center, minimum height=1cm},
  decision/.style={diamond, draw=poporange, fill=poporangeL, thick, aspect=2.2, align=center, text width=2.6cm},
  action/.style={rectangle, draw=popgreen, fill=popgreenL, thick, rounded corners=3pt, text width=2.5cm, align=center, minimum height=0.8cm},
  fleche/.style={-{Stealth[length=2.5mm]}, thick, popdark}
]
\node[etape, align=center] (e1){\textbf{Étape 1}\\ Identifier la source (auteur, média, date, URL)};
\node[etape, below=0.7cm of e1, align=center] (e2){\textbf{Étape 2}\\ Vérifier l'exactitude par recoupement};
\node[etape, below=0.7cm of e2, align=center] (e3){\textbf{Étape 3}\\ Analyser le ton et le vocabulaire};
\node[etape, below=0.7cm of e3, align=center] (e4){\textbf{Étape 4}\\ Noter chaque critère sur 10, calculer la note globale};
\node[decision, below=0.9cm of e4] (d1) {Note globale $\geq 5/10$ ?};
\node[action, below left=0.9cm and 0.2cm of d1] (a1) {Partager ou archiver};
\node[action, below right=0.9cm and 0.2cm of d1] (a2) {Signaler ou ignorer};
\draw[fleche] (e1) -- (e2);
\draw[fleche] (e2) -- (e3);
\draw[fleche] (e3) -- (e4);
\draw[fleche] (e4) -- (d1);
\draw[fleche] (d1.west) -| node[pos=0.3, above left]{oui} (a1.north);
\draw[fleche] (d1.east) -| node[pos=0.3, above right]{non} (a2.north);
\end{tikzpicture}
\end{center}
\legende{Organigramme de la démarche d'évaluation d'une information en cinq étapes, avec la décision finale.}
\end{popfigure}

\courssub{Étape 1 : identifier la source}

La première question à se poser est toujours : \emph{qui me parle ?} Une information sans auteur identifiable est comme une lettre sans signature : on ne peut ni lui faire confiance, ni demander des comptes à son auteur.

\begin{methode}[Identifier la source d'une information]
\begin{enumerate}
  \item Repérer l'\cle{auteur} : nom d'une personne, d'un journaliste, d'une administration (MINESEC, MINPOSTEL), d'une entreprise. Si l'auteur est anonyme ou introuvable, c'est un mauvais signe.
  \item Repérer le \cle{média} : site d'un journal connu, page officielle certifiée (badge de vérification), blog personnel, simple capture d'écran relayée.
  \item Repérer la \cle{date} de publication : une information vraie en 2019 peut être fausse ou dépassée aujourd'hui.
  \item Examiner l'\cle{URL} (adresse du site) : les fausses informations imitent souvent les vraies adresses (\texttt{minedub-gov.com} au lieu d'une adresse officielle en \texttt{.gov.cm}).
\end{enumerate}
\end{methode}

\begin{popfigure}
\begin{center}
\begin{tikzpicture}[font=\small]
% Cadre du post
\draw[thick, popdark, fill=white, rounded corners=4pt] (0,0) rectangle (10.5,6.2);
\draw[popblue, thick, fill=popblueL] (0,5.4) rectangle (10.5,6.2);
\node[anchor=west, text=white, fill=popblue, rounded corners=2pt, inner sep=2pt] at (0.3,5.8) {\textbf{Facebook}};
% Avatar et nom
\draw[fill=poppinkL, draw=poppink] (0.7,4.7) circle (0.35);
\node at (0.7,4.7) {\textbf{?}};
\node[anchor=west] at (1.3,4.9) {\textbf{Infos Buzz Camer 237}};
\node[anchor=west, text=popmuted] at (1.3,4.45) {\footnotesize Page sans badge -- date de publication introuvable};
% Texte du post
\node[anchor=north west, text width=9.6cm, align=left] at (0.3,4.0)
{\textbf{URGENT !!!} Le BEPC 2026 est ANNULÉ à cause des grèves. Partagez vite à tous les candidats avant que le ministère ne supprime l'info !!! Cliquez ici : \texttt{bepc-infos237.xyz}};
% Image factice
\draw[fill=popcream, draw=popline, dashed] (0.4,0.9) rectangle (5.4,2.2);
\node[text=popmuted] at (2.9,1.55) {\footnotesize [image floue : logo officiel recopié]};
% Réactions
\node[anchor=west, text=popmuted] at (5.8,1.55) {\footnotesize 12 458 partages \quad 8 900 mentions J'aime};
% Annotations
\draw[-{Stealth}, poporange, thick] (0.7,4.35) -- (0.7,3.5) node[below, text=poporange, align=center, font=\footnotesize] {auteur\\ anonyme};
\draw[-{Stealth}, poporange, thick] (8.6,3.1) -- (9.6,2.6) node[right=0pt, text=poporange, align=left, font=\footnotesize, text width=3.2cm] {URL suspecte\\ (pas de \texttt{.gov.cm})};
\draw[-{Stealth}, poporange, thick] (8.2,1.3) -- (9.6,0.7) node[right=0pt, text=poporange, align=left, font=\footnotesize, text width=3.2cm] {beaucoup de partages\\ $\neq$ information vraie};
\draw[-{Stealth}, poporange, thick] (2.9,0.9) -- (2.9,0.2) node[below, text=poporange, align=center, font=\footnotesize] {image sans source};
\end{tikzpicture}
\end{center}
\legende{Capture d'écran reconstituée et annotée d'un post Facebook suspect : auteur anonyme, URL douteuse, ton alarmiste, image sans origine.}
\end{popfigure}

\courssub{Étape 2 : vérifier l'exactitude par recoupement}

Une fois la source identifiée, il faut vérifier si les \emph{faits} annoncés sont exacts. Le principe est celui du \cle{recoupement} : une information sérieuse est confirmée par plusieurs sources fiables et indépendantes.

\begin{propriete}[Règle du recoupement]
Une information est considérée comme \cle{exacte} si elle est confirmée par au moins \textbf{deux sources fiables indépendantes} (par exemple un site officiel et un média reconnu), ou par une \textbf{donnée officielle} consultable directement (communiqué du MINESEC, site d'une administration camerounaise).
\end{propriete}

Concrètement, pour notre message sur le BEPC : on consulte le site ou la page officielle du MINESEC, on écoute la radio nationale (CRTV), on regarde si des journaux reconnus (\emph{Cameroon Tribune}, \emph{Le Jour}) en parlent. Si aucune source fiable ne confirme, l'exactitude est très faible.

\courssub{Étape 3 : analyser le ton et le vocabulaire}

Le ton d'un texte en dit long sur ses intentions. Une bonne information est \cle{objective} : elle rapporte des faits et sépare clairement les opinions. Une information douteuse cherche souvent à provoquer une émotion (peur, colère, joie excessive) pour pousser au partage.

Les signaux d'alerte à repérer :
\begin{itemize}
  \item les mots en majuscules et les points d'exclamation répétés (« URGENT !!! », « SCANDALE ») ;
  \item les appels à partager rapidement (« partagez avant que ce soit supprimé ») ;
  \item les généralités sans preuve (« tout le monde sait que... », « on nous cache la vérité ») ;
  \item les \cle{opinions non signalées} présentées comme des faits (« ce ministre est le pire » au lieu de « selon tel syndicat, ce ministre... ») ;
  \item les fautes d'orthographe nombreuses dans un texte qui prétend être officiel.
\end{itemize}

\courssub{Étape 4 : noter chaque critère et calculer la note globale}

Pour rendre l'évaluation objective, on attribue à chacun des six critères une \cle{note sur 10} (0 = totalement mauvais, 10 = parfait), puis on calcule la \cle{note globale} comme la moyenne des six notes :
\[
\text{Note globale} = \frac{\text{Fiabilité} + \text{Exactitude} + \text{Actualité} + \text{Objectivité} + \text{Pertinence} + \text{Clarté}}{6}
\]

\begin{methode}[Utiliser la grille C.R.A.F.O.A fournie en annexe]
La fiche « Grille \cle{C.R.A.F.O.A} » (Crédibilité de la source, Recoupement, Actualité, Fiabilité des faits, Objectivité, Adéquation au besoin) permet de noter systématiquement chaque information :
\begin{enumerate}
  \item Recopier le titre de l'information et sa source en haut de la grille.
  \item Pour chaque ligne (critère), entourer la note de 0 à 10 en suivant les étapes 1 à 3 de la démarche.
  \item Additionner les six notes puis diviser par 6 pour obtenir la note globale.
  \item Reporter la note globale dans la case « décision » et cocher l'action choisie.
\end{enumerate}
Utiliser toujours la même grille évite les oublis et permet de comparer deux informations entre elles.
\end{methode}

\begin{exemplebox}[Exemple chiffré : le message viral « Le BEPC 2026 annulé »]
Un message circule massivement sur les réseaux sociaux : « Le BEPC 2026 est annulé ». Appliquons la grille :
\begin{center}
\begin{tabularx}{\linewidth}{|l|c|X|}
\hline
\rowcolor{popblueL} \textbf{Critère} & \textbf{Note /10} & \textbf{Justification} \\
\hline
Fiabilité de la source & 1 & Source anonyme, page non certifiée. \\
\hline
Exactitude & 1 & Aucun recoupement : ni le MINESEC, ni aucun média reconnu ne confirment. \\
\hline
Actualité & 2 & Le message reprend une rumeur datant de plusieurs années. \\
\hline
Objectivité & 2 & Ton alarmiste, appel au partage, majuscules. \\
\hline
Pertinence & 1 & Ne répond à aucun besoin vérifiable, crée la panique. \\
\hline
Clarté & 1 & Texte confus, fautes, aucune référence. \\
\hline
\end{tabularx}
\end{center}
\[
\text{Note globale} = \frac{1 + 1 + 2 + 2 + 1 + 1}{6} = \frac{8}{6} \approx 1{,}3/10
\]
Conclusion : avec une note globale de $1{,}3/10$, ce message doit être \textbf{signalé} puis \textbf{ignoré} — et surtout pas partagé.
\end{exemplebox}

\courssub{Étape 5 : prendre une décision}

La note globale éclaire la décision finale. On retiendra le guide suivant :

\begin{center}
\begin{tabularx}{\linewidth}{|c|X|}
\hline
\rowcolor{popblueL} \textbf{Note globale} & \textbf{Décision conseillée} \\
\hline
de 7 à 10 & \cle{Partager} (en citant la source) ou \cle{archiver} pour réutilisation. \\
\hline
de 5 à 6,9 & \cle{Archiver} avec prudence ; revérifier plus tard avant tout usage. \\
\hline
de 3 à 4,9 & \cle{Ignorer} : ne pas relayer, ne pas utiliser. \\
\hline
de 0 à 2,9 & \cle{Signaler} (à la plateforme, à un enseignant, à un parent) puis ignorer. \\
\hline
\end{tabularx}
\end{center}

\begin{attention}[Le piège du nombre de « likes »]
Se fier uniquement au nombre de « j'aime » ou de partages est un \cle{piège fréquent}. Une fausse information spectaculaire se partage souvent \emph{plus} vite qu'une information vraie mais banale, car elle joue sur les émotions. Les 12 458 partages de notre post suspect ne prouvent absolument pas sa véracité : ils prouvent seulement sa viralité. La note globale, elle, repose sur des critères objectifs.
\end{attention}

\begin{exoresolu}[Évaluer un message reçu sur WhatsApp]
Énoncé. Tu reçois ce message : « Info sûre : distribution gratuite de 50 000 FCFA par le gouvernement à tous les élèves de 3ème. Inscris-toi vite sur \texttt{aide-eleves237.net} avec ton numéro et celui de tes parents. » Applique la démarche en cinq étapes et conclus.
\tcblower
\textbf{Étape 1 — Source :} aucun auteur identifié, aucun média, pas de date ; l'adresse \texttt{aide-eleves237.net} n'est pas une adresse officielle (les sites de l'État camerounais utilisent \texttt{.gov.cm}). Fiabilité : 1/10.

\textbf{Étape 2 — Recoupement :} aucun communiqué du MINESEC ni de la Primature, aucun article de \emph{Cameroon Tribune} ou de la CRTV ne confirme. Exactitude : 1/10.

\textbf{Étape 3 — Ton :} « Info sûre », « vite », promesse d'argent facile, demande de données personnelles (numéros de téléphone) : signaux classiques d'une \emph{arnaque}. Objectivité : 1/10.

\textbf{Étape 4 — Notes :} actualité 2/10 (rumeur recyclée), pertinence 1/10, clarté 2/10. Note globale :
\[
\frac{1 + 1 + 1 + 2 + 1 + 2}{6} = \frac{8}{6} \approx 1{,}3/10
\]

\textbf{Étape 5 — Décision :} note inférieure à 3 : on \textbf{signale} le message (à la plateforme, à un adulte) et on l'\textbf{ignore}. Surtout, on ne communique jamais les numéros demandés.
\end{exoresolu}

\begin{savaistu}[Vérifier en ligne avec les outils de fact-checking]
De grands médias disposent d'outils de vérification rapide en ligne : par exemple « Factuel », le service de vérification de l'Agence France-Presse (AFP), qui démonte chaque semaine des rumeurs circulant en Afrique francophone, y compris au Cameroun. Tu peux consulter ces pages avec l'aide d'un adulte ou de ton enseignant pour vérifier une information douteuse : c'est un excellent complément à ta grille d'évaluation.
\end{savaistu}

\begin{aretenir}[L'essentiel de la démarche]
\begin{itemize}
  \item La démarche compte \cle{cinq étapes} : identifier la source, recouper les faits, analyser le ton, noter les six critères, décider.
  \item La \cle{note globale} est la moyenne des six notes sur 10 ; elle guide objectivement la décision.
  \item Quatre décisions possibles : \cle{partager}, \cle{archiver}, \cle{signaler}, \cle{ignorer}.
  \item La viralité (likes, partages) n'est \emph{jamais} une preuve de vérité.
\end{itemize}
\end{aretenir}

\begin{exercice}[À toi de jouer]
Un camarade te montre un article intitulé « Un nouveau réseau 6G installé secrètement à Yaoundé rend malade », publié sur un blog inconnu, sans date, sans auteur, mais partagé 3 000 fois. Applique les cinq étapes de la démarche, attribue une note à chaque critère en justifiant, calcule la note globale et rédige ta décision finale en trois lignes.
\end{exercice}

\coursec{Rédiger et justifier une évaluation critique}

Après avoir appris à \emph{repérer} les indices de qualité d'une information (Section 3), nous devons maintenant savoir \emph{communiquer} notre analyse de façon structurée, rigoureuse et convaincante. Une évaluation mentale reste subjective ; une évaluation \textbf{écrite et justifiée} devient un argument partageable, utilisable dans un rapport scolaire, un article de journal scolaire ou une décision administrative. C'est le passage du « je pense que » au « je démontre que ».

\courssub{Structure type : le canevas en trois temps}

Toute évaluation critique formelle suit une architecture immuable en trois parties. Cette structure guide le lecteur (et l'examinateur) pas à pas, de la présentation de l'objet au jugement final.

\begin{definition}[Structure type d'une évaluation critique]
Une évaluation critique standard comprend trois sections obligatoires :
\begin{enumerate}
    \item \textbf{Introduction} : Identifie l'objet analysé (titre, auteur, source, date, URL), précise le contexte (pourquoi cette info ?) et annonce la thèse globale (ex. : « Cette information est globalement fiable mais incomplète sur le plan chiffré »).
    \item \textbf{Développement (Analyse C.R.A.F.O.A)} : Examine \emph{systématiquement} les six critères. Pour chaque critère : \emph{énoncé du critère} $\rightarrow$ \emph{constat factuel (citation/preuve)} $\rightarrow$ \emph{interprétation (impact sur la fiabilité)}. Les critères sont liés par des connecteurs logiques.
    \item \textbf{Conclusion} : Synthétise les forces et faiblesses, donne une \textbf{note globale} (sur 10 ou 20) et un \textbf{jugement motivé} (ex. : « Utilisable avec prudence pour un exposé, à croiser pour un dossier officiel »).
\end{enumerate}
\end{definition}

\courssub{Les connecteurs logiques : tisser le raisonnement}

Une suite de phrases « Critère 1 : ... Critère 2 : ... » n'est pas une démonstration. Les connecteurs logiques montrent la relation \emph{entre} les critères : concession, opposition, conséquence, renforcement.

\begin{propriete}[Principaux connecteurs pour l'évaluation critique]
\begin{center}
\begin{tabularx}{\linewidth}{|l|X|l|}
\hline
\textbf{Relation logique} & \textbf{Connecteurs recommandés} & \textbf{Effet produit} \\ \hline
\rowcolor{popblueL}
Ajout / Renforcement & \texttt{De plus, En outre, Par ailleurs, Également} & Accumule les arguments dans le même sens. \\ \hline
Concession / Nuance & \texttt{Toutefois, Cependant, Néanmoins, Certes... mais} & Signale une faiblesse sans tout rejeter. \\ \hline
\rowcolor{popgreenL}
Conséquence / Déduction & \texttt{Par conséquent, Dès lors, C'est pourquoi, Ainsi} & Lie la preuve au jugement (ex. : « Source anonyme, \textbf{donc} crédibilité faible »). \\ \hline
Cause / Explication & \texttt{Car, Parce que, En effet, Puisque} & Justifie une note basse ou haute. \\ \hline
\rowcolor{poporangeL}
Comparaison / Alternative & \texttt{Alors que, Tandis que, Contrairement à} & Met en perspective (ex. : « Titre accrocheur \textbf{tandis que} contenu vide »). \\ \hline
\end{tabularx}
\end{center}
\end{propriete}

\begin{attention}[Piège : la liste de courses]
Ne rédigez \textbf{jamais} : « La source est connue. La date est récente. L'info est objective. Donc c'est bon. »
\emph{Rédigez plutôt :} « \textbf{Certes} la source est institutionnelle (MINCOMMERCE), \textbf{ce qui assure} une forte crédibilité. \textbf{Toutefois}, l'absence de date de publication \textbf{empêche} de vérifier la réactualisation des prix. \textbf{Par conséquent}, l'information ne peut être utilisée telle quelle pour une étude de marché actuelle. »
\end{attention}

\courssub{Citer ses preuves : l'ancrage dans le texte}

Une affirmation sans preuve est une opinion. En évaluation critique, \textbf{tout jugement sur un critère doit s'appuyer sur une citation explicite} (guillemets + référence précise : §, ligne, minute vidéo, capture d'écran).

\begin{methode}[Citer une preuve dans le développement]
\begin{enumerate}
    \item \textbf{Repérez} l'élément dans le document (phrase, chiffre, nom d'auteur, logo, URL).
    \item \textbf{Citez-le} textuellement entre guillemets français \guillemotleft ~ \guillemotright~ (ou en bloc indenté si > 2 lignes).
    \item \textbf{Référencez} la localisation : (l. 12-14), (§ 3), (capture d'écran Fig. 2), (min 04:12).
    \item \textbf{Interprétez} : expliquez \emph{en quoi} cet extrait prouve le respect ou la violation du critère.
\end{enumerate}
\end{methode}

\begin{exemplebox}[Exemple de citation intégrée]
\textbf{Critère : Fiabilité (vérifiabilité croisée).}
\begin{quote}
\guillemotleft Le prix de la baguette passe de 150 à 175 FCFA \guillemotright~ (l. 5-6).
\end{quote}
\textbf{Analyse :} Ce chiffre précis est \textbf{non sourcé} dans l'article (pas de référence au communiqué officiel du MINCOMMERCE ni au syndicat des boulangers). \textbf{Donc} la fiabilité est affaiblie : impossible de croiser cette donnée avec une source primaire.
\end{exemplebox}

\courssub{Neutralité et relecture : les garanties de rigueur}

\courssub{Neutralité : évaluer le contenu, pas la personne}

\begin{definition}[Neutralité de l'évaluation]
La neutralité consiste à juger \textbf{l'information} (son exactitude, sa source, sa forme) et non \textbf{l'auteur} (ses intentions, sa personnalité, son camp politique). On évite les adjectifs subjectifs (\guillemotleft mensonger \guillemotright, \guillemotleft manipulé \guillemotright, \guillemotleft incompétent \guillemotright) au profit de descripteurs techniques (\guillemotleft non sourcé \guillemotright, \guillemotleft contradictoire \guillemotright, \guillemotleft obsolète \guillemotright).
\end{definition}

\begin{attention}[Jugements de valeur non étayés — Interdits]
\begin{itemize}
    \item \textcolor{popdark}{\textbf{Interdit :}} \guillemotleft Cet article est faux parce que je ne suis pas d'accord avec la hausse. \guillemotright~ (Opinion personnelle).
    \item \textcolor{popdark}{\textbf{Interdit :}} \guillemotleft L'auteur est malhonnête. \guillemotright~ (Attaque \emph{ad hominem}).
    \item \textcolor{popgreen}{\textbf{Correct :}} \guillemotleft L'article affirme une hausse de 25\% sans citer de source officielle, ce qui rend l'information non vérifiable (critère Fiabilité non respecté). \guillemotright
\end{itemize}
\end{attention}

\courssub{Relecture : la grille de contrôle finale}

Avant de rendre votre copie, appliquez cette grille (checklist) :

\begin{aretenir}[Grille de relecture — 5 points de contrôle]
\begin{enumerate}
    \item \textbf{Cohérence structurelle :} Introduction / Développement (6 critères) / Conclusion présents ?
    \item \textbf{Couverture C.R.A.F.O.A :} Les 6 critères sont-ils traités explicitement (même pour dire « non applicable ») ?
    \item \textbf{Preuves citées :} Chaque jugement sur un critère s'appuie-t-il sur une citation localisée ?
    \item \textbf{Connecteurs :} Les phrases s'enchaînent-elles logiquement (pas de « et... et... et... ») ?
    \item \textbf{Orthographe / Syntaxe :} Zéro faute, phrases complètes, ton neutre, note globale chiffrée en conclusion.
\end{enumerate}
\end{aretenir}

\begin{center}
\begin{popfigure}
\begin{tikzpicture}[
    node distance=8mm and 12mm,
    box/.style={draw=popdark, thick, rounded corners=4pt, fill=popcream, text width=0.45\linewidth, align=center, font=\small},
    arrow/.style={-Stealth, thick, popblue},
    labelbox/.style={font=\tiny\bfseries, text=popblue, anchor=south, align=center}
]
% Paragraphe central simulé
\node[box, minimum height=7cm] (paragraphe) {
    \textbf{Introduction}\par\medskip
    L'article \guillemotleft Hausse du pain : 175 FCFA la baguette \guillemotright~ (Journal \emph{Le Quotidien}, 12/03/2025, p.3) annonce une augmentation soudaine.\par\medskip
    \textbf{Développement}\par\medskip
    \underline{Crédibilité} : Source presse écrite connue, mais auteur non signé (pigiste anonyme). \textbf{Donc} crédibilité moyenne.\par\medskip
    \underline{Réactualisation} : Date du jour (12/03). \textbf{Donc} info fraîche.\par\medskip
    \underline{Accessibilité} : Langage clair, titre explicite. \textbf{Donc} accessible.\par\medskip
    \underline{Fiabilité} : Chiffre 175 FCFA cité \guillemotleft selon des sources concordantes \guillemotright~ (l. 8) sans les nommer. \textbf{Toutefois}, pas de lien vers arrêté ministériel. \textbf{Par conséquent}, fiabilité faible.\par\medskip
    \underline{Objectivité} : Ton neutre, pas d'adjectifs péjoratifs. \textbf{Donc} objectif.\par\medskip
    \underline{Accessibilité (web)} : Article payant (paywall). \textbf{Donc} accessibilité réduite.\par\medskip
    \textbf{Conclusion}\par\medskip
    Note globale : \textbf{5,5/10}. Info alerte utile, mais non exploitable sans confirmation officielle (MINCOMMERCE).
};
% Annotations sur les côtés
\node[labelbox] at ($(paragraphe.north west)+(-1.5,0.3)$) {Critère\\C.R.A.F.O.A};
\draw[arrow] ($(paragraphe.west)+(-0.5,2.2)$) -- ++(-1.2,0) node[left, font=\tiny, text=popblue, align=right] {Introduction\\Objet + Source + Thèse};
\draw[arrow] ($(paragraphe.west)+(-0.5,1.0)$) -- ++(-1.2,0) node[left, font=\tiny, text=popgreen, align=right] {Crédibilité\\Preuve + Interprétation};
\draw[arrow] ($(paragraphe.west)+(-0.5,-0.2)$) -- ++(-1.2,0) node[left, font=\tiny, text=poporange, align=right] {Fiabilité\\Citation l.8 +\\Connecteur \texttt{Toutefois}};
\draw[arrow] ($(paragraphe.west)+(-0.5,-1.6)$) -- ++(-1.2,0) node[left, font=\tiny, text=poppurple, align=right] {Conclusion\\Note + Jugement\\Motivé};
% Légendes de couleur
\node[draw=popblue, thick, fill=popblue!10, rounded corners, below=1cm of paragraphe, font=\small] {
    \textcolor{popblue}{\textbf{■ Introduction/Conclusion}} \quad
    \textcolor{popgreen}{\textbf{■ Critère validé}} \quad
    \textcolor{poporange}{\textbf{■ Critère fragilisé (Toutefois)}} \quad
    \textcolor{poppurple}{\textbf{■ Jugement final chiffré}}
};
\end{tikzpicture}
\legende{Modèle annoté d'un paragraphe d'évaluation critique : structure, connecteurs, citations et code couleur des critères C.R.A.F.O.A.}
\end{popfigure}
\end{center}

\courssub{Application : Exemple résolu complet (Contexte camerounais)}

Voici un exemple intégral, prêt à servir de modèle pour le BEPC. L'exercice impose une longueur de 150 à 200 mots.

\begin{exoresolu}[Évaluation d'un article web : \guillemotleft Hausse du prix du pain au Cameroun : ce qu'il faut savoir \guillemotright]
\textbf{Énoncé :} Vous trouvez sur la page Facebook \guillemotleft Actu Cameroun 237 \guillemotright (partagée 2 400 fois, 14/03/2025) un texte affirmant : \guillemotleft Le gouvernement a validé hier la hausse du prix de la baguette à 200 FCFA. Les boulangers menacent de grève si le prix du sac de farine ne baisse pas. \guillemotright Aucune signature, aucun lien vers un communiqué officiel, photo d'une baguette floue en illustration. Évaluez cette information en 180 mots maximum.

\tcblower
\textbf{Solution détaillée :}

\textbf{Introduction}
L'information analysée est une publication Facebook de la page \guillemotleft Actu Cameroun 237 \guillemotright (14/03/2025) annonçant une hausse officielle du pain à 200 FCFA. L'objectif est d'en évaluer la fiabilité pour un usage citoyen.

\textbf{Développement (Analyse C.R.A.F.O.A)}
\textbf{Crédibilité} : La source est une page Facebook grand public, non institutionnelle, sans auteur identifié (pseudo uniquement). \textbf{Donc} crédibilité très faible.
\textbf{Réactualisation} : Datée du jour (14/03), l'info est \emph{a priori} fraîche. \textbf{Cependant}, l'expression \guillemotleft hier \guillemotright est floue sans heure précise.
\textbf{Accessibilité} : Format court, langue française correcte, image présente. \textbf{Donc} accessible.
\textbf{Fiabilité} : L'affirmation \guillemotleft Le gouvernement a validé \guillemotright n'est étayée par \textbf{aucun lien} vers le portail \texttt{www.spm.gov.cm} ni vers un communiqué du MINCOMMERCE. \textbf{Par conséquent}, l'information est non vérifiable (critère majeur défaillant).
\textbf{Objectivité} : Vocabulaire alarmiste (\guillemotleft menacent \guillemotright, \guillemotleft grève \guillemotright) sans citation syndicale nominale. \textbf{Toutefois}, le ton reste informatif.
\textbf{Accessibilité (web)} : Public, gratuit, viral (2 400 partages). \textbf{Donc} forte diffusion mais sans gage de qualité.

\textbf{Conclusion}
Note globale : \textbf{3,5/10}. Information \textbf{non fiable} (absence de source primaire vérifiable). \textbf{Ne pas partager} sans confirmation officielle (site MINCOMMERCE ou CAMTEL/CRTV). Risque de rumeur (fake news) élevé.
\end{exoresolu}

\begin{savaistu}[Le fact-checking au Cameroun : \texttt{DataCheck} et \texttt{StopIntox}]
Au Cameroun, des initiatives comme \textbf{DataCheck Cameroun} (projet de l'association \emph{Internet Sans Frontières}) ou la rubrique \guillemotleft Stop Intox \guillemotright de \emph{CRTV Web} appliquent \emph{exactement} cette méthode C.R.A.F.O.A pour débusquer les fausses informations virales (ex. : rumeurs sur les concours de la Fonction Publique, fausses annonces de bourses d'études). Leur charte éditoriale impose : citation de la source primaire, capture d'écran horodatée, verdict chiffré. C'est la version professionnelle de ce que vous apprenez ici !
\end{savaistu}

\courssub{Illustration : Fiche d'évaluation type (Modèle PDF projeté)}

En examen ou en TP, on vous remettra souvent une fiche pré-structurée. Voici à quoi elle ressemble une fois remplie (simulation).

\begin{center}
\begin{popfigure}
\begin{tikzpicture}[
    form/.style={draw=popdark, thick, rounded corners=3pt, fill=white, inner sep=4pt},
    header/.style={fill=popblue, text=white, font=\bfseries\small, rounded corners=3pt},
    field/.style={font=\small, text=popdark},
    label/.style={font=\tiny\bfseries, text=popmuted, anchor=west},
    note/.style={font=\large\bfseries, text=poporange, anchor=east}
]
% Cadre principal
\node[form, minimum width=\linewidth, minimum height=9.5cm] (cadre) {};
% En-tête
\node[header, minimum width=\linewidth-8pt, anchor=north] at (cadre.north) (entete) {FICHE D'ÉVALUATION CRITIQUE D'INFORMATION — CLASSE DE 3ÈME};
% Champs d'en-tête (positionnés relativement au bord gauche du cadre)
\foreach \y/\lab/\val in {
    8.5/Objet analysé :/Article FB \guillemotleft Hausse pain 200 FCFA \guillemotright — Page \texttt{Actu Cameroun 237},
    7.8/Source / Auteur :/Page FB non vérifiée — Auteur anonyme (pseudo),
    7.1/Date / URL :/14/03/2025 — \texttt{facebook.com/actu237/posts/12345},
    6.4/Contexte :/Recherche info pour exposé éco citoyenneté
} {
    \node[label] at ($(cadre.west) + (1.0, \y)$) {\lab};
    \node[field, anchor=west, text width=0.75\linewidth] at ($(cadre.west) + (3.5, \y)$) {\val};
}
% Grille C.R.A.F.O.A
\foreach \i/\crit/\score/\justif in {
    1/Crédibilité/2/{Source anonyme, non institutionnelle},
    2/Réactualisation/4/{Date jour mais \guillemotleft hier \guillemotright imprécis},
    3/Accessibilité/5/{Clair, court, image, gratuit},
    4/Fiabilité/1/{Pas de source primaire, non vérifiable},
    5/Objectivité/3/{Ton neutre mais vocabulaire alarmiste},
    6/Accessibilité Web/5/{Viral, public, paywall absent}} {
    \pgfmathsetmacro{\yy}{5.5 - \i*0.6};
    \node[label] at ($(cadre.west) + (1.0, \yy)$) {\crit~:};
    \node[field, anchor=west, text width=0.75\linewidth] at ($(cadre.west) + (3.5, \yy)$) {\justif};
    \node[note, anchor=east] at ($(cadre.east) + (-0.5, \yy)$) {\score/5};
}
% Conclusion
\node[label] at ($(cadre.west) + (1.0, 1.5)$) {NOTE GLOBALE :};
\node[font=\Large\bfseries, text=poporange] at ($(cadre.west) + (5.0, 1.5)$) {3,5 / 10};
\node[label] at ($(cadre.west) + (8.0, 1.5)$) {VERDICT :};
\node[field, anchor=west, text width=0.55\linewidth] at ($(cadre.west) + (10.5, 1.5)$) {\textcolor{popdark}{\textbf{Non fiable — Ne pas relayer — Vérifier source officielle (MINCOMMERCE)}}};
\node[label] at ($(cadre.west) + (1.0, 0.7)$) {Nom / Signature :};
\node[field, anchor=west, text width=0.75\linewidth] at ($(cadre.west) + (3.5, 0.7)$) {Élève \_\_\_\_\_\_\_\_\_\_\_\_\_\_\_\_\_\_\_\_\_\_\_\_\_\_\_\_\_\_\_\_\_\_\_\_\_\_\_\_};
\end{tikzpicture}
\legende{Exemple de fiche d'évaluation critique remplie (modèle type BEPC/Concours).}
\end{popfigure}
\end{center}

\begin{experience}[TP Guidéo : « Rédige ton évaluation » (Durée : 20 min)]
\textbf{Objectif :} Produire une évaluation critique écrite complète (180 mots) sur un support réel.
\textbf{Matériel :} Smartphone/PC connecté + fiche vierge (modèle ci-dessus) + stylo.
\textbf{Consignes :}
\begin{enumerate}
    \item Ouvrez la page Facebook \guillemotleft \texttt{Communes du Cameroun} \guillemotright~ ou \guillemotleft \texttt{Journal Officiel du Cameroun} \guillemotright.
    \item Choisissez \textbf{une} publication d'actualité (communiqué, arrêté, alerte) datant de \textbf{moins de 7 jours}.
    \item Remplissez la fiche d'évaluation (critères C.R.A.F.O.A + note + verdict).
    \item Rédigez le \textbf{paragraphe final structuré} (Intro / Développement connecté / Conclusion) sur la face verso.
    \item \textbf{Auto-évaluation :} Échangez votre copie avec un binôme. Il coche la grille de relecture (5 points, §4.4.2). Notez-vous /5.
\end{enumerate}
\textbf{Observations attendues :} Les élèves confondent souvent \guillemotleft Accessibilité \guillemotright (lisibilité) et \guillemotleft Accessibilité Web \guillemotright (libre accès). Le prof précise : \emph{Accessibilité = compréhensibilité ; Accessibilité Web = gratuité/absence de paywall/technique}.
\textbf{Interprétation :} Une note globale $< 5/10$ = info à rejeter. $5-7$ = info à croiser. $> 7$ = info exploitable.
\end{experience}

\begin{exercice}[Entraînement BEPC — Sujet type 2025]
Vous recevez par WhatsApp un message vocal (2 min 13 s) d'une voix d'homme affirmant : \guillemotleft Attention parents ! Le MINEDUB vient d'annoncer que les épreuves écrites du BEPC 2025 sont avancées au 15 mai au lieu du 20 juin. Préparez vos enfants ! \guillemotright
\textbf{Consigne :} En 150 mots maximum, rédigez l'évaluation critique de cette information en suivant le canevas \guillemotleft Introduction — Analyse C.R.A.F.O.A — Conclusion \guillemotright. Notez-la sur 10.
\end{exercice}

\begin{corrige}[Exercice BEPC — Message vocal BEPC 2025]
\textbf{Introduction} : Message vocal WhatsApp anonyme (reçu 10/03/2025) annonçant l'avancement des épreuves écrites BEPC 2025 au 15 mai. Évaluation de sa fiabilité pour une famille d'élève.

\textbf{Développement} :
\textbf{Crédibilité} : Source anonyme (WhatsApp, vocal, pas de nom), canal non officiel. \textbf{Donc} crédibilité nulle.
\textbf{Réactualisation} : Daté du jour (métadonnées WhatsApp), mais contenu invérifiable. \textbf{Toutefois}, la date annoncée (15 mai) contredit le calendrier officiel publié en janvier.
\textbf{Accessibilité} : Format audio, langue parlée accessible. \textbf{Donc} accessible.
\textbf{Fiabilité} : Aucune référence au \texttt{www.minedub.cm}, au Journal Officiel, ni à la circulaire n°XXX. \textbf{Par conséquent}, information non sourcée, invérifiable, contradictoire avec la source primaire.
\textbf{Objectivité} : Ton alarmiste (\guillemotleft Attention parents ! \guillemotright), visée manipulation/panique. \textbf{Donc} non objectif.
\textbf{Accessibilité Web} : Viral, gratuit, facile à transférer. \textbf{Donc} forte diffusion, danger élevé.

\textbf{Conclusion} : Note \textbf{1,5/10}. \textbf{Fake news avérée}. \textbf{Action :} Ne pas partager, signaler (WhatsApp \guillemotleft Signaler \guillemotright), consulter \texttt{minedub.cm} ou l'établissement pour le calendrier officiel.
\end{corrige}

\coursec{Application à des situations camerounaises concrètes}

Après avoir acquis la méthode rigoureuse de rédaction et de justification d'une évaluation critique (section précédente), nous passons maintenant à la mise en situation réelle. C'est ici que la théorie rencontre le terrain : ton téléphone vibre, ton fil d'actualité défile, la radio grésille en fond sonore. Comment réagir en citoyen numérique responsable ? Nous allons disséquer quatre cas typiquement camerounais, puis croiser les regards de trois médias sur un même événement. L'objectif est double : ancrer les six critères \cle{C.R.A.F.O.A.} (\cle{Crédibilité}, \cle{Pertinence}, \cle{Actualité}, \cle{Fiabilité}, \cle{Objectivité}, \cle{Accessibilité}) dans ta mémoire et te doter d'un réflexe de vérification avant tout partage.

\courssub{Rappel express : la grille C.R.A.F.O.A. en action}

Avant d'entrer dans les cas, fixons la grille de lecture. Chaque critère se note sur 5 (1 = très faible, 5 = excellent). Une information \textbf{exploitable} doit idéalement dépasser 3/5 sur \emph{tous} les critères ; une note $\le 2$ sur un critère critique (Fiabilité, Crédibilité, Actualité) doit déclencher une alerte rouge.

\begin{definition}[Grille C.R.A.F.O.A. pour l'évaluation rapide]
\begin{itemize}
    \item \textbf{C -- Crédibilité (Source) :} Qui émet l'information ? Institution connue, journaliste identifié, anonyme, compte parodique ?
    \item \textbf{R -- Pertinence (Rôle/Utilité) :} Cette information répond-elle à mon besoin ? Est-elle adaptée au contexte (camerounais, local) ?
    \item \textbf{A -- Actualité (Fraîcheur) :} Date de publication, date de l'événement. Une information de 2021 partagée en 2025 est une \emph{intox} potentielle.
    \item \textbf{F -- Fiabilité (Preuves) :} Sources citées, données chiffrées vérifiables, liens vers documents officiels, témoins directs.
    \item \textbf{O -- Objectivité (Neutralité) :} Ton factuel ou émotionnel ? Distinction claire entre fait et avis. Absence de conflit d'intérêt évident.
    \item \textbf{A -- Accessibilité (Forme) :} Langue comprise, format lisible, pas de barrière de paiement bloquante, site fonctionnel (pas d'erreur \texttt{404}).
\end{itemize}
\end{definition}

\courssub{Cas 1 : la rumeur WhatsApp -- « Fermeture immédiate de tous les établissements »}

\textbf{Situation :} mardi 14 janvier, 06h42. Tu reçois un message vocal WhatsApp (1 min 12 s) dans le groupe « Parents Élèves Lycée Bilingue Essos ». Une voix masculine, essoufflée, affirme : \og{} Le MINESEC vient de signer l'arrêté de fermeture de tous les lycées et collèges du Cameroun à cause de la grève des enseignants. C'est effectif dès aujourd'hui. N'envoyez pas vos enfants. Transmettez à tous vos groupes. \fg{} Le message porte la mention \texttt{Transféré plusieurs fois} (double flèche).

\begin{exemplebox}[Analyse C.R.A.F.O.A. du Cas 1 -- Rumeur WhatsApp]
\textbf{Évaluation critère par critère :}
\begin{itemize}
    \item \textbf{Crédibilité (1/5) :} source anonyme (voix non identifiée), pas de nom, pas de fonction. La mention \texttt{Transféré plusieurs fois} signale une chaîne de transmission non vérifiée.
    \item \textbf{Pertinence (5/5) :} extrêmement pertinent pour toi (parent/élève), impact direct sur l'organisation familiale.
    \item \textbf{Actualité (1/5) :} aucune date dans le message. « Dès aujourd'hui » est vague. Pas de référence à un numéro d'arrêté, pas de date de signature.
    \item \textbf{Fiabilité (1/5) :} aucune preuve. Pas de lien vers le site \texttt{minesec.gov.cm}, pas de capture d'écran de l'arrêté, pas de citation d'un communiqué de presse officiel.
    \item \textbf{Objectivité (1/5) :} ton alarmiste, impératif (\og{} N'envoyez pas \fg{}), injonction au partage viral (\og{} Transmettez à tous \fg{}). Langage de rumeur, pas d'information.
    \item \textbf{Accessibilité (4/5) :} message vocal en français courant, facile à écouter, mais format audio non indexable, non copiable-collable pour vérification.
\end{itemize}
\textbf{Note globale : 13/30.}

\textbf{Décision :} \textbf{NE PAS PARTAGER. SIGNALER comme fausse information dans le groupe.} \emph{Action correcte :} chercher le communiqué officiel sur \texttt{minesec.gov.cm} ou la page Facebook vérifiée du MINESEC, ou appeler l'administration de l'établissement.
\tcblower
\textbf{Justification rédigée (modèle BEPC) :}
\og{} Ce message vocal obtient une note de 13/30. Les critères critiques (Crédibilité, Actualité, Fiabilité, Objectivité) sont à 1/5. L'absence de source identifiable, de date précise, de référence administrative et le ton alarmiste en font une rumeur non fondée. Le partager contribuerait à la désinformation et à la panique. La démarche responsable est de le signaler et de consulter la source officielle (MINESEC). \fg{}
\end{exemplebox}

\begin{attention}[Piège du « Transféré plusieurs fois »]
Sur WhatsApp, la double flèche indique que le message n'a pas été composé par ton contact direct, mais qu'il a fait au moins deux sauts entre utilisateurs. C'est un indicateur technique fort de \emph{viralité non contrôlée}. Traite tout message \texttt{Transféré plusieurs fois} contenant une annonce officielle comme \textbf{suspect par défaut} tant qu'il n'est pas confirmé par le site \texttt{.gov.cm} correspondant.
\end{attention}

\courssub{Cas 2 : publication officielle MINSANTE -- campagne de vaccination (site .gov.cm)}

\textbf{Situation :} tu navigues sur \texttt{www.minsante.gov.cm} (onglet \og{} Actualités \fg{}). Tu tombes sur un article daté du \textbf{10 octobre 2024} : \og{} Lancement de la campagne de vaccination contre le HPV pour les filles de 9 à 14 ans \fg{}. L'article cite un communiqué de presse numéroté et daté du 08/10/2024, mentionne les districts de santé ciblés (avec liste PDF téléchargeable), les numéros verts et les partenaires techniques (OMS, Gavi).

\begin{propriete}[Analyse C.R.A.F.O.A. du Cas 2 -- Source institutionnelle .gov.cm]
\begin{itemize}
    \item \textbf{Crédibilité (5/5) :} nom de domaine officiel \texttt{.gov.cm}, auteur institutionnel (cellule de communication du MINSANTE), références administratives exactes (numéro et date du communiqué).
    \item \textbf{Pertinence (5/5) :} information de santé publique majeure, ciblée (filles de 9 à 14 ans), actionnable (lieux, numéros).
    \item \textbf{Actualité (5/5) :} date de publication (10/10/2024) cohérente avec l'événement (lancement). Communiqué récent (08/10/2024).
    \item \textbf{Fiabilité (5/5) :} preuves documentées : lien vers le PDF de la liste des districts, numéros verts vérifiables, partenaires internationaux nommés.
    \item \textbf{Objectivité (5/5) :} ton institutionnel, factuel, aucun adjectif qualificatif subjectif. Distinction claire entre le fait (\og{} La campagne vise à... \fg{}) et la position officielle (\og{} Nous encourageons... \fg{}).
    \item \textbf{Accessibilité (4/5) :} site accessible, PDF téléchargeable. \emph{Bémol :} site parfois lent aux heures de pointe ; pas de version en langues nationales pour cette page.
\end{itemize}
\textbf{Note globale : 29/30.}

\textbf{Décision :} \textbf{INFORMATION VALIDÉE. PARTAGE RECOMMANDÉ (avec le lien source).}
\end{propriete}

\begin{attention}[Vigilance sur les sites .gov.cm : actualité et faux sites]
Un nom de domaine \texttt{.gov.cm} garantit l'\textbf{authenticité de l'émetteur} (Crédibilité 5/5), mais \textbf{pas l'actualité du contenu}. Une page peut ne pas avoir été mise à jour depuis deux ans (exemple : un calendrier de concours 2022 encore affiché en 2025). \textbf{Vérifie toujours :} (1) la date de \emph{dernière mise à jour} de la page (souvent en bas), (2) la date du document PDF joint, (3) s'il existe une version plus récente via le moteur de recherche du site. Méfie-toi aussi des faux sites imitant l'adresse officielle : \texttt{minsante-gov-cm.com} (faux) contre \texttt{minsante.gov.cm} (vrai).
\end{attention}

\courssub{Cas 3 : publicité « produit miracle » sur Facebook -- « Tisane qui guérit tout »}

\textbf{Situation :} sur ton fil Facebook, une publication marquée \og{} Sponsorisée \fg{} (petit label sous le nom de la page \texttt{BienÊtreNaturel237}). Vidéo : une femme en pagne danse en tenant une bouteille. Texte : \og{} \textbf{Guérit le diabète, l'hypertension, l'ulcère et l'infertilité en 7 jours !} 100\,\% naturel, zéro effet secondaire. Commandez au 6XX XX XX XX. Livraison partout au Cameroun. \fg{} Commentaires : \og{} Merci docteur, mon mari a guéri ! \fg{}, \og{} Prix svp ? \fg{}, \og{} Arnaque ! \fg{} (réponse de la page : \og{} Ce sont des jaloux, commande et tu verras \fg{}).

\begin{exoresolu}[Analyse C.R.A.F.O.A. du Cas 3 -- Publicité santé non vérifiée]
\textbf{Énoncé :} évalue cette publication et rédige la justification de ton refus d'achat et de partage.
\tcblower
\textbf{Solution détaillée :}
\begin{itemize}
    \item \textbf{Crédibilité (1/5) :} page Facebook générique, pas de raison sociale, pas d'agrément du Ministère de la Santé, pas d'autorisation de mise sur le marché. Le label \texttt{Sponsorisé} signifie publicité payante, pas caution de Facebook.
    \item \textbf{Pertinence (3/5) :} problèmes de santé réels (diabète, hypertension), forte demande. Mais la réponse proposée est inadaptée (un produit unique pour des maladies distinctes).
    \item \textbf{Actualité (2/5) :} pas de date sur la publication. Vidéo potentiellement ancienne, recyclée.
    \item \textbf{Fiabilité (1/5) :} \textbf{drapeau rouge majeur :} une revendication thérapeutique multiple (\og{} guérit le diabète ET l'hypertension ET... \fg{}) est scientifiquement impossible pour une simple tisane. Aucune étude citée, aucun avis médical. Les témoignages en commentaires sont anecdotiques, non vérifiables, potentiellement achetés.
    \item \textbf{Objectivité (1/5) :} langage marketing agressif (\og{} 100\,\% naturel, zéro effet secondaire \fg{} = affirmation mensongère), gestion agressive des commentaires négatifs.
    \item \textbf{Accessibilité (5/5) :} vidéo courte, français simple, numéro WhatsApp direct, paiement Mobile Money (MTN MoMo / Orange Money) intégré.
\end{itemize}
\textbf{Note globale : 13/30.}

\textbf{Justification type BEPC :} \og{} Cette publication obtient 13/30. Elle échoue gravement sur la Fiabilité (1/5) : une tisane ne peut guérir des maladies chroniques distinctes. L'absence d'autorisation sanitaire, de données médicales et le ton marketing disqualifient l'information. C'est une publicité trompeuse, dangereuse pour la santé publique (retard de prise en charge médicale). Je ne partage pas, je signale comme arnaque / produit dangereux. \fg{}
\end{exoresolu}

\courssub{Cas 4 : article de presse -- « Budget 2025 : l'éducation reçoit 3 200 milliards FCFA »}

\textbf{Situation :} article en ligne sur le site d'un journal camerounais de référence, daté du \textbf{15 novembre 2024}, signé par un journaliste identifié et accrédité à l'Assemblée Nationale. Titre : \og{} Loi de Finances 2025 : l'éducation de base et secondaire capte 16,8\,\% du budget national \fg{}. Corps de l'article : chiffres exacts, comparaison avec 2024, citation nominale du ministre des Finances, citation d'un député de l'opposition, lien vers le texte intégral de la Loi de Finances sur un site \texttt{.gov.cm}. Graphique de répartition par sous-secteur.

\begin{propriete}[Analyse C.R.A.F.O.A. du Cas 4 -- Presse écrite de référence]
\begin{itemize}
    \item \textbf{Crédibilité (5/5) :} journal de référence (historique, ligne éditoriale connue), journaliste identifié et accrédité.
    \item \textbf{Pertinence (5/5) :} sujet d'intérêt national, impact direct sur les établissements scolaires (infrastructures, manuels, cantines).
    \item \textbf{Actualité (5/5) :} date précise (15/11/2024), événement récent (vote du budget), perspective année 2025.
    \item \textbf{Fiabilité (5/5) :} données chiffrées sourcées (Loi de Finances), lien direct vers la source primaire, citations croisées (majorité / opposition), graphique sourcé.
    \item \textbf{Objectivité (4/5) :} traitement équilibré : chiffre officiel + réaction de l'opposition + analyse d'expert. \emph{Réserve :} titre légèrement orienté (\og{} capte \fg{}, verbe fort) alors que le corps de l'article reste neutre.
    \item \textbf{Accessibilité (4/5) :} site web adapté au mobile, article gratuit (parfois limité après quelques articles par mois), graphique lisible. Langue française standard.
\end{itemize}
\textbf{Note globale : 28/30.}

\textbf{Décision :} \textbf{INFORMATION DE HAUTE QUALITÉ. PARTAGE AVEC CITATION DE LA SOURCE (nom du journal, date).}
\end{propriete}

\courssub{Analyse comparative : un même événement, trois regards médiatiques}

\textbf{Scénario :} annonce de la hausse du prix du carburant à la pompe (Super : 840 FCFA/L $\to$ 910 FCFA/L ; Gasoil : 720 FCFA/L $\to$ 780 FCFA/L) le 1\ier{} février 2025. Trois médias traitent l'information le même jour à 12h.

\begin{popfigure}
\begin{tikzpicture}[
    screen/.style={rectangle, rounded corners=4pt, minimum width=5.5cm, minimum height=4.5cm, draw=popdark, thick, fill=white},
    header/.style={fill=popblue, text=white, font=\bfseries\small, rounded corners=3pt, minimum height=0.7cm, align=center},
    content/.style={font=\tiny, align=left, text=popdark},
    tag/.style={font=\tiny\bfseries, text=white, fill=poporange, rounded corners=2pt, inner sep=1.5pt}
]
% --- ECRAN 1 : RADIO ---
\node[screen] (radio) at (0,0) {};
\node[header, minimum width=5.5cm] at (radio.north) {RADIO CRTV -- Journal de 12h};
\node[content, anchor=north west, text width=5.1cm] at ([yshift=-0.8cm]radio.north west) {
\textbullet\ \textbf{Ton :} neutre, officiel.\par
\textbullet\ \textbf{Source :} communiqué du MINCOMMERCE lu par le speaker.\par
\textbullet\ \textbf{Contenu :} nouveaux tarifs, date d'application, déclaration du Ministre.\par
\textbullet\ \textbf{Durée :} 45 secondes.\par
\textbullet\ \textbf{Interaction :} aucune (flux linéaire).\par
\textbullet\ \textbf{Accessibilité :} haute (autoradio, téléphone, zone rurale).
};
\node[tag, anchor=south east] at ([xshift=-0.2cm, yshift=0.2cm]radio.south east) {C.R.A.F.O.A. $\approx$ 26/30};

% --- ECRAN 2 : TV ---
\node[screen] (tv) at (7.5,0) {};
\node[header, minimum width=5.5cm, fill=poppurple] at (tv.north) {TÉLÉVISION -- Journal de 13h};
\node[content, anchor=north west, text width=5.1cm] at ([yshift=-0.8cm]tv.north west) {
\textbullet\ \textbf{Ton :} visuel, mis en scène.\par
\textbullet\ \textbf{Source :} reportage à la pompe + interview du Ministre + micro-trottoir.\par
\textbullet\ \textbf{Contenu :} images de files d'attente, calcul (10 L = 9 100 FCFA), réaction d'un usager.\par
\textbullet\ \textbf{Durée :} 3 min 30 s.\par
\textbullet\ \textbf{Interaction :} rediffusion en ligne, commentaires.\par
\textbullet\ \textbf{Accessibilité :} moyenne (électricité, données pour le replay).
};
\node[tag, anchor=south east] at ([xshift=-0.2cm, yshift=0.2cm]tv.south east) {C.R.A.F.O.A. $\approx$ 27/30};

% --- ECRAN 3 : RESEAUX SOCIAUX ---
\node[screen] (fb) at (15,0) {};
\node[header, minimum width=5.5cm, fill=poppink] at (fb.north) {FACEBOOK / WHATSAPP -- Groupes};
\node[content, anchor=north west, text width=5.1cm] at ([yshift=-0.8cm]fb.north west) {
\textbullet\ \textbf{Ton :} émotionnel, viral, partisan.\par
\textbullet\ \textbf{Sources :} capture d'écran floue, mème, audio WhatsApp anonyme, blog obscur.\par
\textbullet\ \textbf{Contenu :} \og{} Encore une arnaque ! \fg{}, \og{} Grève générale lundi ! \fg{}, \og{} Le Super à 1000 F dans 3 jours \fg{} (rumeur).\par
\textbullet\ \textbf{Durée :} éphémère (défilement infini).\par
\textbullet\ \textbf{Interaction :} partages massifs, commentaires houleux.\par
\textbullet\ \textbf{Accessibilité :} maximale (smartphone, données à bas coût).
};
\node[tag, anchor=south east] at ([xshift=-0.2cm, yshift=0.2cm]fb.south east) {C.R.A.F.O.A. $\approx$ 11/30};

% Flèches comparatives
\draw[popblue, thick, -{Stealth[length=3mm]}, bend left=15] (radio.east) to node[midway, above, font=\tiny, text=popblue] {Source primaire officielle} (tv.west);
\draw[poporange, thick, -{Stealth[length=3mm]}, bend right=15] (tv.east) to node[midway, below, font=\tiny, text=poporange] {Reprise + émotion + bruit} (fb.west);
\draw[popgreen, dashed, thick, -{Stealth[length=3mm]}] (radio.north east) to[out=45, in=135] node[midway, above, font=\tiny, text=popgreen] {Vérification croisée recommandée} (fb.north west);
\end{tikzpicture}
\legende{Comparatif visuel : traitement de la hausse du carburant (1\ier{} février 2025) par la radio (officiel), la télévision (contextualisé) et les réseaux sociaux (émotionnel/viral). Notes C.R.A.F.O.A. estimatives.}
\end{popfigure}

\begin{experience}[TP guidé : analyse comparative croisée -- Durée 20 min]
\textbf{Objectif :} appliquer la grille C.R.A.F.O.A. à trois traitements d'un même fait d'actualité camerounaise récente.

\textbf{Matériel :} smartphone ou PC avec connexion. Accès à : (1) le site d'un média public ou d'un journal de référence, (2) la page Facebook d'un grand média, (3) un groupe WhatsApp/Facebook d'actualité (ou une simulation fournie par le professeur).

\textbf{Consignes :}
\begin{enumerate}
    \item Choisis un événement récent (rentrée scolaire, compétition sportive, budget, grève, catastrophe naturelle).
    \item Ouvre les trois sources simultanément (onglets du navigateur).
    \item Remplis le tableau comparatif ci-dessous pour chaque critère C.R.A.F.O.A. (note /5 + preuve ou observation).
    \item Rédige une synthèse de cinq lignes : \og{} Quelle source je cite pour mon exposé ? Pourquoi ? \fg{}
\end{enumerate}

\textbf{Tableau à reproduire dans ton cahier :}
\begin{center}
\begin{tabularx}{\linewidth}{|l|X|X|X|}
\hline
\rowcolor{popblueL}
\textbf{Critère} & \textbf{Radio / Site officiel} & \textbf{TV / Site de presse} & \textbf{Réseaux sociaux} \\
\hline
\textbf{C -- Crédibilité} & Note : /5 \newline Preuve : & Note : /5 \newline Preuve : & Note : /5 \newline Preuve : \\
\hline
\textbf{R -- Pertinence} & Note : /5 \newline Preuve : & Note : /5 \newline Preuve : & Note : /5 \newline Preuve : \\
\hline
\textbf{A -- Actualité} & Note : /5 \newline Preuve : & Note : /5 \newline Preuve : & Note : /5 \newline Preuve : \\
\hline
\textbf{F -- Fiabilité} & Note : /5 \newline Preuve : & Note : /5 \newline Preuve : & Note : /5 \newline Preuve : \\
\hline
\textbf{O -- Objectivité} & Note : /5 \newline Preuve : & Note : /5 \newline Preuve : & Note : /5 \newline Preuve : \\
\hline
\textbf{A -- Accessibilité} & Note : /5 \newline Preuve : & Note : /5 \newline Preuve : & Note : /5 \newline Preuve : \\
\hline
\textbf{TOTAL /30} & \textbf{ } & \textbf{ } & \textbf{ } \\
\hline
\end{tabularx}
\end{center}

\textbf{Observations attendues :} la source officielle gagne sur C, F, A. La presse gagne sur O, R (analyse). Les réseaux sociaux gagnent sur la vitesse et le format, mais perdent sur C, F, O.

\textbf{Interprétation :} l'information \textbf{de référence} (citable) est celle qui valide C, F, A $\ge 4$. L'information \textbf{de compréhension} (contexte, débat) vient de la presse. L'information \textbf{de veille} (réactions, rumeurs) vient des réseaux sociaux -- à traiter avec la plus grande prudence.
\end{experience}

\courssub{Synthèse visuelle : carte mentale -- cas pratiques et critères C.R.A.F.O.A.}

La figure ci-dessous relie chaque cas d'étude aux critères qui ont \textbf{décidé} du verdict (partager / signaler / ignorer). Les nœuds \textbf{verts} indiquent un critère fort ($\ge 4$), les nœuds \textbf{roses} un critère faible ($\le 2$) ayant déclenché l'alerte.

\begin{popfigure}
\begin{tikzpicture}[
    mindmap,
    concept color=popblue,
    level 1 concept/.append style={level distance=4.5cm, sibling angle=90, font=\small\bfseries},
    level 2 concept/.append style={level distance=3.2cm, sibling angle=45, font=\tiny},
    every node/.style={scale=0.85},
    strong/.style={concept color=popgreen!70!popblue, fill=popgreen!20},
    weak/.style={concept color=poppink!70!popblue, fill=poppink!20},
    case/.style={concept color=poporange!70!popblue, fill=poporange!20, font=\small\bfseries},
    crit/.style={font=\tiny\bfseries}
]
\node[case] {Cas camerounais}
    child[grow=30] {
        node[case, align=center] (c1){Cas 1\\Rumeur WA}
        child { node[crit, weak] (c1f) {Fiabilité 1/5} }
        child { node[crit, weak] (c1c) {Crédibilité 1/5} }
        child { node[crit, weak] (c1a) {Actualité 1/5} }
        child { node[crit, weak] (c1o) {Objectivité 1/5} }
        child { node[crit, strong] (c1r) {Pertinence 5/5} }
        child { node[crit] (c1ac) {Access. 4/5} }
    }
    child[grow=150] {
        node[case, align=center] (c2){Cas 2\\MINSANTE .gov}
        child { node[crit, strong] (c2c) {Crédibilité 5/5} }
        child { node[crit, strong] (c2f) {Fiabilité 5/5} }
        child { node[crit, strong] (c2a) {Actualité 5/5} }
        child { node[crit, strong] (c2o) {Objectivité 5/5} }
        child { node[crit, strong] (c2r) {Pertinence 5/5} }
        child { node[crit] (c2ac) {Access. 4/5} }
    }
    child[grow=210] {
        node[case, align=center] (c3){Cas 3\\Pub miracle FB}
        child { node[crit, weak] (c3f) {Fiabilité 1/5} }
        child { node[crit, weak] (c3c) {Crédibilité 1/5} }
        child { node[crit, weak] (c3o) {Objectivité 1/5} }
        child { node[crit] (c3a) {Actualité 2/5} }
        child { node[crit] (c3r) {Pertinence 3/5} }
        child { node[crit, strong] (c3ac) {Access. 5/5} }
    }
    child[grow=330] {
        node[case, align=center] (c4){Cas 4\\Presse budget}
        child { node[crit, strong] (c4c) {Crédibilité 5/5} }
        child { node[crit, strong] (c4f) {Fiabilité 5/5} }
        child { node[crit, strong] (c4a) {Actualité 5/5} }
        child { node[crit] (c4o) {Objectivité 4/5} }
        child { node[crit, strong] (c4r) {Pertinence 5/5} }
        child { node[crit] (c4ac) {Access. 4/5} }
    };

% Légende manuelle
\node[align=left, font=\tiny, anchor=north west, text=popdark] at (8, -3.5) {
\textbf{Légende :} \textcolor{popgreen!70!black}{\rule{2.5mm}{2.5mm}} Critère fort ($\ge 4$) -- \textcolor{poppink!70!black}{\rule{2.5mm}{2.5mm}} Critère faible ($\le 2$ = alerte) -- \rule{2.5mm}{2.5mm} Critère moyen.
\par\textbf{Enseignement :} un seul critère \textcolor{poppink!70!black}{\textbf{FAIBLE}} sur C, F, A ou O suffit pour \textbf{REJETER} l'information (cas 1 et 3).
\par\textbf{Règle d'or :} la \cle{Fiabilité} et la \cle{Crédibilité} sont des \textbf{prérequis non négociables}.
};
\end{tikzpicture}
\legende{Carte mentale reliant les quatre cas camerounais aux six critères C.R.A.F.O.A. Vert = point fort (validé), rose = point de rupture (rejet).}
\end{popfigure}

\courssub{Méthode : construire ta « check-list express » sur ton téléphone}

Tu n'auras pas toujours ta grille complète sous les yeux. Voici comment créer un raccourci mental (et numérique) pour évaluer une information en \textbf{30 secondes chrono} avant de cliquer sur \og{} Partager \fg{}.

\begin{methode}[Ma check-list express « 3 questions / 3 clics » -- à enregistrer dans \texttt{Notes} ou \texttt{Google Keep}]
\textbf{Étape 1 -- LA SOURCE (clic 1 : profil / adresse du site) -- \emph{critères C et F}}
\begin{enumerate}
    \item \textbf{Qui parle ?} Nom affiché $\to$ clic sur le profil $\to$ vérifie : site \texttt{.gov.cm}, journal connu, compte vérifié (badge), ou anonyme ?
    \item \textbf{Preuve ?} Y a-t-il un lien cliquable vers un document officiel (PDF, décret, communiqué) ? \textbf{Pas de lien = suspect.}
\end{enumerate}

\textbf{Étape 2 -- LA DATE (clic 2 : informations de publication) -- \emph{critère A}}
\begin{enumerate}
    \item \textbf{Quand ?} Date de l'article, du post ou de la vidéo. \textbf{Si plus de 6 mois sans mention « mise à jour » = attention} (le contexte a pu changer).
    \item \textbf{Événement ?} La date de l'événement correspond-elle à la date de publication ? (Exemple : une photo de 2019 partagée comme \og{} aujourd'hui \fg{}.)
\end{enumerate}

\textbf{Étape 3 -- LE TON ET L'INTENTION (lecture diagonale de 10 secondes) -- \emph{critères O et R}}
\begin{enumerate}
    \item \textbf{Fait ou opinion ?} Verbes : \og{} le Ministre a signé \fg{} (fait) contre \og{} c'est une honte \fg{} (opinion).
    \item \textbf{Injonction ?} \og{} Partagez urgemment \fg{}, \og{} Achetez maintenant \fg{}, \og{} Signez la pétition \fg{} $\to$ \textbf{manipulation probable.}
\end{enumerate}

\textbf{Décision immédiate :}
\begin{itemize}
    \item \textbf{3 réponses positives} (source fiable, date correcte, ton factuel) $\to$ \textbf{PARTAGE} (avec le lien source).
    \item \textbf{1 réponse négative} sur la source, la date ou la preuve $\to$ \textbf{NE PARTAGE PAS. VÉRIFIE} (moteur de recherche, site officiel).
    \item \textbf{2 réponses négatives ou injonction virale} $\to$ \textbf{SIGNALE} (fausse information, arnaque, danger).
\end{itemize}
\end{methode}

\begin{savaistu}[Le BEPC et l'esprit critique]
Au BEPC, plusieurs épreuves (français, histoire-géographie, éducation à la citoyenneté) te demandent de \textbf{lire un document, de dégager son sens et de donner ton avis argumenté}. La grille C.R.A.F.O.A. est exactement la bonne démarche : identifier la source (C), dater le document (A), repérer les preuves (F), distinguer fait et opinion (O). Au-delà de l'examen, le Cameroun compte des millions d'utilisateurs de WhatsApp et de Facebook : les rumeurs sur la santé, les concours ou la sécurité s'y propagent en quelques minutes. En refusant de relayer une information douteuse, tu protèges concrètement ta famille et ta communauté : c'est cela, la citoyenneté numérique.
\end{savaistu}

\courssub{Exercices d'application -- situations d'examen (type BEPC)}

\begin{exercice}[Cas pratique : l'annonce de recrutement « Fonction Publique 2025 »]
Tu reçois ce message sur un groupe WhatsApp \og{} Emploi Jeunes 237 \fg{} (5 000 membres) :
\og{} \textbf{RECRUTEMENT MINFOPRA 2025 :} 5 000 postes (enseignants, santé, administration). Inscription GRATUITE via le lien : \texttt{www.minfopra-recrutement-2025-cm.com}. Date limite : 30/06/2025. Partagez pour vos frères ! \fg{}
Le lien renvoie vers un formulaire en ligne demandant : nom, prénom, numéro de CNI, diplôme, et \textbf{numéro Mobile Money pour des « frais de dossier » de 2 000 FCFA}.
\begin{enumerate}
    \item Applique la grille C.R.A.F.O.A. (note sur 5 + justification courte par critère).
    \item Identifie \textbf{trois indices} prouvant l'arnaque (adresse du site, demande d'argent, formulaire).
    \item Rédige la réponse à poster dans le groupe pour alerter tes camarades (ton constructif, cinq lignes maximum).
\end{enumerate}
\end{exercice}

\begin{corrige}[Exercice : recrutement MINFOPRA]
\begin{enumerate}
    \item \textbf{C.R.A.F.O.A. :} C (1/5) -- domaine \texttt{.com} non officiel (le vrai site est en \texttt{.gov.cm}) ; F (1/5) -- formulaire anonyme + demande d'argent (un concours public officiel ne demande jamais de paiement par Mobile Money) ; A (2/5) -- date limite lointaine mais aucun communiqué officiel récent ; O (1/5) -- injonction au partage + appât du gain ; R (4/5) -- besoin réel d'emploi ; Accessibilité (4/5) -- formulaire facile d'accès. \textbf{Total : 13/30.}
    \item \textbf{Trois indices :} (1) l'adresse \texttt{minfopra-recrutement-2025-cm.com} imite le site officiel sans en être un (faux domaine). (2) La demande de paiement de 2 000 FCFA par Mobile Money pour des « frais de dossier ». (3) Un simple formulaire en ligne gratuit et anonyme au lieu d'une plateforme officielle de concours.
    \item \textbf{Réponse type :} \og{} \textbf{ALERTE ARNAQUE} : ce lien ne vient PAS du MINFOPRA. Le vrai site du ministère se termine par \texttt{.gov.cm}. Un concours de la fonction publique ne demande jamais de paiement par Mobile Money pour s'inscrire. Ne donnez ni votre numéro de CNI ni votre argent. Signalez ce message. \fg{}
\end{enumerate}
\end{corrige}

\begin{exercice}[Analyse comparative -- grève des transporteurs (Yaoundé/Douala)]
Trois médias traitent la grève du 10 mars 2025 :
\begin{itemize}
    \item \textbf{Doc A (radio publique, 7h) :} \og{} Le groupement des transporteurs a suspendu son mot d'ordre après la rencontre avec le ministre des Transports. La circulation reprend progressivement. \fg{} (30 secondes).
    \item \textbf{Doc B (page Facebook « Actu Cameroun », 9h) :} vidéo tremblotée montrant un taxi brûlé. Légende : \og{} Guerre à Yaoundé ! Grève générale illimitée ! On nous a menti ! Partagez ! \fg{} (15 000 partages).
    \item \textbf{Doc C (journal en ligne, 11h) :} article titré \og{} Transport : accord signé sous réserve. Les syndicats minoritaires maintiennent la grève. Usagers piégés. \fg{}, avec citations de syndicalistes des deux camps, du ministre des Transports, photos de rues calmes et photos d'incidents.
\end{itemize}
\begin{enumerate}
    \item Complète le tableau C.R.A.F.O.A. comparatif (notes /5) pour les trois documents.
    \item Quelle source cites-tu pour ton exposé sur \og{} l'organisation des transports urbains \fg{} ? Justifie en trois arguments C.R.A.F.O.A.
    \item Pourquoi le Doc B est-il \textbf{plus viral} que les Docs A et C malgré sa faible note ? (Réponse en deux lignes : psychologie + algorithme.)
\end{enumerate}
\end{exercice}

\begin{corrige}[Exercice : grève des transporteurs]
\begin{enumerate}
    \item \textbf{Tableau (synthèse) :}
    \begin{center}
    \begin{tabularx}{\linewidth}{|X|c|c|c|}
    \hline
    \rowcolor{popblueL}
    \textbf{Critère} & \textbf{Doc A (Radio off.)} & \textbf{Doc B (FB viral)} & \textbf{Doc C (Presse)} \\
    \hline
    Crédibilité & 5 & 1 & 5 \\
    \hline
    Fiabilité & 4 & 1 & 5 \\
    \hline
    Actualité & 5 & 3 (vidéo non datée) & 5 \\
    \hline
    Objectivité & 5 & 1 & 4 \\
    \hline
    Pertinence & 4 & 5 (émotion) & 5 \\
    \hline
    Accessibilité & 5 & 5 & 4 \\
    \hline
    \textbf{Total} & \textbf{28/30} & \textbf{16/30} & \textbf{28/30} \\
    \hline
    \end{tabularx}
    \end{center}
    \item \textbf{Source citée : Doc C (presse) ou Doc A (radio).} Justification : (1) \textbf{Fiabilité 5/5} : sources croisées, citations nommées, accord documenté. (2) \textbf{Objectivité 4/5} : la parole est donnée aux deux camps syndicaux et au ministère, avec la nuance \og{} accord sous réserve \fg{}. (3) \textbf{Accessibilité et archivage} : article web daté, adresse stable, citable dans une bibliographie (contrairement au flux radio éphémère ou au post Facebook instable).
    \item \textbf{Viralité du Doc B :} (1) l'\textbf{émotion forte} (peur, colère, sentiment d'injustice) provoque un partage réflexe. (2) L'\textbf{algorithme} de Facebook favorise les contenus qui génèrent des réactions (commentaires, partages), ce qui amplifie leur portée. Une vidéo courte avec une légende choc est optimisée pour la viralité, pas pour la vérité.
\end{enumerate}
\end{corrige}

\begin{aretenir}[Bilan de la leçon -- Maîtrise des médias et de l'information]
\begin{itemize}
    \item \textbf{L'information de qualité} se reconnaît aux six piliers \cle{C.R.A.F.O.A.} (Crédibilité, Pertinence, Actualité, Fiabilité, Objectivité, Accessibilité).
    \item \textbf{Contexte camerounais :} les sites \texttt{.gov.cm} sont la référence absolue (C/F/A), mais vérifie toujours la date. La presse de référence apporte analyse et contexte (O/R). La radio et la télévision publiques sont officielles et très accessibles. Les réseaux sociaux sont rapides et émotionnels, mais présentent un \textbf{risque majeur} (faible C/F/O).
    \item \textbf{Réflexe d'examen :} face à un document $\to$ appliquer C.R.A.F.O.A. $\to$ croiser les sources $\to$ produire une conclusion argumentée et justifiée.
    \item \textbf{Outillage :} ta check-list « 3 questions » dans ton téléphone est ton bouclier quotidien contre l'intoxication.
    \item \textbf{Éthique :} ne pas partager une rumeur, c'est protéger sa communauté ; la signaler, c'est faire preuve de citoyenneté numérique active.
\end{itemize}
\end{aretenir}

\newpage

\coursec{Activité d'intégration}

\coursec{Acquisition d'une maîtrise des médias et de l'informatique : Critères d'une bonne information}

\courssub{Objectifs de l'activité}
Cette activité d'intégration vise à mobiliser l'ensemble des acquis de la leçon sur les critères d'une bonne information. À l'issue de ce travail, vous serez capable de :
\begin{itemize}
    \item Appliquer systématiquement la grille \textbf{C.R.A.F.O.A} (\textbf{C}rédibilité, \textbf{R}elévance, \textbf{A}ctualité, \textbf{F}iabilité, \textbf{O}bjectivité, \textbf{A}ccessibilité) à une information réelle circulant sur les réseaux sociaux.
    \item Rédiger un rapport d'évaluation critique structuré, argumenté et concis (environ 150 mots ± 10\%).
    \item Présenter et défendre oralement votre analyse devant la classe, en adoptant une posture citoyenne et responsable face à la désinformation.
    \item Identifier les mécanismes psychologiques et techniques (titre racoleur, absence de source, urgence émotionnelle) qui favorisent la propagation des \emph{fake news}.
\end{itemize}

\courssub{Situation déclenchante}
\begin{popfigure}
\begin{tikzpicture}[node distance=1.2cm, every node/.style={font=\small}]
\node[draw=poporange, thick, rounded corners, fill=poporange!10, text width=0.92\linewidth, align=left] (box) {
\textbf{Contexte :} Nous sommes en période de rentrée scolaire (septembre 2025) à Yaoundé. Depuis 48 heures, une publication virale circule massivement sur les groupes WhatsApp « Parents d'élèves Cameroun », « Bons plans Yaoundé » et sur des pages Facebook comme « Actu Cameroun 237 » (120k abonnés).\par\medskip
\textbf{Contenu du post (capture d'écran simulée) :}\par
\begin{tabular}{@{}p{0.9\linewidth}@{}}
\texttt{\textbf{[URGENT - PARTAGEZ MAX]}} \\
\texttt{Le MINEDUB vient d'annoncer une aide exceptionnelle de \textbf{50 000 FCFA} par élève pour la rentrée 2025-2026. Cette subvention vise à soutenir les parents face à la vie chère.}\\
\texttt{\textbf{Procédure :} Cliquez sur le lien ci-dessous, entrez le nom de l'école, la classe et votre numéro Mobile Money (MTN MoMo ou Orange Money). Le virement est instantané.}\\
\texttt{\textbf{Lien :} \url{http://aide-minedub-2025.cm.ng/verify.php}}\\
\texttt{« Faites passer à tous les parents, ça ne coûte rien d'essayer ! » \#Rentrée2025 \#Minedub \#AideGouvernement}
\end{tabular}
\par\medskip
\textbf{Indices visuels :} L'image jointe montre un logo flou du Ministère de l'Éducation de Base (MINEDUB), un fond aux couleurs du drapeau camerounais, et un bouton rouge clignotant « RECEVOIR MON ARGENT ». Le compte qui a posté s'appelle « \texttt{Info\_Officielle\_237} » (créé il y a 3 jours, 12 publications, 45 abonnés).
};
\end{tikzpicture}
\legende{Simulation d'une capture d'écran d'une publication virale suspecte sur Facebook/WhatsApp.}
\end{popfigure}

\courssub{Consignes de l'activité}
Travail en groupes de 4 élèves. Durée totale : 45 minutes (25 min analyse + 10 min rédaction + 10 min restitution).
\begin{enumerate}
    \item \textbf{Analyse individuelle (5 min) :} Chaque élève lit le post et note spontanément : « Je crois / Je ne crois pas » et \emph{Pourquoi ?} (3 arguments).
    \item \textbf{Mise en commun et grille C.R.A.F.O.A (20 min) :} En groupe, complétez le tableau d'évaluation ci-dessous pour \textbf{chaque critère}. Cochez \texttt{Validé}, \texttt{Partiel} ou \texttt{Non validé} et citez \textbf{un élément de preuve} issu du post ou de votre connaissance du contexte camerounais.
    \item \textbf{Rédaction du rapport (10 min) :} Rédigez collectivement un \textbf{rapport d'évaluation} (150 mots ± 10\%) structuré ainsi : \textbf{Verdict final} (Fiable / Douteux / Faux) -- \textbf{3 arguments majeurs} (basés sur les critères C.R.A.F.O.A les plus discriminants) -- \textbf{Recommandation} (Action à mener : signaler, supprimer, alerter la famille).
    \item \textbf{Présentation orale (10 min) :} Le porte-parole désigné présente le verdict et les 3 arguments en 1 minute 30 max. La classe et le professeur posent 2 questions de clarification.
\end{enumerate}

\courssub{Ressources théoriques et contextuelles}
\begin{definition}[Grille d'évaluation C.R.A.F.O.A -- Rappel synthétique]
\begin{center}
\begin{tabularx}{\linewidth}{|l|X|X|}\hline
\rowcolor{popblueL}
\textbf{Critère} & \textbf{Question clé} & \textbf{Indicateurs de validation} \\\hline
\textbf{C} -- Crédibilité & Qui émet l'info ? Quelle est son autorité ? & Source officielle identifiée (site \texttt{.gov.cm}, compte vérifié), auteur connu, expertise reconnue. \\\hline
\textbf{R} -- Pertinence (Relevance) & L'info répond-elle à mon besoin ? & Contexte adapté, public ciblé clair, utilité immédiate. \\\hline
\textbf{A} -- Actualité & L'info est-elle à jour ? & Date de publication visible, cohérence avec le calendrier réel (rentrée, budget). \\\hline
\textbf{F} -- Fiabilité & L'info est-elle vérifiable, corroborée ? & Sources croisées (site officiel, presse reconnue), données chiffrées précises, lien direct vers l'acte administratif. \\\hline
\textbf{O} -- Objectivité & Le ton est-il neutre, factuel ? & Absence de langage émotionnel, d'impératifs de partage, de majuscules excessives, de fautes d'orthographe. \\\hline
\textbf{A} -- Accessibilité & L'info est-elle accessible sans risque ? & URL officielle (HTTPS, domaine connu), pas de demande de données sensibles (code PIN, OTP) via lien externe. \\\hline
\end{tabularx}
\end{center}
\end{definition}

\begin{propriete}[Références contextuelles camerounaises utiles]
\begin{itemize}
    \item Site officiel MINEDUB : \texttt{www.minedub.gov.cm} (domaine \texttt{.gov.cm} ou \texttt{.cm}).
    \item Communiqués officiels : diffusés via \texttt{CRTV}, \texttt{Cameroon Tribune}, page Facebook vérifiée « \texttt{Ministry of Basic Education Cameroon} » (badge bleu).
    \item Procédures administratives : Aucune allocation financière directe aux parents via lien Mobile Money non authentifié (USSD \texttt{\#150\#} ou \texttt{\#144\#} pour les services officiels MTN/Orange).
    \item Arnaque classique : \emph{Phishing} (hameçonnage) pour voler identifiants MoMo/OM ou installer un malware.
\end{itemize}
\end{propriete}

\courssub{Démarche et corrigé commenté}
\begin{exoresolu}[Évaluation collaborative du post « Aide 50 000 FCFA »]
\textbf{Énoncé :} À l'aide de la grille C.R.A.F.O.A, évaluer la crédibilité de la publication virale présentée ci-dessus et rédiger un rapport d'évaluation structuré (Verdict, 3 arguments, Recommandation).
\tcblower
\textbf{Étape 1 -- Remplissage de la grille C.R.A.F.O.A par le groupe}

\begin{center}
\begin{tabularx}{\linewidth}{|l|c|X|}\hline
\rowcolor{popblueL}
\textbf{Critère} & \textbf{Statut} & \textbf{Preuve / Justification (Extrait du post ou connaissance externe)} \\\hline
\textbf{C -- Crédibilité} & \cellcolor{poporange!30}\texttt{Non validé} & Compte « \texttt{Info\_Officielle\_237} » : 3 jours d'existence, 45 abonnés, pas de badge bleu. Pas de signature « MINEDUB » ou « Ministre ». L'URL \texttt{.cm.ng} (Nigéria) n'est pas un domaine gouvernemental camerounais (\texttt{.gov.cm} ou \texttt{.cm}). \\\hline
\textbf{R -- Pertinence} & \cellcolor{popgreen!30}\texttt{Validé} & Besoin réel des parents (vie chère, rentrée). Cible précise (parents d'élèves). L'arnaque exploite une préoccupation légitime. \\\hline
\textbf{A -- Actualité} & \cellcolor{poporange!30}\texttt{Partiel} & Publié « il y a 2h » (cohérence temporelle rentrée), mais \textbf{aucune date officielle} de signature d'arrêté ou de communiqué n'est mentionnée. \\\hline
\textbf{F -- Fiabilité} & \cellcolor{poporange!30}\texttt{Non validé} & \textbf{Zéro source croisée.} Lien vers \texttt{.ng} (étranger). Demande de données sensibles (numéro MoMo/OM) sur page non sécurisée (HTTP probable, pas HTTPS visible). Aucun numéro d'arrêté, aucun décret cité. \\\hline
\textbf{O -- Objectivité} & \cellcolor{poporange!30}\texttt{Non validé} & Langage \textbf{urgentiste} : « URGENT », « PARTAGEZ MAX », « instantané », « ça ne coûte rien d'essayer ». Majuscules abusives. Ton incitatif, pas informatif. \\\hline
\textbf{A -- Accessibilité} & \cellcolor{popred!30}\texttt{Non validé} & Lien court masquant l'URL réelle (\texttt{.ng}). Demande de numéro de téléphone \textbf{avant} toute authentification. Risque \textbf{élevé} de vol de session MoMo (OTP) ou d'installation d'APK malveillant. \\\hline
\end{tabularx}
\end{center}

\textbf{Étape 2 -- Rédaction du rapport d'évaluation (Modèle attendu)}

\begin{quote}
\textbf{VERDICT : FAUX -- ARNAQUE CONFIRMÉE (PHISHING MOBILE MONEY)}

\textbf{Arguments majeurs :}
1. \textbf{Crédibilité nulle (Critère C) :} L'émetteur est un compte anonyme récent (3 jours, 45 abonnés) sans badge de vérification. Le domaine \texttt{.cm.ng} appartient au Nigeria, incompatible avec un site officiel camerounais (\texttt{.gov.cm}). Le logo MINEDUB est flou, usurpé.
2. \textbf{Fiabilité absente \& Accessibilité dangereuse (Critères F \& A) :} Aucune référence administrative (numéro d'arrêté, lien vers \texttt{minedub.gov.cm}). La procédure demande le numéro Mobile Money sur un lien externe non sécurisé : technique classique de \emph{phishing} pour détourner des fonds via OTP ou usurpation d'identité.
3. \textbf{Objectivité manipulatoire (Critère O) :} L'usage de majuscules (« URGENT »), d'impératifs viraux (« PARTAGEZ MAX ») et de la promesse de gain facile (« instantané », « ça ne coûte rien ») vise à court-circuiter l'esprit critique par l'urgence et l'appât du gain.

\textbf{Recommandation :} \textbf{NE PAS CLICQUER. NE PAS PARTAGER. SIGNALER} le post sur Facebook (menu « Signaler » $\rightarrow$ « Arnaque »). Avertir immédiatement les groupes de parents : « C'est une arnaque pour voler vos codes MoMo/Orange Money. Le MINEDUB ne verse jamais d'argent par lien Facebook. »
\end{quote}

\textbf{Étape 3 -- Points de vigilance pour la présentation orale (Débrief professeur)}
\begin{itemize}
    \item \textbf{Ne pas confondre Pertinence et Fiabilité :} Le fait que l'info « arrange » les parents (Pertinence validée) ne la rend pas vraie. C'est le piège n°1 de l'ingénierie sociale.
    \item \textbf{L'URL est une preuve technique majeure :} \texttt{.cm.ng} $\neq$ \texttt{.cm} ou \texttt{.gov.cm}. Apprenez à lire la structure d'un nom de domaine (TLD).
    \item \textbf{La demande de numéro MoMo/OM sur un lien web est un \textbf{Drapeau Rouge} absolu.} Les opérateurs (MTN, Orange) n'utilisent que l'USSD (\texttt{\#150\#}, \texttt{\#144\#}) ou leurs applications officielles pour les transactions.
    \item \textbf{Orthographe et grammaire :} L'absence de date précise sur l'annonce officielle, couplée à une mise en forme approximative (majuscules excessives), est un signe de rédaction amateur ou automatisée.
\end{itemize}
\end{exoresolu}

\courssub{Bilan et transfert de compétences}
Cette enquête collaborative a permis de transformer des connaissances théoriques (la liste C.R.A.F.O.A) en \textbf{compétence opérationnelle de citoyenneté numérique}. Les points clés acquis sont :
\begin{enumerate}
    \item \textbf{La méthode prime sur l'intuition :} Sans la grille, l'appât du gain (50 000 FCFA) et l'urgence (« instantané ») auraient pu l'emporter. L'analyse structurée neutralise l'émotion.
    \item \textbf{La technique valide le fond :} L'examen de l'URL (\texttt{.ng} vs \texttt{.cm}), de l'âge du compte et du protocole (HTTP vs HTTPS, USSD vs Web) fournit des preuves \emph{irréfutables}, indépendantes du contenu textuel.
    \item \textbf{La responsabilité collective :} Le rapport rédigé et la présentation orale outillent l'élève pour devenir un « lanceur d'alerte » protecteur dans son entourage familial (parents, frères/sœurs moins méfiants).
    \item \textbf{Transférabilité :} Cette grille s'applique à \textbf{toute} information : offre d'emploi suspecte sur LinkedIn, rumeur sanitaire sur WhatsApp, promotion « trop belle » sur Jumia/Kaymu, actualité politique non sourcée.
\end{enumerate}

\begin{aretenir}[Réflexe « 3 Questions » avant de partager]
Avant de cliquer sur « Partager » ou « Transférer » sur WhatsApp/Facebook, imposez-vous ce délai de 30 secondes :
\begin{enumerate}
    \item \textbf{Source ?} Est-ce le site officiel (\texttt{.gov.cm}), la CRTV, un journal reconnu (\texttt{Cameroon Tribune}, \texttt{Le Jour}) ou un compte anonyme « Info\_237 » ?
    \item \textbf{Preuve ?} Y a-t-il un numéro de décret, un lien direct vers le PDF officiel, une citation nominative du Ministre ?
    \item \textbf{Risque ?} On me demande mon code PIN, OTP, numéro de téléphone, ou de cliquer sur un lien raccourci (bit.ly, tinyurl) ? $\rightarrow$ \textbf{STOP. C'est une arnaque.}
\end{enumerate}
\end{aretenir}

\newpage

\coursec{Méthodes et exercices résolus}

\courssub{Méthode 1 : Identifier une information et la distinguer d'une simple donnée}

\begin{definition}[Donnée et information]
Une \cle{donnée} est un élément brut, sans interprétation : un chiffre, un mot, une image, un son. Une \cle{information} est une donnée placée dans un contexte qui lui donne un sens et qui apprend quelque chose à celui qui la reçoit. On retiendra la formule : \textbf{information = donnée + sens + contexte}.
\end{definition}

\begin{methode}[Reconnaître une information]
Pour déterminer si un message constitue une information, on procède par étapes :
\begin{enumerate}
\item \textbf{Lire attentivement} le message ou l'élément présenté (texte, image, chiffre, son).
\item \textbf{Se poser la question} : cet élément m'apprend-il quelque chose de nouveau ? Répond-il à une question que je me pose ?
\item \textbf{Vérifier le contexte} : une donnée brute (un chiffre seul, par exemple « 37 ») ne devient une information que si on sait de quoi elle parle (« la température à Yaoundé est de 37 degrés »).
\item \textbf{Conclure} : si l'élément a un sens pour celui qui le reçoit et peut l'aider à comprendre ou à décider, c'est une \cle{information}.
\end{enumerate}
\textbf{Réflexe à retenir} : une information = une donnée + un sens + un contexte. Un chiffre isolé n'est qu'une donnée ; placé dans une phrase claire, il devient une information.
\end{methode}

\begin{exoresolu}[Donnée ou information ?]
Dans le carnet de santé d'un élève, on lit les éléments suivants : (a) « 12 » ; (b) « L'élève a 12 ans » ; (c) « rouge » ; (d) « Le vaccin contre la fièvre jaune a été fait le 12 mars 2024 au Centre de Santé de Biyem-Assi ». Pour chaque élément, dire s'il s'agit d'une donnée brute ou d'une information, en justifiant.
\tcblower
Appliquons la méthode : un élément est une information s'il a un sens clair dans un contexte et apprend quelque chose au lecteur.
\begin{itemize}
\item \textbf{(a) « 12 »} : c'est une donnée brute. Seul, ce chiffre ne dit rien : est-ce un âge, une note, un numéro de maison ? Il manque le contexte. Ce n'est donc pas une information.
\item \textbf{(b) « L'élève a 12 ans »} : c'est une information. La phrase donne un sens au chiffre 12 (un âge) et précise de qui on parle (l'élève). Elle apprend quelque chose de précis au lecteur.
\item \textbf{(c) « rouge »} : c'est une donnée brute. On ne sait pas ce qui est rouge : un carnet, une alerte, un vêtement ? Sans contexte, ce mot n'informe pas.
\item \textbf{(d) « Le vaccin contre la fièvre jaune a été fait le 12 mars 2024 au Centre de Santé de Biyem-Assi »} : c'est une information complète. Elle précise l'événement (vaccination), la nature (fièvre jaune), la date (12 mars 2024) et le lieu (Centre de Santé de Biyem-Assi). Elle peut aider le médecin ou les parents à prendre une décision (par exemple savoir si un rappel est nécessaire).
\end{itemize}
\textbf{Conclusion} : les éléments (a) et (c) sont de simples données brutes, tandis que (b) et (d) sont des informations, car elles associent une donnée à un sens et à un contexte clairs.
\end{exoresolu}

\courssub{Méthode 2 : Mémoriser et définir les six critères d'une bonne information}

\begin{methode}[Les six critères à connaître par cœur]
Au BEPC, il faut savoir citer et définir les \cle{six critères d'une bonne information}. Pour chacun, retiens une définition simple et un exemple :
\begin{enumerate}
\item \textbf{Exactitude} : l'information est vraie, conforme à la réalité, sans erreur.
\item \textbf{Pertinence} : l'information est utile, elle correspond au besoin de celui qui la recherche.
\item \textbf{Actualité (ou fraîcheur)} : l'information est récente, à jour.
\item \textbf{Exhaustivité (ou complétude)} : l'information est complète, rien d'important ne manque.
\item \textbf{Fiabilité (ou crédibilité de la source)} : l'information provient d'une source digne de confiance.
\item \textbf{Clarté (ou compréhensibilité)} : l'information est présentée de façon claire, facile à comprendre.
\end{enumerate}
\textbf{Moyen mnémotechnique} : retiens la phrase « \textbf{E}lle \textbf{P}art \textbf{A}vec \textbf{E}nfants \textbf{F}idèles et \textbf{C}lairs » (Exactitude, Pertinence, Actualité, Exhaustivité, Fiabilité, Clarté).
\end{methode}

\begin{aretenir}[L'essentiel sur les critères]
Une \cle{bonne information} est une information qui respecte les six critères : elle est \textbf{exacte} (vraie), \textbf{pertinente} (utile), \textbf{actuelle} (à jour), \textbf{exhaustive} (complète), \textbf{fiable} (source crédible) et \textbf{claire} (compréhensible). Si un seul de ces critères fait défaut, l'information doit être vérifiée avant d'être crue, utilisée ou partagée.
\end{aretenir}

\begin{exoresolu}[Associer chaque situation au bon critère]
Pour chacune des situations suivantes, citer le critère d'une bonne information qui est respecté ou qui fait défaut, et justifier en une phrase.
\begin{enumerate}
\item Un site annonce : « Le Cameroun compte 8 régions. »
\item Pour préparer un exposé sur la coupe du monde 2022, un élève consulte un article sur la coupe du monde 2022.
\item Une affiche de la mairie indique la date et le lieu de la fête de la jeunesse, mais pas l'heure.
\item Un journal publie les résultats du BEPC en se basant sur le communiqué officiel du MINESEC.
\end{enumerate}
\tcblower
On analyse chaque situation en comparant avec les définitions des six critères.
\begin{enumerate}
\item \textbf{Critère en défaut : l'exactitude.} L'information est fausse : le Cameroun compte 10 régions et non 8. Une bonne information doit être conforme à la réalité.
\item \textbf{Critère respecté : la pertinence.} L'article correspond exactement au besoin de l'élève (le sujet de son exposé). Une information pertinente répond au besoin de celui qui la cherche.
\item \textbf{Critère en défaut : l'exhaustivité.} L'information est incomplète : il manque l'heure, un élément essentiel pour se rendre à la fête. Une bonne information ne doit rien omettre d'important.
\item \textbf{Critère respecté : la fiabilité.} La source (le communiqué officiel du MINESEC) est digne de confiance, car c'est l'organisme qui organise l'examen. Une information fiable vient d'une source crédible.
\end{enumerate}
\textbf{Conclusion} : chaque situation illustre un critère précis ; savoir nommer le critère ET le justifier par la situation est indispensable pour obtenir tous les points.
\end{exoresolu}

\courssub{Méthode 3 : Évaluer une information rencontrée dans la vie courante}

\begin{methode}[Appliquer les critères à une situation concrète]
Face à une information reçue (message WhatsApp, affiche, annonce radio, publication Facebook), suis cette démarche en 5 étapes :
\begin{enumerate}
\item \textbf{Identifier la source} : qui a produit cette information ? Est-ce une personne ou un organisme connu et sérieux ? (critère de fiabilité)
\item \textbf{Vérifier la date} : quand l'information a-t-elle été publiée ? Est-elle encore valable aujourd'hui ? (critère d'actualité)
\item \textbf{Contrôler la véracité} : les faits annoncés sont-ils vrais ? Peut-on les confirmer ailleurs ? (critère d'exactitude)
\item \textbf{Examiner le contenu} : l'information est-elle complète ? Est-elle facile à comprendre ? (critères d'exhaustivité et de clarté)
\item \textbf{Se demander si elle nous est utile} : répond-elle à notre besoin du moment ? (critère de pertinence)
\end{enumerate}
\textbf{Conclusion de la démarche} : si l'information échoue à un ou plusieurs critères, il faut la vérifier auprès d'une autre source avant de la croire ou de la partager.
\end{methode}

\begin{exoresolu}[Évaluer un message reçu sur WhatsApp]
Amina, élève en classe de 3ème à Douala, reçoit sur WhatsApp le message suivant, transféré de nombreuses fois : « URGENT !!! Le BEPC est annulé cette année à cause des travaux dans les lycées. Partage à tous tes camarades !!! » En appliquant la démarche d'évaluation, dire si Amina doit croire et partager ce message. Justifier la réponse en citant au moins trois critères.
\tcblower
Appliquons les 5 étapes de la démarche.
\begin{enumerate}
\item \textbf{Source} : le message est anonyme et simplement transféré. Il ne provient ni du MINESEC, ni du principal du collège, ni d'un organe officiel. \textbf{Le critère de fiabilité n'est pas respecté.}
\item \textbf{Date} : le message ne porte aucune date ni aucune référence à une décision officielle. On ne peut pas vérifier son actualité.
\item \textbf{Véracité} : une décision aussi grave (annulation d'un examen national) serait annoncée par voie officielle (communiqué du MINESEC, radio nationale, affichage au collège). Rien de tel n'existe. \textbf{Le critère d'exactitude n'est pas vérifié.}
\item \textbf{Contenu} : le message est vague (quels travaux ? quels lycées ? quelle année exactement ?). \textbf{Les critères d'exhaustivité et de clarté ne sont pas respectés.}
\item \textbf{Utilité} : l'information concerne bien Amina (elle passe le BEPC), donc elle est pertinente pour elle, mais cela ne suffit pas à la rendre bonne.
\end{enumerate}
\textbf{Conclusion} : ce message ne respecte ni la fiabilité, ni l'exactitude, ni l'exhaustivité. Amina ne doit ni le croire ni le partager ; elle doit vérifier auprès de son établissement ou des canaux officiels du MINESEC. Partager un tel message reviendrait à propager une fausse information (rumeur).
\end{exoresolu}

\courssub{Méthode 4 : Rédiger et justifier une évaluation critique d'une information}

\begin{methode}[Rédiger une réponse justifiée au BEPC]
Pour rédiger une évaluation critique claire et complète, adopte le plan suivant :
\begin{enumerate}
\item \textbf{Rappeler l'information étudiée} en une phrase (de quoi parle-t-elle ? d'où vient-elle ?).
\item \textbf{Analyser critère par critère} : pour chaque critère concerné, écris une phrase du type « Le critère de \ldots{} est (ou n'est pas) respecté car \ldots ». Appuie-toi toujours sur un élément précis de l'énoncé.
\item \textbf{Conclure par un jugement global} : l'information est-elle bonne, moyenne ou mauvaise ?
\item \textbf{Proposer une conduite à tenir} : que doit faire la personne (croire, vérifier, ignorer, signaler) ?
\end{enumerate}
\textbf{Réflexes de rédaction} : citer le nom exact du critère ; justifier chaque affirmation par un fait tiré de la situation ; rédiger en phrases complètes, pas en mots isolés ; soigner l'orthographe des termes techniques (exactitude, pertinence, fiabilité\ldots).
\end{methode}

\begin{exoresolu}[Rédiger l'évaluation critique d'une affiche]
Dans un quartier de Yaoundé, une affiche annonce : « Grande campagne gratuite de vaccination contre la méningite, samedi prochain. » L'affiche ne porte ni logo, ni nom d'organisme, ni lieu précis, et le papier est déchiré à l'endroit de la date exacte. Rédige un court texte (8 à 10 lignes) évaluant cette information et indiquant la conduite à tenir pour un habitant du quartier.
\tcblower
\textbf{Évaluation critique de l'affiche}

L'affiche étudiée annonce une campagne gratuite de vaccination contre la méningite, prévue « samedi prochain » dans le quartier.

Analysons cette information critère par critère. D'abord, le critère de fiabilité n'est pas respecté : l'affiche ne porte ni logo ni nom d'organisme (ministère de la Santé publique, centre de santé, mairie), on ne sait donc pas qui organise cette campagne. Ensuite, le critère d'exhaustivité n'est pas respecté non plus : il manque le lieu précis de la vaccination et la date exacte est illisible car le papier est déchiré. Le critère d'actualité ne peut pas être vérifié, puisque « samedi prochain » ne précise aucune date : l'affiche est peut-être ancienne. En revanche, le critère de pertinence est respecté, car une campagne de vaccination gratuite est une information utile pour les habitants.

En conclusion, cette information est de mauvaise qualité : incomplète et de source inconnue. La conduite à tenir est de ne pas s'y rendre sur la seule base de cette affiche, mais de vérifier l'information auprès du centre de santé le plus proche ou de la mairie avant tout déplacement.
\end{exoresolu}

\courssub{Méthode 5 : Vérifier une information en croisant les sources}

\begin{methode}[Le réflexe du croisement des sources]
Pour s'assurer qu'une information est bonne avant de l'utiliser ou de la transmettre :
\begin{enumerate}
\item \textbf{Noter l'affirmation principale} de l'information reçue (le fait essentiel annoncé).
\item \textbf{Chercher la même information dans au moins deux sources différentes} et indépendantes (site officiel, radio nationale, enseignant, document imprimé officiel).
\item \textbf{Comparer} : les sources disent-elles la même chose ? Les dates, les chiffres, les noms concordent-ils ?
\item \textbf{Privilégier les sources officielles} : pour un examen, le MINESEC ; pour la santé, le MINSANTE ; pour la météo, la météorologie nationale.
\item \textbf{Conclure} : si les sources concordent et sont fiables, l'information peut être considérée comme exacte ; sinon, elle reste douteuse.
\end{enumerate}
\textbf{Réflexe} : une information confirmée par au moins deux sources fiables et indépendantes est beaucoup plus crédible qu'un message isolé.
\end{methode}

\begin{exoresolu}[Croiser les sources avant de décider]
Le père de Marlyse entend à la radio locale de Bafoussam que « les inscriptions au concours d'entrée en classe de 2nde au lycée technique sont ouvertes jusqu'au 30 septembre, moyennant des frais de 5\,000 FCFA ». Un voisin lui affirme que les frais sont de 10\,000 FCFA et que la date limite est le 15 septembre. Explique la démarche que doit suivre le père de Marlyse pour obtenir une bonne information, et dis quelle source il doit privilégier.
\tcblower
Appliquons la méthode du croisement des sources.
\begin{enumerate}
\item \textbf{Affirmation principale} : deux informations contradictoires circulent sur les inscriptions au concours (date limite : 30 septembre ou 15 septembre ; frais : 5\,000 FCFA ou 10\,000 FCFA).
\item \textbf{Sources disponibles} : la radio locale (source médiatique, plutôt fiable mais pouvant se tromper) et le voisin (source non officielle, simple rumeur possible). Aucune des deux n'est la source directe de l'information.
\item \textbf{Source officielle à consulter} : pour un concours d'entrée dans un lycée public, l'information officielle vient du lycée lui-même (affichage au secrétariat) ou du MINESEC (arrêté ou communiqué publiant le calendrier et les frais des concours).
\item \textbf{Démarche concrète} : le père de Marlyse doit se rendre au secrétariat du lycée technique concerné ou consulter le communiqué officiel du MINESEC, puis comparer les dates et les montants annoncés avec ce qu'il a entendu.
\item \textbf{Conclusion} : il doit privilégier la source officielle (le lycée ou le MINESEC), car c'est elle qui organise le concours : c'est la seule source dont la fiabilité est garantie. Tant qu'il n'a pas vérifié, il ne doit ni payer ni transmettre l'information à d'autres parents.
\end{enumerate}
\textbf{Conclusion} : face à deux versions contradictoires, seul le croisement avec une source officielle permet d'obtenir une information exacte et fiable.
\end{exoresolu}

\courssub{Méthode 6 : Appliquer les critères à la recherche d'information sur Internet}

\begin{methode}[Bien choisir ses sources sur Internet]
Internet regorge d'informations de qualité très variable. Pour un travail scolaire :
\begin{enumerate}
\item \textbf{Préférer les sites officiels et éducatifs} : sites des ministères (.gov.cm), des organismes reconnus, des encyclopédies sérieuses.
\item \textbf{Vérifier l'auteur et la date de publication} de la page consultée : une page sans auteur ni date est suspecte.
\item \textbf{Se méfier des réseaux sociaux} : tout le monde peut y publier n'importe quoi ; un grand nombre de partages ne prouve pas qu'une information est vraie.
\item \textbf{Comparer au moins deux sites} avant de retenir un fait ou un chiffre.
\item \textbf{Citer ses sources} dans le travail rendu (nom du site, titre de la page, date de consultation).
\end{enumerate}
\end{methode}

\begin{exoresolu}[Choisir la bonne source pour un exposé]
Pour un exposé sur « la population du Cameroun », trois élèves proposent chacun une source : (a) un commentaire anonyme sous une vidéo ; (b) le site de l'Institut National de la Statistique (INS) du Cameroun ; (c) un ancien manuel scolaire publié il y a 25 ans. Pour chaque source, évaluer sa qualité en citant le ou les critères concernés, puis indiquer laquelle choisir.
\tcblower
Évaluons chaque source avec les six critères.
\begin{itemize}
\item \textbf{(a) Commentaire anonyme sous une vidéo} : source de très mauvaise qualité. L'auteur est inconnu (\textbf{fiabilité} non respectée), le chiffre donné n'est vérifiable par personne (\textbf{exactitude} non garantie) et on ne connaît pas sa date (\textbf{actualité} invérifiable).
\item \textbf{(b) Site de l'INS du Cameroun} : source de bonne qualité. L'INS est l'organisme officiel chargé des statistiques au Cameroun (\textbf{fiabilité} respectée) ; ses chiffres sont issus de recensements et d'études sérieuses (\textbf{exactitude}) et le site publie des données récentes (\textbf{actualité}). C'est aussi une source complète et claire (\textbf{exhaustivité}, \textbf{clarté}).
\item \textbf{(c) Ancien manuel de 25 ans} : source moyenne. Le manuel est un document sérieux (\textbf{fiabilité} correcte) et ses chiffres étaient exacts à l'époque, mais la population du Cameroun a beaucoup augmenté depuis : le critère d'\textbf{actualité} n'est pas respecté, donc l'information est devenue fausse pour aujourd'hui.
\end{itemize}
\textbf{Conclusion} : il faut choisir la source (b), le site de l'INS, car c'est la seule qui respecte à la fois la fiabilité, l'exactitude et l'actualité. Cet exemple montre qu'une information autrefois exacte peut devenir mauvaise si elle n'est plus à jour.
\end{exoresolu}

\begin{savaistu}[Les fausses informations, un danger réel]
Les fausses informations (appelées aussi \emph{fake news} ou intox) circulent très vite sur les réseaux sociaux. Au Cameroun comme ailleurs, des rumeurs sur la santé, les examens ou la sécurité ont déjà provoqué des paniques. La loi camerounaise punit d'ailleurs la diffusion de fausses nouvelles. Vérifier une information avant de la partager, c'est donc à la fois un réflexe de bon élève, un devoir de citoyen et une protection contre les sanctions.
\end{savaistu}

\begin{attention}[Les 5 erreurs qui coûtent des points au BEPC]
\begin{enumerate}
\item \textbf{Confondre donnée et information.} Dire « 12 est une information » est faux : sans contexte, c'est une donnée brute. Exige toujours un sens et un contexte avant de parler d'information.
\item \textbf{Citer le critère sans le justifier.} Écrire « le critère est la fiabilité » sans expliquer pourquoi (source inconnue, organisme officiel\ldots) fait perdre la moitié des points : toute réponse doit être appuyée par un élément précis de l'énoncé.
\item \textbf{Mélanger les critères entre eux.} Erreurs classiques : confondre \emph{exactitude} (l'information est vraie) avec \emph{fiabilité} (la source est digne de confiance), ou \emph{actualité} (récente) avec \emph{pertinence} (utile au besoin). Revois les définitions avant l'épreuve.
\item \textbf{Oublier de conclure et de proposer une conduite à tenir.} Une évaluation critique se termine toujours par un jugement global (« cette information est de mauvaise qualité ») ET par un conseil pratique (vérifier auprès de\ldots, ne pas partager\ldots).
\item \textbf{Croire qu'un message très partagé est forcément vrai.} Le nombre de partages sur WhatsApp ou Facebook n'est pas un critère de qualité : une rumeur peut circuler des milliers de fois. Seuls les six critères (et le croisement des sources) permettent de juger une information.
\end{enumerate}
\end{attention}

\newpage

\coursec{Exercices}

\courssub{Je vérifie mes connaissances}

\begin{exercice}[QCM : les critères d'une bonne information]
Pour chaque question, entoure la (ou les) bonne(s) réponse(s).
\begin{enumerate}
\item Une information est dite \cle{fiable} lorsque :
\begin{enumerate}
\item elle provient d'une source digne de confiance ;
\item elle est partagée par beaucoup de personnes ;
\item elle est véridique et vérifiable.
\end{enumerate}
\item Le critère qui consiste à vérifier que l'information est récente et encore valable est :
\begin{enumerate}
\item l'exactitude ;
\item l'actualité (la fraîcheur) ;
\item la pertinence.
\end{enumerate}
\item Une information \cle{pertinente} est une information :
\begin{enumerate}
\item qui répond au besoin de celui qui la cherche ;
\item qui est écrite en bon français ;
\item qui est très longue et détaillée.
\end{enumerate}
\item Une information \cle{objective} est une information :
\begin{enumerate}
\item qui présente les faits sans parti pris ni sentiment personnel ;
\item qui défend une seule opinion ;
\item qui est publiée sur un site officiel.
\end{enumerate}
\end{enumerate}
\end{exercice}

\begin{exercice}[Vrai ou faux ? Justifie ta réponse]
Réponds par Vrai ou Faux, puis justifie ta réponse en une ou deux phrases.
\begin{enumerate}
\item Une information publiée sur un site gouvernemental est toujours exacte et ne doit jamais être vérifiée.
\item Une information peut être exacte mais ne pas être pertinente pour mon besoin.
\item La date de publication fait partie des éléments à contrôler pour évaluer une information.
\item Une information partagée des milliers de fois sur WhatsApp est forcément une bonne information.
\end{enumerate}
\end{exercice}

\begin{exercice}[Définitions]
Définis clairement, en une ou deux phrases chacune, les termes suivants :
\begin{enumerate}
\item information ;
\item source d'information ;
\item exactitude d'une information ;
\item complétude d'une information.
\end{enumerate}
\end{exercice}

\begin{exercice}[Associer critère et situation]
Recopie le tableau ci-dessous et associe chaque situation au critère de bonne information qu'elle met en défaut (fiabilité, exactitude, actualité, pertinence, objectivité, complétude).

\begin{center}
\begin{tabularx}{\linewidth}{|c|X|X|}
\hline
\rowcolor{popblueL} \textbf{N°} & \textbf{Situation} & \textbf{Critère non respecté} \\
\hline
1 & Un site annonce les dates de la rentrée scolaire 2019-2020 alors que nous sommes en 2026. & \\
\hline
2 & Un texte affirme que le Cameroun compte 58 régions administratives. & \\
\hline
3 & Un article vantant un produit de beauté est rédigé par le vendeur de ce produit. & \\
\hline
4 & Un message anonyme, sans auteur ni source, annonce la fermeture de tous les lycées du pays. & \\
\hline
5 & Un article sur les résultats du BEPC ne donne que les taux de Yaoundé et oublie les neuf autres régions. & \\
\hline
6 & Pour réviser la leçon sur les réseaux, je lis un article sur les recettes de cuisine. & \\
\hline
\end{tabularx}
\end{center}
\end{exercice}

\courssub{Je m'entraîne}

\begin{exercice}[Analyse d'un message sur les réseaux sociaux]
Voici la description d'une capture d'écran d'un tweet : « ALERTE !!! Grève générale dans tous les établissements scolaires du Cameroun à partir de lundi. Partagez massivement pour sauver nos enfants ! ». Le compte qui publie ce message s'appelle « InfosVraies237 », créé il y a trois jours, sans photo d'identité, et le message ne cite aucune source officielle ni aucune date précise d'année.

Analyse ce message à l'aide des six critères d'une bonne information : pour chaque critère, indique s'il est respecté ou non, et justifie ta réponse. Conclus en indiquant si tu partagerais ce message et pourquoi.
\end{exercice}

\begin{exercice}[Comparer deux sources]
On cherche à connaître le calendrier officiel de l'année scolaire en cours (dates de rentrée, de congés et d'examens). On dispose de deux sources :
\begin{itemize}
\item \textbf{Source A} : le site officiel du MINEDUB (Ministère de l'Éducation de base), page mise à jour le mois dernier, signée par un service identifié ;
\item \textbf{Source B} : le blog personnel d'un enseignant, article publié il y a deux ans, donnant des dates « probables » selon lui.
\end{itemize}
\begin{enumerate}
\item Recopie et complète le tableau comparatif ci-dessous en évaluant chaque source selon les critères : fiabilité de la source, exactitude, actualité, objectivité, complétude, pertinence (note chaque critère : respecté / partiellement respecté / non respecté, avec une courte justification).
\item Indique quelle source tu utiliseras pour préparer ton année scolaire, et justifie ton choix en trois phrases.
\end{enumerate}

\begin{center}
\begin{tabularx}{\linewidth}{|l|X|X|}
\hline
\rowcolor{popblueL} \textbf{Critère} & \textbf{Source A (MINEDUB)} & \textbf{Source B (blog)} \\
\hline
Fiabilité de la source & & \\
\hline
Exactitude & & \\
\hline
Actualité & & \\
\hline
Objectivité & & \\
\hline
Complétude & & \\
\hline
Pertinence & & \\
\hline
\end{tabularx}
\end{center}
\end{exercice}

\begin{exercice}[Ma check-list avant de partager]
Crée ta check-list personnelle de dix questions à te poser avant de partager une information sur les réseaux sociaux (WhatsApp, Facebook, TikTok...). Chaque question doit être formulée clairement, commencer par une majuscule et se terminer par un point d'interrogation, et correspondre à l'un des critères d'une bonne information ou à une règle d'éthique numérique. Présente ta check-list sous forme de liste numérotée, avec un titre.
\end{exercice}

\begin{exercice}[Chasse à la fausse information]
Pendant les vacances, ton petit frère reçoit sur WhatsApp le message suivant : « Le Gouvernement offre 50 000 FCFA à tous les élèves de 3ème qui s'inscrivent sur le lien www.bourses-gratuites-cm.xyz avant ce soir minuit. Envoyez votre nom, votre école et le numéro de téléphone de vos parents. »
\begin{enumerate}
\item Releve dans ce message au moins quatre indices qui doivent te rendre méfiant.
\item Pour chaque indice, indique le critère d'une bonne information qui est mis en cause.
\item Explique en quelques phrases ce que tu conseilles à ton frère de faire (et de ne pas faire) avec ce message.
\end{enumerate}
\end{exercice}

\courssub{Je me prépare au BEPC}

\begin{exercice}[L'article sur la réforme du système éducatif]
Ton professeur d'informatique te demande d'évaluer un article publié sur le site d'information en ligne « LeCamerounaisNet.com », intitulé « Réforme du système éducatif : tout ce qui va changer à la rentrée prochaine ». L'article est signé par un journaliste identifié, il cite un communiqué du MINESEC en donnant la date exacte de celui-ci, il a été publié il y a deux semaines, il présente les arguments du ministère et ceux d'un syndicat d'enseignants, mais il ne précise pas si la réforme concerne aussi les écoles privées.

\begin{enumerate}
\item Rappelle la liste des six critères d'une bonne information étudiés en classe.
\item Évalue l'article critère par critère, en t'appuyant sur les éléments de l'énoncé.
\item Rédige un texte argumenté d'environ 180 mots dans lequel tu justifies ton jugement global sur la fiabilité de cet article : peux-tu l'utiliser pour préparer un exposé sur la réforme du système éducatif ? Ta conclusion devra indiquer clairement ta position et les vérifications complémentaires que tu effectuerais.
\end{enumerate}
\end{exercice}

\begin{exercice}[Enquête au CDI du collège]
Dans le cadre d'un exposé sur « l'accès à l'eau potable dans la région du Centre », tu trouves au CDI de ton collège les trois documents suivants :

\begin{itemize}
\item \textbf{Document 1} : un rapport du Ministère de l'Eau et de l'Énergie (MINEE) publié cette année, avec des statistiques détaillées par commune ;
\item \textbf{Document 2} : un article d'un journal national datant de dix ans, très complet mais ancien ;
\item \textbf{Document 3} : une page Facebook d'une association inconnue affirmant, sans chiffre ni référence, que « 90 \% des habitants du Centre n'ont pas l'eau potable ».
\end{itemize}

\begin{enumerate}
\item Pour chacun des trois documents, évalue les critères de fiabilité, d'actualité et d'exactitude, en justifiant chaque réponse.
\item Classe les trois documents du plus fiable au moins fiable pour ton exposé, en expliquant ton classement.
\item Le Document 2 est ancien mais très complet : explique dans quelles conditions tu pourrais quand même l'utiliser (par exemple pour comparer l'évolution de la situation).
\item Rédige une courte conclusion (5 à 8 lignes) présentant la démarche qu'un élève doit toujours suivre avant d'utiliser un document dans un travail scolaire.
\end{enumerate}
\end{exercice}

\begin{exercice}[Rumeur de fermeture des marchés à Douala]
Un soir, un message vocal circule dans les groupes WhatsApp des parents d'élèves de ton quartier à Douala : une voix anonyme annonce que « tous les marchés de la ville seront fermés pendant un mois à partir de demain ». En quelques heures, le message est transféré des centaines de fois et certains habitants se précipitent pour faire des provisions. Le lendemain, la radio régionale CRTV Littoral dément l'information en citant un communiqué de la Communauté Urbaine de Douala : aucune fermeture n'est prévue.

\begin{enumerate}
\item Identifie dans cette situation au moins trois signaux qui auraient dû alerter les destinataires du message vocal, en précisant pour chacun le critère de bonne information concerné.
\item Explique pourquoi la rapidité de diffusion d'une information sur les réseaux sociaux ne prouve pas son exactitude.
\item Décris, sous forme d'une démarche en étapes numérotées, ce qu'aurait dû faire un parent d'élève prudent en recevant ce message vocal, avant de le transférer ou d'agir.
\item Rédige un court message (5 lignes maximum) que tu enverrais dans le groupe WhatsApp pour alerter calmement les membres sur les dangers de ce type de rumeur et leur proposer un réflexe simple à adopter.
\end{enumerate}
\end{exercice}

\newpage

\coursec{Corrigés des exercices}

\coursec{Corrigés des exercices — Acquisition d'une maîtrise des médias et de l'informatique}

\begin{corrige}[Exercice 1 — QCM : les critères d'une bonne information]
\begin{enumerate}
\item Réponses \textbf{a} et \textbf{c}. La \cle{fiabilité} repose sur la confiance accordée à la source et sur le caractère véridique et vérifiable de l'information. Le nombre de partages (réponse b) ne prouve rien : une rumeur peut être partagée des millions de fois.
\item Réponse \textbf{b}. L'\cle{actualité} (ou fraîcheur) est le critère qui exige que l'information soit récente et encore valable au moment où on l'utilise.
\item Réponse \textbf{a}. La \cle{pertinence} mesure l'adéquation entre l'information et le besoin de celui qui la recherche.
\item Réponse \textbf{a}. L'\cle{objectivité} consiste à présenter les faits sans parti pris, sans sentiment personnel et sans chercher à défendre une seule opinion.
\end{enumerate}
\textbf{Points de vigilance} : ne pas confondre popularité et fiabilité ; bien distinguer \cle{exactitude} (conformité aux faits) et \cle{actualité} (date).
\end{corrige}

\begin{corrige}[Exercice 2 — Vrai ou faux ? Justifie ta réponse]
\begin{enumerate}
\item \textbf{Faux.} Un site gouvernemental est une source généralement fiable, mais aucune source n'est à l'abri d'une erreur ou d'une information périmée : il faut toujours vérifier la date et croiser avec d'autres sources.
\item \textbf{Vrai.} Par exemple, la liste exacte des résultats du BEPC en Norvège est exacte mais ne répond pas à mon besoin si je cherche les résultats de mon collège : elle n'est pas \cle{pertinente} pour moi.
\item \textbf{Vrai.} La date de publication permet de juger de l'\cle{actualité} (la fraîcheur) de l'information : une information ancienne peut être devenue fausse ou inutilisable.
\item \textbf{Faux.} Le nombre de partages mesure la popularité d'un message, pas sa qualité : une fausse information peut circuler très vite et très largement.
\end{enumerate}
\textbf{Points de vigilance} : la justification est obligatoire ; une réponse « Vrai » ou « Faux » seule ne rapporte pas tous les points.
\end{corrige}

\begin{corrige}[Exercice 3 — Définitions]
\begin{enumerate}
\item \textbf{Information} : ensemble de faits, de données ou de nouvelles portés à la connaissance de quelqu'un, par un texte, une image, un son ou une vidéo.
\item \textbf{Source d'information} : personne, organisme ou support (site web, journal, radio, livre...) d'où provient une information.
\item \textbf{Exactitude} : qualité d'une information conforme à la réalité des faits, sans erreur ni exagération.
\item \textbf{Complétude} : qualité d'une information complète, qui contient tous les éléments nécessaires à la compréhension du sujet, sans omission importante.
\end{enumerate}
\textbf{Points de vigilance} : une définition doit être courte, précise et ne pas employer le mot défini dans sa propre définition.
\end{corrige}

\begin{corrige}[Exercice 4 — Associer critère et situation]
\begin{enumerate}
\item \textbf{Actualité} : l'information est périmée, elle date de plusieurs années.
\item \textbf{Exactitude} : le Cameroun compte 10 régions, le chiffre annoncé est faux.
\item \textbf{Objectivité} : le vendeur défend son propre intérêt, il y a parti pris.
\item \textbf{Fiabilité} : la source est anonyme et invérifiable, on ne peut pas lui accorder de confiance.
\item \textbf{Complétude} : l'information est incomplète, il manque les données de neuf régions.
\item \textbf{Pertinence} : l'article ne répond pas au besoin (réviser la leçon sur les réseaux).
\end{enumerate}
\textbf{Points de vigilance} : chaque critère n'est utilisé qu'une seule fois ; bien lire toute la situation avant de choisir.
\end{corrige}

\begin{corrige}[Exercice 5 — Analyse d'un message sur les réseaux sociaux]
Analyse critère par critère :
\begin{enumerate}
\item \textbf{Fiabilité : non respecté.} Le compte « InfosVraies237 » est créé il y a trois jours, sans photo d'identité, et ne cite aucune source officielle : on ne peut pas lui accorder de confiance.
\item \textbf{Exactitude : non vérifiable.} Aucun communiqué officiel n'est cité ; l'affirmation ne peut pas être confirmée.
\item \textbf{Actualité : partiellement respecté.} Le message annonce un événement « à partir de lundi », mais sans année précise : on ne sait pas de quel lundi il s'agit.
\item \textbf{Pertinence : respecté en apparence.} Le sujet (grève scolaire) concerne les élèves et les parents, mais une information fausse, même sur un sujet qui nous concerne, reste inutilisable.
\item \textbf{Objectivité : non respecté.} Le ton alarmiste (« ALERTE !!! », « sauver nos enfants ») joue sur les émotions et cherche à provoquer un partage massif.
\item \textbf{Complétude : non respecté.} Il manque la date exacte, les régions concernées, les motifs de la grève et la source officielle.
\end{enumerate}
\textbf{Conclusion} : je ne partagerais pas ce message. Cinq critères sur six ne sont pas respectés ; il présente tous les signes d'une rumeur. Je vérifierais d'abord sur le site du MINESEC ou auprès de mon établissement avant toute diffusion.

\textbf{Points de vigilance} : chaque critère doit être justifié par un élément précis de l'énoncé (nom du compte, date de création, absence de source, ton du message).
\end{corrige}

\begin{corrige}[Exercice 6 — Comparer deux sources]
\begin{enumerate}
\item Tableau comparatif complété :
\begin{center}
\begin{tabularx}{\linewidth}{|l|X|X|}
\hline
\rowcolor{popblueL} \textbf{Critère} & \textbf{Source A (MINEDUB)} & \textbf{Source B (blog)} \\
\hline
\cle{Fiabilité} de la source & Respecté : site officiel, service identifié & Partiellement respecté : enseignant identifié mais simple particulier \\
\hline
\cle{Exactitude} & Respecté : dates officielles & Non respecté : dates « probables », donc incertaines \\
\hline
\cle{Actualité} & Respecté : mise à jour le mois dernier & Non respecté : article vieux de deux ans \\
\hline
\cle{Objectivité} & Respecté : information administrative neutre & Partiellement respecté : avis personnel de l'auteur \\
\hline
\cle{Complétude} & Respecté : calendrier complet et officiel & Partiellement respecté : informations partielles \\
\hline
\cle{Pertinence} & Respecté : répond exactement au besoin & Partiellement respecté : sujet proche mais données dépassées \\
\hline
\end{tabularx}
\end{center}
\item J'utiliserai la \textbf{source A}, le site du MINEDUB. C'est la source officielle chargée de publier le calendrier scolaire : elle est donc fiable et exacte. Elle est récente, complète et répond exactement à mon besoin. La source B, ancienne et fondée sur des suppositions personnelles, ne peut servir que d'indication secondaire.
\end{enumerate}
\textbf{Points de vigilance} : chaque case du tableau doit contenir une note ET une courte justification ; le choix final doit s'appuyer sur au moins deux critères.
\end{corrige}

\begin{corrige}[Exercice 7 — Ma check-list avant de partager]
Exemple de check-list attendue :

\textbf{Ma check-list avant de partager une information}
\begin{enumerate}
\item Qui est l'auteur de cette information et est-il identifiable ? (\cle{fiabilité})
\item La source est-elle connue et digne de confiance ? (\cle{fiabilité})
\item Les faits annoncés sont-ils vérifiables sur un site officiel ? (\cle{exactitude})
\item Quelle est la date de publication de cette information ? (\cle{actualité})
\item Cette information est-elle encore valable aujourd'hui ? (\cle{actualité})
\item Cette information répond-elle vraiment à mon besoin ou à celui de mes contacts ? (\cle{pertinence})
\item Le message présente-t-il les faits sans exagération ni appel aux émotions ? (\cle{objectivité})
\item Tous les éléments importants du sujet sont-ils donnés ? (\cle{complétude})
\item Ce message respecte-t-il la vie privée et la dignité des personnes ? (\cle{éthique numérique})
\item Ai-je croisé cette information avec au moins une autre source indépendante ? (démarche de vérification)
\end{enumerate}
\textbf{Points de vigilance} : dix questions bien formulées (majuscule, point d'interrogation), couvrant les six critères et au moins une règle d'éthique numérique ; toute check-list cohérente et complète est acceptée.
\end{corrige}

\begin{corrige}[Exercice 8 — Chasse à la fausse information]
\begin{enumerate}
\item Indices de méfiance (au moins quatre) :
\begin{itemize}
\item l'adresse du site est inconnue et suspecte (\texttt{www.bourses-gratuites-cm.xyz}), ce n'est pas un site officiel du Gouvernement ;
\item le message promet une somme d'argent (50 000 FCFA) « à tous les élèves », promesse trop belle pour être vraie ;
\item il impose une urgence (« avant ce soir minuit ») pour empêcher toute réflexion ;
\item il demande des données personnelles (nom, école, numéro de téléphone des parents) ;
\item aucune source officielle, aucun communiqué n'est cité.
\end{itemize}
\item Critères mis en cause :
\begin{itemize}
\item site inconnu et absence de source officielle : \cle{fiabilité} ;
\item promesse invérifiable : \cle{exactitude} ;
\item urgence artificielle et ton manipulateur : \cle{objectivité} ;
\item absence de détails officiels (conditions, texte de loi) : \cle{complétude}.
\end{itemize}
\item Conseils à donner : ne \textbf{pas} cliquer sur le lien, ne \textbf{pas} communiquer les informations personnelles demandées, ne \textbf{pas} transférer le message. Vérifier l'information sur les sites officiels (MINESEC, MINEDUB) ou auprès du directeur de l'établissement, puis supprimer le message et, si possible, le signaler. Il s'agit très probablement d'une tentative d'escroquerie visant à collecter des données personnelles.
\end{enumerate}
\textbf{Points de vigilance} : chaque indice relevé doit être associé explicitement à un critère ; les conseils doivent comporter des actions à faire ET à ne pas faire.
\end{corrige}

\begin{corrige}[Exercice 9 — L'article sur la réforme du système éducatif]
\begin{enumerate}
\item Les six critères d'une bonne information sont : la \cle{fiabilité} de la source, l'\cle{exactitude}, l'\cle{actualité} (fraîcheur), la \cle{pertinence}, l'\cle{objectivité} et la \cle{complétude}.
\item Évaluation critère par critère :
\begin{itemize}
\item \textbf{Fiabilité : respecté.} L'article est signé par un journaliste identifié et publié sur un site d'information connu.
\item \textbf{Exactitude : respecté.} L'article cite un communiqué du MINESEC en donnant sa date exacte : les faits sont vérifiables.
\item \textbf{Actualité : respecté.} Publié il y a deux semaines, l'article est récent.
\item \textbf{Pertinence : respecté.} Le sujet correspond exactement au thème de l'exposé.
\item \textbf{Objectivité : respecté.} L'article présente les arguments du ministère et ceux d'un syndicat d'enseignants : les deux points de vue sont donnés.
\item \textbf{Complétude : partiellement respecté.} L'article ne précise pas si la réforme concerne aussi les écoles privées : il manque un élément.
\end{itemize}
\item \textbf{Texte argumenté (environ 180 mots) :}

Cet article me semble globalement fiable et utilisable pour préparer mon exposé. D'abord, sa source est identifiable : il est signé par un journaliste et publié sur un site d'information en ligne connu, ce qui garantit la \cle{fiabilité}. Ensuite, l'\cle{exactitude} est assurée, car l'article s'appuie sur un communiqué officiel du MINESEC dont il donne la date précise : je peux donc vérifier les faits annoncés. L'\cle{actualité} est également respectée, puisque l'article a été publié il y a seulement deux semaines. De plus, l'\cle{objectivité} est remarquable : l'auteur présente à la fois les arguments du ministère et ceux d'un syndicat d'enseignants, sans prendre parti. Cependant, l'article n'est pas parfaitement \cle{complet} : il ne dit pas si la réforme s'applique aux écoles privées. En conclusion, je peux utiliser cet article comme source principale de mon exposé, mais je consulterai en complément le communiqué original du MINESEC et un autre article pour combler cette lacune et confirmer les informations.

\item[] \textbf{Barème indicatif (sur 20)} : liste des six critères (4 pts) ; évaluation critère par critère justifiée par l'énoncé (8 pts) ; texte argumenté structuré avec position claire et vérifications complémentaires (6 pts) ; qualité de la langue (2 pts).
\item[] \textbf{Points de vigilance} : la conclusion doit contenir une position explicite (« je peux / je ne peux pas l'utiliser ») et citer au moins une vérification complémentaire (consulter le communiqué original, croiser les sources).
\end{enumerate}
\end{corrige}

\begin{corrige}[Exercice 10 — Enquête au CDI du collège]
\begin{enumerate}
\item Évaluation des trois documents :
\begin{itemize}
\item \textbf{Document 1 (rapport du MINEE)} : \cle{fiabilité} \textbf{respectée} (ministère officiel) ; \cle{actualité} \textbf{respectée} (publié cette année) ; \cle{exactitude} \textbf{respectée} (statistiques détaillées par commune, vérifiables).
\item \textbf{Document 2 (article de journal, dix ans)} : \cle{fiabilité} \textbf{respectée} (journal national) ; \cle{actualité} \textbf{non respectée} (dix ans, situation ayant pu changer) ; \cle{exactitude} \textbf{probablement respectée pour son époque}, mais plus garantie aujourd'hui.
\item \textbf{Document 3 (page Facebook)} : \cle{fiabilité} \textbf{non respectée} (association inconnue) ; \cle{actualité} \textbf{indéterminée} (pas de date claire) ; \cle{exactitude} \textbf{non vérifiable} (aucun chiffre sourcé, aucune référence).
\end{itemize}
\item Classement du plus fiable au moins fiable : \textbf{Document 1}, puis \textbf{Document 2}, puis \textbf{Document 3}. Le Document 1 est officiel, récent et détaillé ; le Document 2 est fiable mais périmé ; le Document 3 n'offre aucune garantie de source ni de preuve.
\item Le Document 2 peut être utilisé \textbf{à titre de comparaison historique} : par exemple pour montrer l'évolution de l'accès à l'eau potable entre il y a dix ans et aujourd'hui, en le confrontant aux chiffres récents du Document 1. Il ne doit jamais être présenté comme décrivant la situation actuelle.
\item \textbf{Conclusion (démarche)} : avant d'utiliser un document dans un travail scolaire, un élève doit d'abord identifier la source et vérifier qu'elle est digne de confiance. Il doit ensuite contrôler la date de publication pour s'assurer que l'information est encore valable. Il doit enfin vérifier l'exactitude des faits en croisant au moins deux sources indépendantes, et s'assurer que le document répond bien à son besoin. Ce n'est qu'après ces vérifications qu'il peut citer le document dans son travail, en mentionnant sa source.
\end{enumerate}
\textbf{Barème indicatif (sur 20)} : évaluation des trois documents (9 pts) ; classement justifié (3 pts) ; conditions d'utilisation du Document 2 (3 pts) ; conclusion rédigée (5 pts).

\textbf{Points de vigilance} : ne pas rejeter totalement un document ancien, mais préciser à quelles conditions il reste exploitable ; toujours citer ses sources dans un travail scolaire.
\end{corrige}

\begin{corrige}[Exercice 11 — Rumeur de fermeture des marchés à Douala]
\begin{enumerate}
\item Signaux d'alerte :
\begin{itemize}
\item la voix est \textbf{anonyme} : impossible d'identifier l'auteur $\rightarrow$ critère de \cle{fiabilité} ;
\item aucune \textbf{source officielle} n'est citée (ni communiqué, ni autorité nommée) $\rightarrow$ critères de \cle{fiabilité} et d'\cle{exactitude} ;
\item l'annonce est invérifiable et alarmante (« tous les marchés », « pendant un mois ») $\rightarrow$ critère d'\cle{objectivité} (appel à la peur) ;
\item aucune précision sur les motifs ou les modalités $\rightarrow$ critère de \cle{complétude}.
\end{itemize}
\item La rapidité de diffusion ne prouve pas l'exactitude, car transférer un message ne demande aucune vérification : chaque personne le transmet en quelques secondes, souvent par émotion ou par peur. Un message peut donc atteindre des milliers de personnes sans que personne n'ait contrôlé les faits. La popularité d'un message mesure sa circulation, pas sa vérité.
\item Démarche d'un parent prudent :
\begin{enumerate}
\item Ne pas transférer le message immédiatement.
\item Identifier l'auteur du message vocal : est-il connu et identifiable ?
\item Chercher une confirmation officielle : écouter la radio (CRTV), consulter les sites ou pages officiels de la Communauté Urbaine de Douala ou du Gouvernement.
\item Croiser l'information avec au moins deux sources indépendantes et fiables.
\item N'agir et ne partager que si l'information est confirmée par une source officielle ; sinon, avertir le groupe qu'il s'agit d'une rumeur.
\end{enumerate}
\item \textbf{Message pour le groupe WhatsApp :}

« Chers parents, le message vocal sur la fermeture des marchés est une fausse information : la CRTV Littoral a démenti ce matin en citant la Communauté Urbaine de Douala. Partager de telles rumeurs crée une panique inutile. Avant de transférer un message, vérifions toujours auprès d'une source officielle. Un réflexe simple : pas de source officielle, pas de partage ! »
\end{enumerate}
\textbf{Barème indicatif (sur 20)} : trois signaux avec critères associés (6 pts) ; explication diffusion/exactitude (4 pts) ; démarche en étapes numérotées (6 pts) ; message rédigé, calme et correct (4 pts).

\textbf{Points de vigilance} : la démarche doit être présentée en étapes numérotées et ordonnées ; le message final doit rester court (5 lignes maximum), calme et proposer un réflexe concret.
\end{corrige}

\newpage

\coursec{Fiche bilan}

\begin{popfigure}
\begin{tikzpicture}[mindmap, grow cyclic, every node/.style={concept, text=white, font=\small\bfseries},
root concept/.append style={concept color=popblue, font=\large\bfseries},
level 1/.append style={level distance=3.4cm, sibling angle=72},
level 2/.append style={level distance=2.1cm, sibling angle=45, font=\scriptsize}]
\node[root concept]{La bonne information}
  child[concept color=poporange]{ node {Définition}
    child { node {Donnée + sens} }
    child { node {Utile pour décider} }
  }
  child[concept color=popgreen]{ node {Exactitude}
    child { node {Vraie, sans erreur} }
  }
  child[concept color=poppurple]{ node {Pertinence}
    child { node {Utile au besoin} }
  }
  child[concept color=poppink]{ node {Actualité}
    child { node {Récente, à jour} }
  }
  child[concept color=popgold]{ node {Fiabilité}
    child { node {Source crédible} }
  }
  child[concept color=popmuted]{ node {Complétude et clarté}
    child { node {Complète} }
    child { node {Compréhensible} }
  };
\end{tikzpicture}
\legende{Carte mentale de la leçon : l'information et ses critères de qualité.}
\end{popfigure}

\begin{bilan}
\textbf{Définitions clés}
\begin{itemize}
  \item \cle{Information} : donnée à laquelle on donne un \textbf{sens}, utile pour s'informer, décider ou agir.
  \item \cle{Bonne information} : information de qualité, sur laquelle on peut s'appuyer en confiance.
  \item \cle{Source} : personne, document ou site d'où provient l'information (journal officiel, site du gouvernement, média reconnu, inconnu sur les réseaux sociaux\ldots).
  \item \cle{Intoxication} (ou \emph{fake news}) : information fausse diffusée volontairement pour tromper.
\end{itemize}

\textbf{Les six critères d'une bonne information (à connaître par c\oe ur)}

\begin{center}
\begin{tabularx}{\linewidth}{|l|X|}
\hline
\rowcolor{popblueL} \textbf{Critère} & \textbf{Question à se poser} \\
\hline
\cle{Exactitude} & L'information est-elle vraie, sans erreur ? \\
\hline
\cle{Pertinence} & Répond-elle à mon besoin ? \\
\hline
\cle{Actualité} & Est-elle récente et à jour ? \\
\hline
\cle{Fiabilité} & La source est-elle crédible et identifiable ? \\
\hline
\cle{Complétude} & Est-elle complète, sans omission importante ? \\
\hline
\cle{Clarté} & Est-elle claire et compréhensible ? \\
\hline
\end{tabularx}
\end{center}

\textbf{Méthode express : évaluer une information en 4 étapes}
\begin{enumerate}
  \item \textbf{Identifier} la source (qui ? où ? quand ?).
  \item \textbf{Vérifier} chaque critère (exactitude, pertinence, actualité, fiabilité, complétude, clarté).
  \item \textbf{Croiser} avec au moins une autre source fiable.
  \item \textbf{Conclure et justifier} : « Cette information est bonne / douteuse car\ldots » en citant les critères respectés ou non respectés.
\end{enumerate}
\end{bilan}

\begin{aretenir}[Checklist avant le BEPC]
\begin{itemize}
  \item $\square$ Je sais définir une information et la distinguer d'une simple donnée.
  \item $\square$ Je connais par c\oe ur les six critères d'une bonne information.
  \item $\square$ Je sais expliquer chaque critère avec mes propres mots.
  \item $\square$ Je sais identifier la source d'une information (auteur, date, support).
  \item $\square$ Je vérifie l'exactitude en croisant au moins deux sources.
  \item $\square$ Je contrôle la date de publication pour juger de l'actualité.
  \item $\square$ Je me méfie des informations anonymes diffusées sur les réseaux sociaux.
  \item $\square$ Je sais repérer une information incomplète ou présentée de façon trompeuse.
  \item $\square$ Je sais rédiger une conclusion justifiée en citant les critères.
  \item $\square$ Je sais appliquer la démarche à une situation concrète (rumeur WhatsApp, affiche scolaire, annonce officielle).
\end{itemize}
\end{aretenir}

\findecours

\end{document}
